\chapter{图像识别与检索}
\label{chap:vision-image-recognition}

一张路口照片中有两辆车，一辆红色货车靠近相机，一辆蓝色轿车停在公交站旁。说出“这是街道”，还没有说明车在哪里；给出两个框，还没有说明哪辆是红色、哪辆被遮挡、离相机多远；在图库中找到另一张相似的街景，也未必找到了同一辆车的其他照片。本章要回答的问题是：怎样把一幅图像变成这些逐渐具体、而且可以核验的结果，即语义判断、空间范围、测量量和跨图对应。

本章用到两类已有材料。第一部分讲过的经典方法给出了不经学习也能工作的做法：\ref{chap:vision-image-formation}章的相机模型说明像素与三维点怎样对应，\ref{chap:vision-image-transforms}章的金字塔与滤波说明怎样在多个尺度上观察图像，\ref{chap:vision-image-features}章的边缘、局部特征、经典分割和基于人工特征的识别说明怎样从像素中提出可比较的证据。\ref{chap:vision-image-video-models}章和\ref{chap:vision-multimodal-models}章则给出了可复用的视觉表示和图文表示。本章不重讲这些基础，而是研究在具体任务中还缺什么：判断针对整幅图、一个对象还是一个像素；监督约束哪一种输出；网络输出怎样还原成原图坐标、对象编号或米制距离；结果靠什么核验。每类任务先回顾经典路线怎样处理同一问题、在什么条件下失效，再讲学习方法改动了哪一步。

读完本章，读者应能说明各类识别与检索任务的输入、输出与评价指标，指出一个任务的主要困难来自外观、尺度、分辨率、几何还是开放类别，并能沿着本章的模型结构表复核一个检测器、分割器或几何模型的张量形状。第一节先把这些任务和困难摆清楚，后续各节再逐一回答。

街景中的两辆车贯穿对象与空间任务，正常零件与划痕用于说明没有完整缺陷类别表时怎样判断异常，建筑照片用于说明检索与几何恢复。以下算例均为教学构造，数值用于复核机制，不是模型实验结果。

\section{图像识别与检索的问题与挑战}
\label{sec:vision-recognition-challenges}

图像识别与检索不是一个任务，而是一组输出形式不同的任务。它们的输入都是一幅或多幅图像，有时附带类别表、文字、点击或图库；输出则从一个类别名，到一组框、一张逐像素的类别图、若干关节坐标、一张深度图，再到一份按相关程度排序的图库列表。输出形式不同，评价指标也不同，而指标决定了模型被要求做对什么。下面先把各任务的输入、输出和指标列清楚，再讨论使它们变难的五类挑战。

表\ref{tab:vision-recognition-task-metrics}汇总了本章各任务的输入、输出与主要指标。表中每个指标都在下文给出定义；同一名称在不同基准中的细节可能不同，比较结果时必须说明所用协议。

\begin{table*}[tbp]
 \centering
 \small
 \begin{tabularx}{\linewidth}{@{}lXXX@{}}
 \toprule
 任务 & 输入 & 输出 & 主要指标\\
 \midrule
 类别与场景判定 & 整图或对象裁剪，可附类别名 & 一个或多个类别 & 准确率、前$k$准确率、类别平均准确率\\
 异常判定 & 零件图像，只有正常样本可供参照 & 整图异常分数与异常热图 & 图像级与像素级AUROC\\
 边界框预测 & 图像，可附类别表或文字 & 框、类别与得分的集合 & IoU阈值下的AP与mAP\\
 像素类别预测 & 图像，或两期配准影像 & 逐像素类别图 & 类别IoU与mIoU\\
 对象掩码预测 & 图像，可附类别、点框或短语 & 实例掩码集合 & 掩码AP、PQ\\
 前景透明度预测 & 图像与三分图 & 逐像素$\alpha$ & SAD、MSE\\
 计数 & 图像与计入规则 & 数量，可附位置 & MAE、RMSE\\
 姿态估计 & 含人图像 & 二维或三维关节，或人体网格 & OKS下的AP、MPJPE\\
 深度与几何估计 & 一幅或多幅图像，可附相机 & 深度图、点图与相机 & AbsRel、$\delta<1.25$、旋转与平移误差\\
 图像检索 & 查询图与图库 & 图库图像的排序列表 & mAP、Recall@$K$\\
 \bottomrule
 \end{tabularx}
 \caption{本章各任务的输入、输出与主要评价指标。指标的协议细节（阈值、对齐方式、忽略项）随基准而变，报告时应一并给出。}
 \label{tab:vision-recognition-task-metrics}
\end{table*}

类别判定最直接的指标是\term{准确率}（accuracy），即预测类别$\hat y_i$等于标注$y_i$的样本比例$(1/n)\sum_i\mathbf 1\{\hat y_i=y_i\}$。类别很多且彼此相近时，常同时报告前$k$准确率：只要真实类别出现在得分最高的$k$个类别中就算正确。样本在类别间不均衡时，样本平均准确率会被大类主导，此时还要报告各类准确率的平均值，即类别平均准确率。异常判定的分数需要选阈值才能变成“正常／异常”，不同阈值给出不同的误报与漏报。为了在选阈值之前比较排序能力，使用ROC曲线下面积（area under the ROC curve, AUROC），它等于随机抽一个异常样本和一个正常样本、前者分数更高的概率。图像级AUROC按整图分数计算，像素级AUROC把每个像素当作一个样本，检查热图能否把缺陷像素排在正常像素之前。

检测和分割要先回答“预测与哪个真实对象对应”。两个区域的重叠用\term{交并比}（intersection over union, IoU）衡量，它将交集面积除以并集面积：
\begin{equation}
 \operatorname{IoU}(\bm{b},\bm{b}')
 =\frac{|\bm{b}\cap\bm{b}'|}
 {|\bm{b}|+|\bm{b}'|-|\bm{b}\cap\bm{b}'|}\text{。}
 \label{eq:vision-iou}
\end{equation}
$\bm{b},\bm{b}'$可以是矩形框，也可以是掩码覆盖的像素集合。取两个面积均为800的框$(10,10,50,30)$和$(20,10,60,30)$，交集面积为600，IoU为$600/1000=0.6$。检测评价先固定一个IoU阈值$t$，把预测按得分从高到低逐个与同类、尚未被匹配的真实对象配对，IoU不低于$t$的配对算真阳性，其余预测算假阳性，未被配上的真实对象算漏检。逐步降低得分门限，就得到一串精确率$P$（真阳性占全部预测的比例）和召回率$R$（真阳性占全部真实对象的比例），连成精确率—召回率曲线。\term{平均精度}（average precision, AP）是这条曲线下的面积；实际计算时先把精确率改成单调不增的包络$P_{\mathrm{interp}}(r)=\max_{r'\geq r}P(r')$，再在若干召回点上取平均。对所有类别的AP再求平均，就是\term{平均精度均值}（mean average precision, mAP）。COCO检测基准在召回率0至1之间均匀取101个点做插值，并把IoU阈值从0.50到0.95每隔0.05取一次、共10个阈值，对类别和阈值同时求平均；它不区分AP与mAP的写法，直接把这个平均值称为AP。高阈值上的AP奖励定位精细的框，0.5阈值上的AP则主要检查有没有找到对象。

逐像素输出的评价改成按类别累计像素。对类别$r$，令$\mathrm{TP}_r,\mathrm{FP}_r,\mathrm{FN}_r$分别为预测和标注都属于$r$、只有预测属于$r$、只有标注属于$r$的像素数，则类别IoU为$\mathrm{TP}_r/(\mathrm{TP}_r+\mathrm{FP}_r+\mathrm{FN}_r)$，对$K$个类别取平均得到\term{平均交并比}（mean intersection over union, mIoU）。它以类别为单位平均，因此路杆这类小类别与道路这类大类别权重相同。实例掩码沿用检测的AP，只是把框的IoU换成掩码的IoU。全景分割既有可数对象又有道路天空等区域，\term{全景质量}（panoptic quality, PQ）把匹配区域的IoU与对象错误合在一起：对同类、IoU超过0.5的区域建立匹配集合$T$，令$F_{\mathrm{p}},F_{\mathrm{n}}$为多报和漏报区域集合，
\begin{equation}
 \mathrm{PQ}=
 \frac{\sum_{(p,g)\in T}\operatorname{IoU}(p,g)}
 {|T|+\frac12|F_{\mathrm{p}}|+\frac12|F_{\mathrm{n}}|}\text{。}
 \label{eq:vision-pq}
\end{equation}
IoU超过0.5保证一个真实区域至多与一个预测区域匹配，所以匹配无须另作优化\scite{vision-1801-00868}。例如匹配区域IoU为0.8和0.6，另有一个漏报且无多报，则PQ为$1.4/2.5=0.56$；只报告匹配区域的平均IoU会得到0.7，掩盖了漏掉的区域。透明度预测在三分图的未知区内比较连续值，常用绝对误差和$\mathrm{SAD}=\sum_p|\hat\alpha_p-\alpha_p|$与均方误差。

测量任务的指标带单位。计数比较每幅图的预测数$\hat c_i$与真值$c_i$，平均绝对误差为$\mathrm{MAE}=(1/n)\sum_i|\hat c_i-c_i|$，均方根误差为$\mathrm{RMSE}=\sqrt{(1/n)\sum_i(\hat c_i-c_i)^2}$，后者对少数大误差更敏感。二维姿态沿用检测的AP，只是把IoU换成\term{对象关键点相似度}（object keypoint similarity, OKS）：
\begin{equation}
 \mathrm{OKS}=
 \frac{\sum_k\exp\!\left(-d_k^2/(2s^2\kappa_k^2)\right)\mathbf 1\{v_k>0\}}
      {\sum_k\mathbf 1\{v_k>0\}}\text{，}
 \label{eq:vision-oks}
\end{equation}
其中$d_k$是第$k$个关节的预测与标注之间的距离，$s$是该人的尺度（COCO取标注区域面积的平方根），$\kappa_k$是按关节类型设定的容差常数，$v_k>0$表示该关节有标注。OKS把每个关节的像素距离换算成0到1之间的相似度，再在有标注的关节上平均；COCO关键点基准同样在0.50到0.95的10个OKS阈值上平均AP。三维姿态用\term{平均关节位置误差}（mean per-joint position error, MPJPE），即各关节预测位置与真实位置三维距离的平均，单位通常是毫米。深度的常用指标来自早期单目深度工作\scite{vision-1406-02283}：对有效像素集合$\Omega$，绝对相对误差$\mathrm{AbsRel}=(1/|\Omega|)\sum_p|\hat z_p-z_p|/z_p$，阈值准确率$\delta<1.25$是满足$\max(\hat z_p/z_p,z_p/\hat z_p)<1.25$的像素比例，另有线性与对数尺度的RMSE。相机估计比较旋转与平移：两个旋转矩阵$\hat{\bm R},\bm R$之间的角度误差为$\arccos\{[\operatorname{tr}(\hat{\bm R}^\top\bm R)-1]/2\}$，平移在尺度未知时只比较方向夹角。

检索按相关性检查排名。若某个查询有$R$个相关项，排名$k$处的相关指示为$r_k$，截至该位置的精确率为$P(k)$，非插值平均精度为
\begin{equation}
 \mathrm{AP}=\frac1R\sum_k P(k)r_k\text{，}
 \label{eq:vision-retrieval-ap}
\end{equation}
对所有查询取平均即检索的mAP。三个正例出现在第1、3、5位时，AP为$(1+2/3+3/5)/3\approx0.756$。地点识别等任务更关心前几名里有没有正确结果，此时用Recall@$K$：前$K$个结果中至少含一个相关项的查询所占比例。AP既关心是否找到，也关心出现得早不早；Recall@$K$只关心前$K$名，适合候选还要交给后续步骤复核的场合。

有了输入、输出和指标，就可以看清这些任务为什么难。下面五个小节各讨论一类挑战。它们都不只属于某一个任务，但在不同任务上以不同形式出现。

\subsection{外观变化带来的类别歧义}
\label{sec:vision-recognition-challenge-appearance}

同一类别的物体，在图像上可以很不相同；不同类别的物体，又可能很相像。路口的红色货车在阴影里变暗，正对相机时和侧对相机时轮廓完全不同；同一型号的轿车喷成蓝色和红色，像素值几乎没有共同之处。反过来，起重机和鹤在英文里同名，两栋玻璃幕墙建筑在检索中很难分开。识别要求对视角、光照、形变和类内差异保持不变，同时对区分类别或个体的细节保持敏感，这两个要求方向相反。

经典方法用人工设计的特征缓和这一矛盾。\ref{sec:vision-features-template-matching}节的模板匹配直接比较像素，归一化互相关能抵消整体亮度和对比度的变化，却对视角和形变毫无办法。\ref{sec:vision-features-sift}节的SIFT和\ref{sec:vision-features-hog}节的HOG改为比较梯度方向的统计：梯度方向对光照的加性和乘性变化不敏感，局部直方图又容忍小的位移。\ref{sec:vision-features-bovw}节的视觉词袋再把一幅图的局部特征量化成“视觉词”的频率，用支持向量机等分类器判别，也用于在图库中检索同一地点。这些特征的不变性是设计者事先选定的，选定之外的变化（大视角、遮挡、细粒度差异）就无法处理；分类边界也只能在人工特征张成的空间里画。

学习方法把“对什么不变、对什么敏感”交给数据决定。卷积网络和视觉Transformer从大量图像中学习表示，自监督或图文对比训练又让同一表示可以被多个任务复用。代价是表示可能学到数据中的捷径，例如用背景判断类别，或在训练少见的领域失效。\ref{sec:vision-recognition-semantics}节在固定类别表与文字类别表两种条件下讨论类别判定，并把属性限制到对象区域内；\ref{sec:vision-recognition-retrieval}节讨论检索中的同一问题：区分同一实例与相似实例。

\subsection{尺度与遮挡带来的定位困难}
\label{sec:vision-recognition-challenge-localization}

定位要回答“对象在哪里”，而对象在图中的大小和可见程度变化极大。靠近相机的货车占画面三分之一，路口远处的轿车可能只有$20\times10$个像素；公交车站旁的车被行人挡住一半，相邻两车的框大面积重叠。一幅图里的对象数量也不固定，可以是零个，也可以是几十个。定位器必须在所有位置、所有尺度上搜索，又要给出一个不重复、不遗漏的集合。

经典检测器用穷举加剪枝的方式处理。\ref{sec:vision-features-viola-jones}节的Viola–Jones检测器在\ref{sec:vision-transforms-pyramid}节的图像金字塔各层上滑动一个固定大小的窗口，用积分图快速计算类Haar特征，再用级联分类器尽早拒绝明显不是人脸的窗口；尺度变化由金字塔覆盖，海量的背景窗口由级联的前几级过滤。\ref{sec:vision-features-dpm}节的可变形部件模型在HOG金字塔上用一个根滤波器描述整体、若干部件滤波器描述局部，并允许部件相对根位置有代价地移动，从而容忍形变和部分遮挡。两者都在得分图上用非极大值抑制去掉重复窗口。这条路线的问题是：窗口和部件的形状要人工设定，每个尺度要重算特征，被大面积遮挡或比例罕见的对象容易漏掉。

学习式检测器保留了“在所有位置和尺度上打分”的思想，但改变了实现方式。卷积网络在整幅图上一次前向就得到所有位置的预测，相当于把滑动窗口写成卷积；图像金字塔换成网络内部的特征金字塔，各层共享一次计算；级联对容易背景的剪枝，换成了训练时降低容易样本权重的损失；去重从后处理规则逐步变成训练目标。\ref{sec:vision-recognition-detection}节沿着这条线讲六个检测模型，\ref{sec:vision-recognition-mask-sets}节把同样的问题推到实例掩码。

\subsection{稠密输出与分辨率的矛盾}
\label{sec:vision-recognition-challenge-dense}

逐像素的输出要同时满足两个相互冲突的要求。判断道路中央的一块灰色是路面还是车顶，需要看到周围很大范围的上下文；画出路杆、车窗和发丝的细边界，又需要接近原图的分辨率。卷积网络靠逐层下采样扩大感受范围，代价正是丢掉位置细节；如果不下采样，计算和存储又会随分辨率迅速增长。对象掩码还要在像素层面分开相邻的同类对象，透明度预测则要给出连续的混合比例。

经典分割从低层线索把像素组织成区域。\ref{sec:vision-features-otsu}节的阈值法按灰度把像素分为两类，\ref{sec:vision-features-watershed}节的分水岭沿梯度的“山脊”切开区域，\ref{sec:vision-features-graph-segmentation}节的图割把分割写成能量最小化，用像素与标记的颜色相似度和相邻像素的平滑项共同决定边界，\ref{sec:vision-features-superpixel}节的超像素把颜色和位置相近的像素聚成小块。这些方法得到的是边界贴合的区域，却不知道区域属于什么类别；颜色相同的车顶和路面会被并到一起，纹理复杂的同一物体又会被切碎。

学习方法让每个像素的类别由网络在上下文中判断，再设法补回分辨率：空洞卷积在不降低分辨率的前提下扩大感受范围，编码器—解码器用浅层特征恢复边界，查询式分割让每个对象拥有一张整图掩码。经典线索并没有消失，图割、超像素和颜色传播的思想仍出现在交互分割与抠图中。\ref{sec:vision-recognition-segmentation}节到\ref{sec:vision-recognition-matting}节依次讨论像素类别、对象掩码和前景透明度三种稠密输出。

\subsection{单目观测带来的几何不确定}
\label{sec:vision-recognition-challenge-geometry}

一幅照片把三维场景投影到二维平面，深度信息在投影中丢失。\ref{sec:vision-formation-camera-intrinsics}节的针孔模型说明，一个像素只确定一条从相机中心出发的射线：同一像素可以对应4米外的点，也可以对应8米外的点。手腕在图上的位置相同，可能是向前伸也可能是向后伸；把整个场景连同相机平移放大一倍，所有照片保持不变。因此从图像得到深度、三维关节或相机位置，必须引入额外的观测或先验。

经典方法用多视图增加观测。\ref{sec:vision-formation-multi-view}节的双目立体在两台已知基线的相机之间沿极线搜索对应点，由视差反比得到深度；运动恢复结构则对一组无序照片先匹配特征点、估计两视图几何，再三角化三维点，最后用光束法平差联合优化所有相机和点的重投影误差。这些方法有明确的几何依据，在纹理丰富、重叠充分时很准，但依赖可靠的特征匹配：弱纹理、反光、重复图案和视图太少都会使流程中断，而且多步骤串联，前一步的错误会传到后面。单幅图像则无从三角化，经典方法只能借助消失点、已知物体尺寸等少量线索。

学习方法从两个方向改变了这一局面。一是用大规模数据学习单目先验，直接从一幅图预测相对深度或人体三维姿态；二是把多视图几何的计算步骤写成可学习的网络，例如在学习的特征上构造代价体，或让网络一次前馈输出所有相机和点图。前者给出完整但只确定到尺度的结果，后者需要多幅图但能利用跨视图约束。\ref{sec:vision-recognition-measurement}节按测量量依次讨论计数、姿态、深度和相机几何，并在每一处说明尺度从哪里来。

\subsection{开放类别与长尾带来的泛化困难}
\label{sec:vision-recognition-challenge-open}

训练时的类别表总是有限的，使用时却会遇到新东西。仓库画面里出现一个没有名字的搬运装置；零件上出现一种从未见过的划痕；用户要找“公交车旁的红色货车”，这个组合不在任何类别表中。即使在类别表之内，样本数量也极不均衡：轿车有成千上万张，少见车型只有几张。闭集模型对新类别只能二选一地误报成已知类或当作背景，对长尾类别则倾向于预测常见类别。

经典方法基本是闭集的。每个类别单独训练一个检测器或分类器，新类别要重新收集样本并训练；视觉词袋的词汇表在建库时固定，检索时只能比较这些词的频率，再用\ref{sec:vision-recognition-ransac}节的RANSAC几何验证排除外观相似但几何不一致的候选。几何验证不需要类别，这是经典路线中少有的“开放”环节。

学习方法提供了三种扩展方式。第一种用文字指定类别，例如CLIP、Grounding DINO与SAM 3把类别表写成文字，使用时再给出；第二种不命名类别，只判断“像不像一个对象”或“是否偏离正常”，例如OWOBJ的对象性与PatchCore的正常记忆；第三种在没有人工标注时自动构造监督，例如GeM用三维重建挖掘同一建筑的照片对。三者都扩大了可处理的范围，却也带来新的风险：文字条件可能被部分忽略，对象性会把背景提成候选，正常记忆分不清无害的外观变化。这些风险分散在\ref{sec:vision-recognition-semantics}节、\ref{sec:vision-recognition-regions}节和\ref{sec:vision-recognition-retrieval}节中讨论，\ref{sec:vision-recognition-evaluation}节再给出相应的检验协议。

五类挑战按输出逐步变具体的顺序在后续各节中得到回答。\ref{sec:vision-recognition-semantics}节先处理外观歧义下的“是什么”，并在异常判定中第一次面对开放类别；\ref{sec:vision-recognition-regions}节把判断落到框、像素和掩码上，同时回答尺度遮挡与稠密分辨率两类困难，开放词汇检测与未知对象发现也在这里出现；\ref{sec:vision-recognition-measurement}节在定位之上加入数量、关节和三维坐标，正面处理单目观测的几何不确定；\ref{sec:vision-recognition-retrieval}节把一幅图与图库中的其他图对应起来，需要的正是前面形成的全局表示、局部点和几何约束。每一处都先回顾经典路线的做法和失效条件，再讲学习模型改变了哪一步。最后一节把这些输出放在同一组评价条件下检验。

\section{语义与属性判定}
\label{sec:vision-recognition-semantics}

本节处理三种语义输出：类别说明“属于什么”，属性说明“具有什么特征”，异常判断说明“是否偏离规定的正常状态”。输入都是一幅图或一个裁剪区域，输出都是有限个标签或分数，难点却各不相同。类别判定的难点在于类别表从哪里来：训练时固定，还是使用时由文字临时给出。属性判定的难点在于归属：整图里同时有红色和蓝色，必须说清哪个颜色属于哪辆车。异常判定的难点在于缺少反例：缺陷类型无法穷举，训练数据主要是正常样本。

三个难点对应三种不同的设计选择。类别表固定时，可以在冻结的视觉表示上训练一个分类头，DINOv3是这条路线的代表；类别表在使用时才给出时，可以把类别名也编码成向量，用CLIP的图文相似度排序。属性判定沿用这两种接口，但先把表示限制在对象区域内。异常判定改变了监督形式：不学习“缺陷长什么样”，而是保存正常外观，再度量新观测离正常外观有多远。

\subsection{类别与场景的判定}
\label{sec:vision-recognition-classification}

类别判定的输入是整幅图像或对象裁剪，输出是类别表中的一项或多项。在\term{单标签分类}（single-label classification）中，一个样本只对应类别表中的一项；\term{多标签分类}（multi-label classification）允许多项同时成立。例如车辆裁剪图可以按主要车型给一个标签，整张街景却可以同时含公交车、轿车和行人。损失函数必须遵守这一区别：互斥类别用softmax在类别间归一化，可共存的对象则分别输出“是否存在”的得分。场景分类把整幅画面概括为街道或厨房，它与对象分类不是同一个问题，厨房这一场景名不能展开为一份完整的物品清单。

两条路线的差别在于类别表进入计算的方式。第一条路线把类别表写进分类头的输出维度：每个类别对应一行权重，用带标签的样本学习这些权重，视觉编码器可以冻结，也可以微调。第二条路线把类别表写成文字：每个类别对应一条描述，由文字编码器把描述变成向量，再与图像向量比较，不为当前类别表训练新参数。前者需要标注数据，但可以针对具体领域学到细致的判别边界；后者不需要标注，却受限于预训练图文数据覆盖了哪些概念、描述写得是否准确。下面分别用DINOv3和CLIP说明这两种用法。

\subsubsection{DINOv3}
\label{sec:vision-recognition-dinov3}

DINOv3的结构与自监督训练已在\ref{sec:vision-ivm-dinov3}节展开。这里只需要它的几个结论：DINOv3是在大规模无标注图像上自监督训练的视觉Transformer系列，输入图像被切成固定边长的图像块，网络输出一个汇总整图的类别词元和每个图像块各自的特征；训练中加入的格拉姆锚定（Gram anchoring）约束使长时间训练后的逐块特征仍保持局部区分能力。原论文强调，这些表示在多种任务上不经微调即可使用\scite{vision-2508-10104}。本节关心的是拿到这样的编码器后，怎样完成一个有标签的车辆或场景分类任务。

输入的组织先决定特征的形状。设所选检查点的图像块边长为16，输入缩放到$224\times224$，则网络得到$14\times14=196$个块特征和一个类别词元，每个向量为$d$维（$d$由检查点规模决定）。全局任务通常取类别词元输出作为$\bm{z}\in\mathbb{R}^{d}$，也可以把它与块特征的平均拼接；空间任务则保留$14\times14\times d$的块特征网格。缩放、裁剪和归一化必须按检查点的约定进行，并保证标签仍对应被保留的内容：若裁剪切掉了车头，“轿车”标签就可能失去依据。训练、调参和测试图像应按车辆、拍摄批次或采集条件分开，避免同一对象的近重复图同时出现在训练和测试中。

拿到特征后有三种使用方式，差别在于更新哪些参数。\term{线性探测}（linear probing）冻结编码器，只训练一个线性分类头；它把编码器当作固定的特征提取器，所以训练快、显存小，同一份特征还能供多个任务共用。$k$近邻分类连分类头也不训练，直接在带标签的训练特征中找最相近的样本投票。\term{微调}（fine-tuning）则同时更新编码器参数，通常使用比分类头小得多的学习率，或只解冻最后几层。设冻结编码器$f_{\bm{\theta}}$得到$\bm{z}=f_{\bm{\theta}}(\bm{X})$，分类头包含$\bm{W}\in\mathbb{R}^{k\times d}$与$\bm{b}\in\mathbb{R}^{k}$。对$k$个互斥类别，计算
\begin{equation}
 p(y=r\mid\bm{X})=
 \frac{\exp(\bm{w}_r^\top\bm{z}+b_r)}
 {\sum_{j=1}^{k}\exp(\bm{w}_j^\top\bm{z}+b_j)}\text{，}
 \label{eq:vision-linear-head}
\end{equation}
其中$\bm{w}_r^\top$是$\bm{W}$的第$r$行。训练目标为平均负对数似然$L=-n^{-1}\sum_i\log p(y_i\mid\bm X_i)$。线性探测只更新$\bm{W},\bm{b}$；微调还更新$\bm{\theta}$，所需梯度存储随之增加，训练数据较少时还可能破坏预训练表示中对当前任务有用的结构。

对一个样本，记类别得分为$a_r=\bm{w}_r^\top\bm{z}+b_r$、概率为$p_r$，则$\partial(-\log p_y)/\partial a_r=p_r-\mathbf{1}\{r=y\}$，头部梯度为$(p_r-\mathbf{1}\{r=y\})\bm{z}$。假设三类车型的得分为$(2,1,0)$，softmax概率依次约为0.665、0.245和0.090；若真实类别为第二类，得分梯度约为$(0.665,-0.755,0.090)$，下降更新提高第二类相对于其他类的得分。这些概率是在给定类别集合内归一化的结果，不能证明图像中一定存在可辨车辆。

选择哪种方式，取决于数据量、领域差距和资源。训练样本少、目标领域与预训练图像接近时，线性探测通常已足够，而且不会改动编码器；领域差距大（例如显微图像或遥感图像），或任务依赖预训练中少见的细节时，微调可能带来改进，但需要更多标注和计算，并要用验证集确认改进确实存在。调参时需要固定的超参数包括输入分辨率、所用特征层、头部学习率和训练轮数；微调还要加上编码器学习率和解冻层数。空间任务可以在每个块特征上接同一个线性头，得到$14\times14$的类别图，这正是\ref{sec:vision-recognition-dense-pixels}节像素类别预测的最简形式。

常见失败有三种。第一，若编码器读取整幅街景，判定左车车型的得分可能依赖道路背景；任务要针对某一辆车时，应先取得该车区域。第二，外观接近的细粒度类别依赖小部位：全局背景足以区分拍摄地点，却可能不足以区分鸟喙或叶缘。iNaturalist包含细粒度类别和不均衡样本，评价时应同时检查小部位分辨率和类别覆盖\scite{vision-1707-06642}。例如常见类别有90张且全对，少见类别有10张且全错，样本平均准确率为90\%，两个类别的平均准确率却为50\%。第三，使用时必须保留训练时的编码器版本、预处理、分类头和类别顺序；需要“未知”输出时，再加入由验证集确定的拒绝阈值。

\subsubsection{CLIP}
\label{sec:vision-recognition-clip}

如果使用时才给出“轿车、货车、公交车”，固定输出维度的分类头不能直接容纳新类别。CLIP的双编码器和对比训练已在\ref{sec:vision-models-contrastive}节展开，其结论是：图像编码器和文字编码器把匹配的图文对映射到单位球面上相近的位置，两者的内积可以比较图像与任意一条描述的对应程度\scite{radford2021}。\term{零样本分类}（zero-shot classification）利用这一点，把类别表写成文字，不为当前类别表训练分类头。

输入的组织分两侧。文字侧，把每个类别名填入模板，例如“一张{类别}的照片”，编码后归一化；同一类别可以使用多个模板，对其单位文字向量求平均后再归一化，CLIP原文在ImageNet上即对多条模板做了这种集成。这样得到的$k$个文字向量排成矩阵，起到线性分类头权重的作用，只是权重来自文字而非标注。图像侧，按检查点约定预处理后得到单位图像向量。令$\bar{\bm{z}}$与$\bar{\bm{t}}_r$分别为单位图像向量和第$r$类的文字向量，$s_r=\bar{\bm{z}}^\top\bar{\bm{t}}_r$，使用温度$\tau>0$计算候选内归一化得分
\begin{equation}
 q_r=\frac{\exp(s_r/\tau)}
 {\sum_{j\in C}\exp(s_j/\tau)}\text{，}
 \label{eq:vision-clip-classification}
\end{equation}
其中$C$是这次输入的非空候选索引集合。CLIP训练时把$1/\tau$作为可学习参数，并把它的上限设为100。分母随候选集合变化，所以$q_r$不是不依赖候选的存在概率。模板集合须在评价前固定，换模板就是换分类器。

CLIP默认不更新任何参数。若有少量标注，可以在冻结的CLIP图像向量上做线性探测，这时CLIP退化为第一条路线中的一个编码器；也可以用文字向量初始化分类头再微调。超参数主要是模板集合、类别名的写法和拒绝阈值。类别名有歧义时（例如“crane”既是鹤也是起重机），模板需要加入上下文。

设四项相似度为$0.31,0.28,0.19,0.12$，最后一项是不在图中的“帆船”。最大值只说明第一项在这四项中领先；把前三项全部换成不相关的词，模型仍会返回一个最大值。开放使用要另设含未知对象的验证协议、拒绝阈值或目标存在判断。整图相似度还无法保证“红色”绑定到货车而非旁边的轿车，也不擅长回答“图中有几辆车”这类计数问题。

两种用法可以在同一车型分类任务、同一测试集上比较，差别集中在类别来源、训练需求和失败方式上，见表\ref{tab:vision-recognition-classification-compare}。CLIP把类别表的指定从训练阶段推迟到使用阶段，换来了不需标注的灵活性，代价是判别边界由图文预训练决定，难以针对领域细节调整；DINOv3加分类头则反过来。二者并不互斥：CLIP图像编码器同样可以做线性探测，DINOv3的块特征也可以接入文字对齐的模块。

\begin{table}[htbp]
 \centering
 \small
 \begin{tabularx}{\linewidth}{@{}lXX@{}}
 \toprule
 比较项 & DINOv3加分类头 & CLIP零样本分类\\
 \midrule
 类别来源 & 训练时固定，写入头部输出维度 & 使用时以文字给出\\
 需要的数据 & 每类带标签的样本 & 无需当前任务标注\\
 更新的参数 & 分类头，或连同编码器微调 & 默认不更新\\
 空间输出 & 块特征可接逐位置分类头 & 主要为整图向量\\
 典型失败 & 背景捷径、细粒度部位不足 & 候选闭合、属性错绑、模板敏感\\
 \bottomrule
 \end{tabularx}
 \caption{同一车型分类任务下两种类别接入方式的比较。表中列出的是接口差别，不是两个模型在某一数据集上的精度结论。}
 \label{tab:vision-recognition-classification-compare}
\end{table}

\subsection{对象属性的判定}
\label{sec:vision-recognition-attributes}

要记录“红色货车”和“蓝色轿车”，只输出整图属性集合$\{\text{红色},\text{蓝色}\}$会丢掉归属。属性判定的输入因而不是整幅图，而是一个对象区域；输出是$(\text{对象编号},\text{属性},\text{取值})$记录。难点在于把属性证据限制在目标对象内：裁剪过小会切掉车灯，过大又会纳入旁车颜色。

路线有两条。一条是整图多标签预测，实现简单，但无法表达归属，只适合“图中是否出现红色物体”这类问题。另一条先定位对象，再对各区域提取表示并接属性头；区域表示可以来自DINOv3块特征在框内的汇总，也可以来自\ref{sec:vision-recognition-mask-rcnn}节的区域对齐特征。CLIP也可以用“一辆红色的车”作为文字查询，但应当对对象裁剪而不是整图计算相似度，否则同样会出现错绑。位置与属性共享对象编号，后续计数、搜索或人工核验才能指向同一个对象。

属性之间是否互斥由标签定义决定。“红色车身”和“车灯开启”可以同时成立，“完整车身的主色”则可以规定为单标签字段。对于可共存的$k$个二值属性，令区域向量$\bm{z}_i$经过每属性一个线性得分$h_{i,r}$，$p_{i,r}=1/(1+e^{-h_{i,r}})$为预测，$a_{i,r}\in\{0,1\}$为标注，$m_{i,r}$表示该属性是否有可靠标注，则
\begin{equation}
 L_{\mathrm{attr}}=
 -\frac{\sum_{i,r}m_{i,r}
 [a_{i,r}\log p_{i,r}+(1-a_{i,r})\log(1-p_{i,r})]}
 {\sum_{i,r}m_{i,r}}\text{。}
 \label{eq:vision-attribute-loss}
\end{equation}
分母要求至少有一个有效标注，全无标注的批次不计算此项。掩码$m_{i,r}$的作用是区分“未标注”和“不存在”：把未标注直接置为不存在，会训练出系统性的假负例。单个有效标注对得分的梯度是$p_{i,r}-a_{i,r}$。冻结区域编码器时只训练属性头；各属性阈值由独立验证记录选择，输出保留每车的属性分数和可见性。若主色被定义为互斥字段，则另用softmax头，不把多个二元输出强行当成唯一主色。

属性所需的部位不可见时，应输出“不可判定”而不是猜一个值。对象编号并不能修正错误的视觉证据，但它使属性错绑能够被单独定位和统计。

\subsection{正常与异常状态的判定}
\label{sec:vision-recognition-anomaly}

工业零件的输入是一幅零件图像，输出是整图的异常分数和标出可疑位置的异常热图。难点在于监督不对称：新的缺陷类型不断出现，训练集主要是正常样本，无法像车型分类那样为每种缺陷准备标签。\term{异常检测}（anomaly detection）因此改为学习或保存正常外观，再度量新观测与正常外观的偏离。

不同路线对“偏离”给出不同的度量。重建路线训练网络重建正常图像，用重建误差作为异常分数；它的困难是网络泛化过好时也能重建缺陷，泛化不足时又会把正常的细纹理重建模糊。分布路线为正常特征拟合一个密度，用低密度表示异常；单个高斯等简单分布可能无法覆盖正常外观中的多种模式。记忆路线保存正常样本的局部特征，用最近邻距离表示异常，不需要假设分布形式，代价是记忆的规模和检索时间。PatchCore属于第三条路线\scite{vision-2106-08265}。

PatchCore针对的具体问题是：整图特征比较会把一处细小划痕淹没在整幅零件的正常外观中，而直接用最深层特征又过于偏向ImageNet类别语义。它的主要想法是保存正常图像中\emph{每个局部位置}的中层特征，查询时逐位置找最近的正常片段；局部比较让细小划痕得到单独的分数，中层特征兼顾纹理细节和一定的上下文。原工作默认使用ImageNet预训练的WideResNet-50，取第2、3个残差阶段的特征。对每个位置，先在$3\times3$邻域内做自适应平均池化，把邻近上下文汇入该位置；再把第3阶段的特征双线性放大到第2阶段的空间尺寸并与之合并，得到每个位置一个固定维的片段向量。原文把图像缩放到$256\times256$后中心裁剪为$224\times224$，不做数据增强。WideResNet-50第2阶段的步幅为8，因此第2阶段的特征网格是$28\times28$，每幅正常图贡献784个片段向量；第3阶段步幅为16，网格$14\times14$，放大两倍后与第2阶段逐位置对齐。合并后的片段维数由自适应平均池化统一到一个预先设定的值，原文未给出这个数值。全部正常图的片段汇成特征集合$\bm{M}$。

PatchCore没有梯度训练，编码器始终冻结，“学习”就是建立记忆。全部片段可能非常多，因而从$\bm{M}$中保留代表子集$\bm{M}_{\mathrm{c}}$，即\term{核心集}（coreset），其覆盖目标为
\begin{equation}
 \min_{\substack{\bm{M}_{\mathrm{c}}\subseteq\bm{M}\\
 |\bm{M}_{\mathrm{c}}|=m}}
 \max_{\bm{z}\in\bm{M}}\min_{\bm{u}\in\bm{M}_{\mathrm{c}}}
 \lVert\bm{z}-\bm{u}\rVert_2\text{。}
 \label{eq:vision-patchcore-coreset}
\end{equation}
最内层最小值为每个正常片段寻找一个代表，外层最大值检查最难覆盖的正常片段。若只最小化平均距离，少量但正常的外观模式可能在记忆中完全缺席，查询时就会被误报为异常。精确求解需要穷举子集，原方法改用近似贪心选择：从一个起点开始，反复加入离当前子集最远的特征，并为每个正常特征维护到当前子集的最近距离；新加入$\bm{u}$后更新为原距离与$\lVert\bm{z}-\bm{u}\rVert_2$的较小者，达到预算$m$就停止。原工作用随机投影降低选择时的距离计算成本，并报告了保留25\%、10\%和1\%的设置。建立记忆的数据必须符合“正常”定义，异常混入记忆会把某些缺陷变成近邻参照。

查询图像位置$p$的局部分数为$d_p=\min_{\bm{u}\in\bm{M}_{\mathrm{c}}}\lVert\bm{z}_p-\bm{u}\rVert_2$。把这些分数按空间排列、双线性放大到输入分辨率并用$\sigma=4$的高斯核平滑，就得到异常热图。整图分数从最大局部距离$s_*=\max_p d_p$开始，再根据最异常片段$\bm{z}_*$的最近正常片段$\bm{u}_*$周围是否稠密重新加权。令$\bm{N}_b(\bm{u}_*)$包含$\bm{u}_*$在记忆中的$b$个近邻，则
\begin{equation}
 s=\left(1-
 \frac{\exp\lVert\bm{z}_*-\bm{u}_*\rVert_2}
 {\sum_{\bm{u}\in\bm{N}_b(\bm{u}_*)}
  \exp\lVert\bm{z}_*-\bm{u}\rVert_2}\right)s_*\text{。}
 \label{eq:vision-patchcore-score}
\end{equation}
邻域大小$b$须大于1，否则分子分母相同，分数退化为零。若最近正常片段较孤立，其他近邻的距离更大，括号中的系数更接近1，最大距离受到的压低更少；也就是说，靠近一个稀少正常模式的查询会被判得更可疑。

假设一维正常记忆为$\{0,1,4\}$，查询局部特征为$\{0.2,0.9,6\}$，对应距离是$0.2,0.1,2$，最大值由第三个片段给出。若$b=2$，取最近正常片段4及其近邻1，整图分数为$2[1-\exp(2)/(\exp(2)+\exp(5))]\approx1.905$。阈值应在独立验证条件下确定。若新设备使所有特征整体平移，正常零件也可能远离记忆；高分只能表示对当前正常参照的偏离，不能直接命名缺陷，也不自动等于业务不合格。

\section{对象定位与区域划分}
\label{sec:vision-recognition-regions}

上一节的判断都没有说明位置。定位要把判断接回图上，输入是图像以及“要找什么”的指定（固定类别表、一段文字或一个点框提示），输出按精细程度分为四种形式：边界框给出对象的大致矩形范围，像素类别图给出每个像素属于哪一类，对象掩码进一步把同类像素分给不同实例，前景透明度则为发丝等混合边界给出连续的贡献比例。四种输出回答同一个“在哪里”的问题，但保存的信息逐级增加，监督和计算代价也逐级增加。本节按输出形式划分小节，每个小节内再比较不同模型对同一输出的不同解法。

各小节共有两个难点。其一是目标数量不定：一幅图里可能有零辆、一辆或二十辆车，网络却只能输出固定形状的张量，因此需要规定哪些输出位置负责哪些目标，以及怎样去掉重复。其二是目标由谁指定：类别表固定时可以把类别写进输出维度，文字或点击指定目标时，需要让查询条件进入定位计算。下面各模型的差别，大多就是对这两个问题的不同回答。

\subsection{对象边界框的预测}
\label{sec:vision-recognition-detection}

\term{目标检测}（object detection）的输入是一幅图像，输出是若干$(\bm{b},c,s)$组成的集合，其中$\bm{b}$是框，$c$是类别，$s$是得分。\term{边界框}（bounding box）用一个矩形包围目标，以下采用连续边界$\bm{b}=(x_1,y_1,x_2,y_2)$，宽高分别为$x_2-x_1$和$y_2-y_1$。标签应说明框包围可见部分还是完整对象。两个框的重叠程度用式\eqref{eq:vision-iou}的IoU衡量。节首例子中IoU为0.6的两个框，若是对同一辆车的两个预测，应当去重；若属于两辆互相遮挡的车，直接去掉一个则会漏检。

在学习式检测器之前，经典路线已经确立了检测的基本框架：在所有位置和尺度上打分，再去掉重复。\ref{sec:vision-features-viola-jones}节的Viola–Jones检测器以$24\times24$左右的固定窗口在图像金字塔的每一层上滑动，每个窗口用积分图在常数时间内算出类Haar特征，再交给由AdaBoost训练的级联分类器；前几级只用少数特征就拒绝绝大多数背景窗口，只有少量窗口走完全部级联。它在正面人脸这类外观稳定的目标上实现了实时检测，但类Haar特征描述不了外观变化大的类别。\ref{sec:vision-features-hog}节的HOG加线性支持向量机把特征换成梯度方向直方图，对行人这类轮廓明确的目标更稳健；\ref{sec:vision-features-dpm}节的可变形部件模型再在HOG金字塔上组合一个粗分辨率的根滤波器和若干细分辨率的部件滤波器，部件位置可以偏离锚点、偏离量受二次代价约束，部件位置作为隐变量用隐支持向量机训练。三者的共同结构是“特征金字塔$\to$逐窗口打分$\to$非极大值抑制”，差别在特征和打分函数。它们的局限也在这三处：每个尺度都要重算人工特征，窗口和部件的形状要人为设定，每个类别单独训练一个检测器。

下面的学习式检测器保留了这套结构，但逐处替换。YOLO与RetinaNet用卷积网络在整幅图上一次前向得到所有位置的得分，滑动窗口变成卷积的逐位置输出，人工特征变成学习的特征，多个类别共用一个网络；RetinaNet还把图像金字塔换成网络内部的特征金字塔，各层共享主干计算，并用焦点损失降低大量容易背景的训练权重，在作用上对应级联对容易窗口的提前拒绝。可变形部件的“允许局部偏移”在卷积网络中没有显式对应，而是由多层非线性特征隐式吸收。去重则直到DETR才从非极大值抑制移进训练目标。

检测的难点集中在“哪个输出负责哪个目标”。本小节的六个模型可以看成对这个问题的一串回答，每一步只改动一个设计。YOLO把图像划成网格，让目标中心所在的格子负责它，用一次前向得到所有框，但每个格子容量有限，重复框要靠后处理去掉。RetinaNet在多尺度特征的每个位置放置一组预设的参考框（锚框），让与目标重叠足够大的锚框负责它；这样小目标和拥挤目标有了更多候选，代价是绝大多数锚框是背景，需要焦点损失平衡训练。FCOS去掉预设锚框，让落在目标内的特征位置直接预测到四条边的距离，减少了锚框尺寸和比例这类人工设定，但仍需规则分配和非极大值抑制。DETR进一步取消密集候选，用固定数量的查询向量和二分匹配直接输出一个无重复的集合，把去重写进了训练目标。Grounding DINO沿用DETR式的查询解码，但把类别表换成任意文字，使“要找什么”在使用时给出。OWOBJ则面对连名字都没有的对象，在检测器上增加对象性信号，使未知对象也能作为候选保留下来。

\subsubsection{YOLO}
\label{sec:vision-recognition-yolo}

YOLO之前的主流检测器先生成上千个候选区域，再对每个区域分别做分类，同一幅图要经过多次网络计算，速度难以满足实时视频。原始YOLO的想法是把检测写成一次回归：整幅图经过一个网络，直接输出所有格子的框和类别得分\scite{vision-1506-02640}。由于每个格子都在整图特征上作出预测，网络在判断时能看到全局上下文。

网络是24个卷积层加2个全连接层，逐层结构见表\ref{tab:vision-recognition-yolo-layers}。它借鉴GoogLeNet的思路，但不用Inception模块，而是在$3\times3$卷积前反复插入$1\times1$卷积：$1\times1$卷积先把通道压缩一半，$3\times3$卷积再在较少的通道上扩大感受范围，这样在加深网络的同时控制计算量。以第4组为例，单个子模块的计算顺序是“$1\times1$卷积降到256通道$\to$带泄漏的ReLU$\to3\times3$卷积回到512通道$\to$带泄漏的ReLU”，重复4次。网络用4个$2\times2$最大池化和2个步长为2的卷积把$448\times448$的输入缩小64倍，得到$7\times7\times1024$的特征；第一个全连接层把它展平后映射到4096维，其后接比例为0.5的随机失活，第二个全连接层输出$7\times7\times30$。除最后一层是线性输出外，各层都用斜率0.1的带泄漏的ReLU，即$x>0$时取$x$、否则取$0.1x$。原文还给出一个只有9个卷积层、滤波器更少的Fast YOLO，具体通道数论文未给出。

\begin{table*}[tbp]
 \centering
 \small
 \begin{tabular}{@{}lllrrl@{}}
 \toprule
 阶段或层组 & 层类型 & 卷积核与步长 & 输出通道 & 重复次数 & 输出尺寸\\
 \midrule
 输入 & — & — & 3 & — & $448\times448$\\
 第1组 & 卷积 & $7\times7$，步长2 & 64 & 1 & $224\times224$\\
       & 最大池化 & $2\times2$，步长2 & 64 & 1 & $112\times112$\\
 第2组 & 卷积 & $3\times3$ & 192 & 1 & $112\times112$\\
       & 最大池化 & $2\times2$，步长2 & 192 & 1 & $56\times56$\\
 第3组 & 卷积 & $1\times1$；$3\times3$ & 128；256 & 1 & $56\times56$\\
       & 卷积 & $1\times1$；$3\times3$ & 256；512 & 1 & $56\times56$\\
       & 最大池化 & $2\times2$，步长2 & 512 & 1 & $28\times28$\\
 第4组 & 卷积 & $1\times1$；$3\times3$ & 256；512 & 4 & $28\times28$\\
       & 卷积 & $1\times1$；$3\times3$ & 512；1024 & 1 & $28\times28$\\
       & 最大池化 & $2\times2$，步长2 & 1024 & 1 & $14\times14$\\
 第5组 & 卷积 & $1\times1$；$3\times3$ & 512；1024 & 2 & $14\times14$\\
       & 卷积 & $3\times3$ & 1024 & 1 & $14\times14$\\
       & 卷积 & $3\times3$，步长2 & 1024 & 1 & $7\times7$\\
 第6组 & 卷积 & $3\times3$ & 1024 & 2 & $7\times7$\\
 全连接 & 全连接 & — & 4096 & 1 & $1\times1$\\
 输出 & 全连接，重排 & — & 30 & 1 & $7\times7$\\
 \bottomrule
 \end{tabular}
 \caption{YOLO（第一版）的逐层结构。层配置取自原文图3，卷积共24层（第1至6组依次为1、1、4、10、6、2层）；输出尺寸按$448\times448$输入逐层推算，原文只在图中画出各组的张量形状。未注明步长的卷积步长为1并保持尺寸。“$1\times1$；$3\times3$”表示依次相连的两层，重复次数指这一对的重复次数。ImageNet预训练使用前20个卷积层，检测时再加上后4个卷积层和2个全连接层。}
 \label{tab:vision-recognition-yolo-layers}
\end{table*}

全连接层输出一个$S\times S\times(5B+C)$的张量；原文取$S=7$、每格$B=2$个框、$C=20$个类别，输出为$7\times7\times30$。每个格子的30个数中，前10个是两个框各自的$(u,v,w,h)$和置信度，后20个是该格子共享的条件类别概率$\mathbb{P}(c\mid\mathrm{obj})$。框中心$(u,v)$相对于格子定位，宽高相对于整图归一化。若格子列行为$(j,i)$，输入宽高为$(W_{\mathrm{in}},H_{\mathrm{in}})$，则中心还原为$((j+u)W_{\mathrm{in}}/S,(i+v)H_{\mathrm{in}}/S)$。

训练按“目标中心落在哪个格子，就由该格子负责”组织监督；在负责格子的$B$个框中，再选当前与真实框IoU最高的那个承担坐标监督。记格子为$a$、框槽为$b$，$I^{\mathrm{resp}}_{a,b}$标出负责框，$I^{\mathrm{other}}_{a,b}$标出其他框槽，$I^{\mathrm{obj}}_a$标出含目标中心的格子。预测中心、宽高和置信度分别为$(\hat x,\hat y)$、$(\hat w,\hat h)$和$\hat C$，真实值去掉帽子。原始目标全部用平方误差，先写位置项：
\begin{equation}
 \begin{aligned}
 e_{a,b}^{\mathrm{pos}}&=(\hat x_{a,b}-x_a)^2
                       +(\hat y_{a,b}-y_a)^2,\\
 e_{a,b}^{\mathrm{size}}&=(\sqrt{\hat w_{a,b}}-\sqrt{w_a})^2
                        +(\sqrt{\hat h_{a,b}}-\sqrt{h_a})^2,\\
 L_{\mathrm{coord}}&=\lambda_{\mathrm c}\sum_{a,b}
 I^{\mathrm{resp}}_{a,b}(e_{a,b}^{\mathrm{pos}}+e_{a,b}^{\mathrm{size}})\text{。}
 \end{aligned}
 \label{eq:vision-yolo-training-process}
\end{equation}
置信度和类别项为
\[
 \begin{aligned}
 L_{\mathrm{conf}}&=\sum_{a,b}I^{\mathrm{resp}}_{a,b}
 (\hat C_{a,b}-\operatorname{IoU}_{a,b})^2
 +\lambda_{\mathrm n}\sum_{a,b}
 I^{\mathrm{other}}_{a,b}\hat C_{a,b}^{\,2},\\
 L_{\mathrm{class}}&=\sum_a I^{\mathrm{obj}}_a\sum_r
 (\hat p_a(r)-\mathbf{1}\{r=c_a\})^2\text{。}
 \end{aligned}
\]
总目标是三项之和。原文取$\lambda_{\mathrm c}=5$、$\lambda_{\mathrm n}=0.5$：大多数框槽不含目标，若不降低它们的权重，“没有目标”的信号会压倒少数负责框的定位信号。宽高取平方根，使同样的像素误差在小框上产生更大的代价。训练分两个阶段：先在ImageNet上以$224\times224$输入预训练前20个卷积层，再加上其余卷积和全连接层，以$448\times448$输入用检测数据训练全部参数。

推理时，每框置信度表示$\mathbb{P}(\mathrm{obj})\operatorname{IoU}$，乘上共享的$\mathbb{P}(c\mid\mathrm{obj})$得到类别相关得分。例如格子$(2,3)$的中心偏移$(0.5,0.25)$还原为$(160,208)$；归一化宽高$(0.2,0.1)$对应像素宽高$(89.6,44.8)$；置信度0.8与货车条件概率0.75相乘，货车框分数为0.6。输入曾缩放或加边时，还须去掉加边并除缩放比例，才是原图上的框。然后按阈值过滤，再用\term{非极大值抑制}（non-maximum suppression, NMS）去重：按得分从高到低取一个框，删除同类中与它IoU超过阈值的候选，对余下框重复。对前述IoU为0.6、得分0.9与0.7的两个框，阈值0.5会删除第二个。

原始YOLO的限制来自它的分配方式。一个格子只预测一组类别概率和两个框，两辆车中心落在同一格时至少会漏掉一辆；7\ensuremath{\times}7的粗网格对远处小车尤其不利。NMS在拥挤场景中还可能删掉真实存在的相邻目标。后来的YOLO版本改变了预测、分配与去重设计，不能沿用此处全部细节。

\subsubsection{RetinaNet}
\label{sec:vision-recognition-fpn}
\label{sec:vision-recognition-retinanet}

YOLO的单一粗网格留下了两个问题：小目标在深层特征中可能已经消失，每个位置的候选又太少。RetinaNet从两处入手：用特征金字塔在多个尺度上预测，并在每个位置放置多个参考框。但候选一多，训练中背景候选就以数量压倒前景，RetinaNet为此提出焦点损失，使单阶段密集检测在训练上不再被大量容易背景主导\scite{vision-1708-02002}。

多尺度特征来自\term{特征金字塔网络}（feature pyramid network, FPN），它是RetinaNet的组成部分\scite{vision-1612-03144}。远处车辆在图上只占几个像素，深层特征步幅大，位置细节丢失；浅层特征保留位置，却缺少语义。经典检测器的做法是\ref{sec:vision-transforms-pyramid}节的图像金字塔：把输入缩放成多个尺寸，每个尺寸分别计算特征，各层语义水平相同，代价是特征计算重复多次。FPN改为只运行一次主干，利用卷积网络本身逐层下采样形成的特征层级，并补上一条自顶向下的路径，使高分辨率的层也带上深层语义。具体做法是把深层特征上采样，与相应浅层特征相加。记主干第$\ell$阶段输出为$\bm{C}_\ell$，自顶向下的中间特征为$\bm{M}_\ell$，则
$\bm{M}_\ell=\operatorname{Conv}_{1\times1}(\bm{C}_\ell)+\operatorname{Up}(\bm{M}_{\ell+1})$，再用$3\times3$卷积得到$\bm{P}_\ell$。$1\times1$横向卷积把通道数统一为256，最近邻上采样把空间尺寸放大两倍，否则两项无法相加。这些附加卷积层都不带非线性。RetinaNet使用$\bm{P}_3$到$\bm{P}_7$五层，$\bm{P}_\ell$的分辨率是输入的$1/2^\ell$，步幅从8到128。其中$\bm{P}_3$至$\bm{P}_5$由主干的$\bm{C}_3$至$\bm{C}_5$经上述自顶向下路径得到；$\bm{P}_6$由$\bm{C}_5$经一个步长为2的$3\times3$卷积得到，$\bm{P}_7$再对$\bm{P}_6$施加ReLU和一个步长为2的$3\times3$卷积；不使用步幅为4的$\bm{P}_2$，以减少计算。较密的层负责小目标，较疏的层负责大目标；若目标在输入缩小时已经丢失，特征融合不能补回观测。

整个检测器按主干、颈部和头部分段，各段的输入输出形状见表\ref{tab:vision-recognition-retinanet-structure}。主干的残差块结构在\ref{sec:vision-ivm-resnet}节已经讲过，这里只用到它各阶段输出的步幅和通道数。

\begin{table*}[tbp]
 \centering
 \small
 \begin{tabularx}{\linewidth}{@{}llXl@{}}
 \toprule
 部分 & 输入形状 & 结构与层数 & 输出形状\\
 \midrule
 主干ResNet-50 & $800\times800\times3$ & 去掉全连接层，取第3至5阶段输出$\bm C_3,\bm C_4,\bm C_5$ & $100^2\times512$，$50^2\times1024$，$25^2\times2048$\\
 颈部FPN & $\bm C_3$至$\bm C_5$ & 每层$1\times1$横向卷积到256通道，最近邻上采样两倍后相加，再经$3\times3$卷积；$\bm P_6$、$\bm P_7$各由一个步长2的$3\times3$卷积得到 & $\bm P_3$至$\bm P_7$：$100^2,50^2,25^2,13^2,7^2$，均为256通道\\
 锚框 & 每层每个位置 & 面积$32^2$（$\bm P_3$）至$512^2$（$\bm P_7$），每层尺度$\{2^0,2^{1/3},2^{2/3}\}$乘长宽比$\{1{:}2,1{:}1,2{:}1\}$，$A=9$ & 覆盖约32至813像素\\
 分类子网 & 每层$h\times w\times256$ & 4个$3\times3$卷积（256通道，各接ReLU），再接$3\times3$卷积与sigmoid；各层共享参数 & $h\times w\times KA$\\
 回归子网 & 每层$h\times w\times256$ & 结构同分类子网但参数独立，末层线性输出；与类别无关 & $h\times w\times4A$\\
 \bottomrule
 \end{tabularx}
 \caption{RetinaNet（ResNet-50-FPN）的分段结构。层数、通道与锚框配置取自原文；输出尺寸按$800\times800$输入推算（原文消融实验多用短边600像素，主结果用短边800像素）。$K$为类别数，$A$为每个位置的锚框数。}
 \label{tab:vision-recognition-retinanet-structure}
\end{table*}

在金字塔每个位置，RetinaNet放置$A=9$个\term{锚框}（anchor box），即3种尺度乘3种长宽比的参考框。两个小型全卷积子网在所有层间共享：分类子网经4层$3\times3$卷积后输出$KA$个通道，每个锚框对每个类别给一个独立sigmoid得分；回归子网结构相同，输出$4A$个通道，给出相对于锚框的修正。两个子网参数不共享，回归子网对所有类别给出同一组修正。设锚框为$(x_a,y_a,w_a,h_a)$，目标为$(x_*,y_*,w_*,h_*)$，回归标签为$((x_*-x_a)/w_a,(y_*-y_a)/h_a,\log(w_*/w_a),\log(h_*/h_a))$；给定预测修正$\bm{\delta}$，逆变换为
\begin{equation}
 \begin{aligned}
 \hat x&=x_a+w_a\delta_x,\qquad \hat y=y_a+h_a\delta_y,\\
 \hat w&=w_a e^{\delta_w},\qquad \hat h=h_a e^{\delta_h}\text{。}
 \end{aligned}
 \label{eq:vision-retinanet-decode}
\end{equation}
中心差除以参考尺寸、宽高比取对数，使标签无像素单位且能表达乘性缩放。教学锚框$(50,30,40,20)$配目标$(60,30,80,20)$，标签为$(0.25,0,\log2,0)$，逆算恢复原目标。

训练按锚框与真实框的IoU分配：IoU不低于0.5的锚框为前景，低于0.4为背景，介于两者之间的忽略。一幅图约有十万个锚框，绝大多数是很容易判断的背景。每个背景锚框的交叉熵很小，但数量巨大，总和仍可能主导梯度。令$p_t$为模型分给真实二值标签的概率，\term{焦点损失}（focal loss）为
\begin{equation}
 L_{\mathrm{focal}}=-\alpha_t(1-p_t)^\gamma\log p_t\text{，}
 \qquad\gamma\geq0\text{。}
 \label{eq:vision-focal-loss}
\end{equation}
其中$\alpha_t$调整正负样本权重，原文取$\gamma=2$、$\alpha=0.25$。$\gamma=0$时退回加权交叉熵；$\gamma=2$时，$p_t=0.99$的容易样本乘上$10^{-4}$，而$p_t=0.5$的困难样本乘上$0.25$。总目标是所有锚框的焦点损失加上正锚框的平滑$L_1$回归损失，再除以正锚框数。分类子网最后一层的偏置按先验前景概率0.01初始化，避免训练初期背景损失过大而不稳定。主干从ImageNet分类权重初始化，金字塔与子网由检测监督训练；原文用16张图的小批量做9万次随机梯度下降迭代。推理时先以0.05的得分阈值过滤，每层最多保留得分最高的1000个锚框并按式\eqref{eq:vision-retinanet-decode}解码，各层合并后以IoU阈值0.5做NMS，最后逆变换到原图。

焦点损失改变的是每个候选的训练贡献，不增加少见类别的观测，因此前景背景失衡与类别长尾仍须分别检查。锚框尺寸与比例是人工设定，与数据中目标形状不匹配时会降低召回，这正是下一个模型要去掉的设计。

\subsubsection{FCOS}
\label{sec:vision-recognition-fcos}

锚框带来了一组需要逐数据集调整的超参数：尺寸、比例、正负IoU阈值；它们还使每个位置的输出数量乘以$A$。FCOS的想法是让特征位置本身充当参考点：落在目标框内的位置直接预测它到框四条边的距离$(l,t,r,b)$，不再以锚框为坐标起点\scite{vision-1904-01355}。

FCOS沿用RetinaNet式的FPN（$\bm{P}_3$至$\bm{P}_7$，步幅8至128）和共享检测头，结构上只有三处改动。第一，$\bm{P}_6$、$\bm{P}_7$改由$\bm{P}_5$（而非$\bm{C}_5$）经步长为2的卷积逐级得到。第二，头部仍是分类与回归两个分支，各含4个256通道的$3\times3$卷积，新加的卷积层后使用组归一化（最后的预测层除外）；分类分支末端输出$h\times w\times K$，回归分支末端输出$h\times w\times4$。第三，与分类输出并列增加一个单层的中心度分支，输出$h\times w\times1$（原文脚注指出把它移到回归分支效果更好，但默认结构仍放在分类分支旁）。于是每个位置只输出一组预测：$K$个类别得分、4个距离和1个中心度，而RetinaNet每个位置要输出$A=9$组。步幅为$s$的特征位置$(i,j)$对应原图点$(x,y)=(\lfloor s/2\rfloor+js,\lfloor s/2\rfloor+is)$，预测框为$(x-l,y-t,x+r,y+b)$。各层共享同一个头，而不同层负责的距离范围相差很大，因此距离分支的输出$x$经$\exp(s_\ell x)$映射，$s_\ell$是每层一个可训练标量：指数保证距离为正，$s_\ell$让共享的头在不同层上自动调整输出量级。默认主干为ResNet-50，训练输入短边800像素、长边不超过1333像素。

监督仍需分配，只是规则换成了几何条件。落在某个真实框内的位置是正例候选；原方法再按层限定可回归的最大距离范围（$\bm{P}_3$至$\bm{P}_6$的上界依次为64、128、256、512像素，$\bm{P}_7$不设上界），把不同大小的目标分到不同层；同一位置落在多个框内时选面积最小的框。训练目标包括类别的焦点损失、正例的IoU损失和中心度的二元交叉熵。\term{中心度}（centerness）衡量位置是否接近框中心，其目标为
\begin{equation}
 c_*=\sqrt{\frac{\min(l,r)}{\max(l,r)}
                 \frac{\min(t,b)}{\max(t,b)}}\text{。}
 \label{eq:vision-fcos-centerness}
\end{equation}
中心点四侧距离均衡时$c_*=1$，靠近边界时其中一个比例趋于零。对于框$(20,10,80,50)$，中心$(50,30)$的距离为$(30,20,30,20)$，中心度为1；位置$(25,30)$的距离为$(5,20,55,20)$，中心度为$\sqrt{1/11}\approx0.302$。图\ref{fig:vision-box-points}把这两个位置与框画在同一坐标系中：两点都可以精确回归同一矩形，中心度却不同，因此它不是预测框IoU的直接测量。

\begin{figure}[htbp]
 \centering
 % 连续像素坐标下的FCOS教学例；所有长度与正文算例一致。
\begin{tikzpicture}[x=1mm,y=-1mm,font=\figurefont\small,
  dist/.style={-{Stealth[length=2mm]},draw=BookBlue,line width=0.9pt}]
 \draw[draw=BookMuted,line width=0.7pt,->] (10,0)--(90,0) node[right] {$x$};
 \draw[draw=BookMuted,line width=0.7pt,->] (10,0)--(10,58) node[below] {$y$};
 \draw[draw=BookTeal,fill=BookTeal!7,line width=1pt] (20,10) rectangle (80,50);
 \draw[draw=BookMuted,dashed,line width=0.7pt] (50,10)--(50,50);
 \draw[draw=BookMuted,dashed,line width=0.7pt] (20,30)--(80,30);
 \fill[BookInk] (50,30) circle (1mm);
 \fill[BookBlue] (25,30) circle (1mm);
 \node[anchor=north west] at (51,33) {中心：$c_*=1$};
 \node[anchor=south west] at (26,27) {边缘：$c_*\approx0.302$};
 \draw[dist] (25,18)--(20,18);
 \draw[dist] (25,18)--(80,18);
 \node[anchor=south] at (22.5,17) {$5$};
 \node[anchor=south] at (57,17) {$55$};
 \draw[draw=BookBlue,line width=0.7pt] (25,18)--(25,29);
 \node[anchor=north] at (20,51) {$(20,50)$};
 \node[anchor=south] at (80,9) {$(80,10)$};
\end{tikzpicture}

 \caption{同一真实框内的中心位置与边缘位置。实线箭头标出边缘位置到左右边界的距离，虚线标出中心轴；两点都可以回归同一矩形，但中心度目标分别为1与约0.302。坐标采用向右、向下为正的连续像素边界。}
 \label{fig:vision-box-points}
\end{figure}

推理时，类别得分乘以预测中心度作为排序分数，降低远离中心的低质量框，再做NMS。FCOS减少了锚框相关的超参数，每个位置的输出也少了$A$倍；但分配规则、尺度范围和NMS仍在，拥挤场景中的重复与漏检问题没有因去掉锚框而消失。

\subsubsection{DETR}
\label{sec:vision-recognition-detr}

前三个模型都先产生密集候选，再用NMS去重。NMS是训练之外的手工规则：阈值高了留下重复框，低了删掉相邻的真目标。DETR的想法是把检测直接视为集合预测：网络输出固定数量的预测，训练时用二分匹配让每个真实目标恰好对应一个预测，其余预测学习“无对象”；这样重复预测在训练中就受到惩罚，推理时不再需要NMS\scite{vision-2005-12872}。

DETR的结构由三段组成。主干（原文为ResNet-50）把$H\times W\times3$的图像变成$H/32\times W/32\times2048$的特征，$1\times1$卷积降到$d=256$通道，展平为$HW/1024$个词元并加正弦位置编码。6层Transformer编码器在这些位置之间交换上下文。解码器读入$N=100$个\term{对象查询}（object query），即可学习的$d$维向量，经6层自注意力和对图像特征的交叉注意力更新；查询之间的自注意力让各查询“知道”其他查询在预测什么，这是避免重复的结构基础。每个查询最后经共享的线性层输出$K+1$类的softmax（多出的一类是“无对象”），经共享的三层前馈网络输出归一化的中心宽高$(c_x,c_y,w,h)\in[0,1]^4$。查询是学习得到的向量，并非预先指定给某辆车的身份，预测序列也没有必须对应目标编号的顺序。Transformer结构本身可参考\BookChapterRef{chap:transformer}{卷五的Transformer章}。

DETR的Transformer部分沿用原始的后置层归一化写法。记编码器第$l$层的输入为$\bm Z^{(l)}\in\mathbb R^{n\times256}$，位置编码为$\bm E$，一个编码层依次计算
\[
 \begin{aligned}
 \bm Z'&=\operatorname{LN}\bigl(\bm Z^{(l)}+\operatorname{MSA}(\bm Z^{(l)}+\bm E,\ \bm Z^{(l)}+\bm E,\ \bm Z^{(l)})\bigr),\\
 \bm Z^{(l+1)}&=\operatorname{LN}\bigl(\bm Z'+\operatorname{FFN}(\bm Z')\bigr)\text{，}
 \end{aligned}
\]
其中$\operatorname{MSA}$的三个参数依次是查询、键和值：位置编码只加在查询和键上，值不含位置。这样位置只影响“读哪里”，不直接混入被读取的内容。解码器的输入在第一层初始化为零，对象查询作为“输出位置编码”在每个注意力层加到查询上；每个解码层在查询自注意力与前馈网络之间多一次交叉注意力，其键为编码器输出加空间位置编码，值为编码器输出。前馈网络是两层带ReLU的逐位置变换，原文未给出其中间宽度。表\ref{tab:vision-recognition-detr-config}按配置表的格式列出原文的四个规模。

\begin{table*}[tbp]
 \centering
 \small
 \begin{tabular}{@{}llccccl@{}}
 \toprule
 模型 & 主干与特征步幅 & 词元数 & 编码／解码层数 & 宽度 & 头数 & 参数量\\
 \midrule
 DETR & ResNet-50，32 & $HW/32^2$ & 6／6 & 256 & 8 & 41M\\
 DETR-DC5 & ResNet-50，16 & $HW/16^2$ & 6／6 & 256 & 8 & 41M\\
 DETR-R101 & ResNet-101，32 & $HW/32^2$ & 6／6 & 256 & 8 & 60M\\
 DETR-DC5-R101 & ResNet-101，16 & $HW/16^2$ & 6／6 & 256 & 8 & 60M\\
 \bottomrule
 \end{tabular}
 \caption{DETR的配置。DETR不切图像块，词元来自主干末层特征的每个位置，因此“图像块大小”一列换成特征步幅；DC5表示把主干最后一阶段的步幅去掉、改用空洞卷积，使特征分辨率加倍。对象查询数均为100，随机失活比例0.1；前馈网络的中间宽度论文未给出。原文说明DETR（R50）的4130万参数中，主干占2350万，Transformer占1780万。以$800\times1216$的输入为例，步幅32给出$25\times38=950$个词元，步幅16则为3800个。}
 \label{tab:vision-recognition-detr-config}
\end{table*}

训练的关键是匹配。标注有$m\leq N$个对象，第$i$个为$(c_i,\bm{b}_i)$，第$j$个预测为$(\hat{\bm{p}}_j,\hat{\bm{b}}_j)$。DETR寻找从真实对象到不同预测的单射$\sigma$，使总匹配代价最小：
\begin{equation}
 \hat{\sigma}=\argmin_\sigma
 \sum_{i=1}^{m}\left[
 -\hat p_{\sigma(i)}(c_i)
 +\lambda_1\lVert\bm{b}_i-\hat{\bm{b}}_{\sigma(i)}\rVert_1
 +\lambda_{\mathrm{g}}L_{\mathrm{GIoU}}
    (\bm{b}_i,\hat{\bm{b}}_{\sigma(i)})\right]\text{。}
 \label{eq:vision-detr-matching}
\end{equation}
这就是\term{二分匹配}（bipartite matching）：一侧是真实对象，一侧是预测，每条边的权重是匹配代价，要求每个真实对象恰好配一个不同的预测。它可以用匈牙利算法在多项式时间内精确求解。原文取$\lambda_1=5$、$\lambda_{\mathrm g}=2$。广义交并比$\operatorname{GIoU}=\operatorname{IoU}-|E\setminus(\bm b\cup\bm b')|/|E|$，其中$E$是两框的最小包围矩形，定位损失为$L_{\mathrm{GIoU}}=1-\operatorname{GIoU}$；不相交框的IoU都为零，额外的空白比例仍能区分隔得近和隔得远的框。

假设两个真实对象对三个查询的总代价矩阵为
\[
 \bm{C}=\begin{pmatrix}1&2&9\\1.2&5&8\end{pmatrix}\text{。}
\]
逐行独立取最小都会选第一个查询，违反唯一对应。二分匹配选择第一对象到第二查询、第二对象到第一查询，总代价为$3.2$，小于反过来分配的$6$。改变标注的排列只改变匹配的写法，不改变最佳集合代价。

匹配选定后视为本轮的固定监督，不对离散匹配本身求梯度。匹配的查询承担负对数分类损失和两项框损失，未匹配的查询学习“无对象”，该类的权重降为0.1，以免大量空查询压倒少数目标。每个解码层的输出都接同样的预测头并计算辅助损失，帮助深层解码器训练。主干从ImageNet权重初始化并冻结批归一化统计，其余部分由检测数据训练；优化器为AdamW，Transformer学习率$10^{-4}$、主干$10^{-5}$。基线模型训练300轮，在16块V100上约需3天；与Faster R-CNN比较时延长到500轮。推理时直接丢弃“无对象”和低分输出，不执行NMS。

DETR把去重从后处理移进了训练目标，代价也随之而来：全局注意力的计算随特征位置数平方增长，难以使用高分辨率特征，小目标效果较弱；每个查询初期没有空间先验，匹配在训练早期不稳定，导致收敛很慢；查询数$N$还限定了一次可输出的对象上限。后续的可变形注意力和以锚点初始化查询的改进主要针对这几点，下面的Grounding DINO即建立在这一系列改进之上。

\subsubsection{Grounding DINO}
\label{sec:vision-recognition-grounding-dino}

前四个模型的类别表都写在分类层的输出维度里。当目标变成“红色货车”或“公交车旁的货车”时，固定类别头无法表达；先用固定检测器找出所有车、再用CLIP对每个裁剪打分，又会把文字与视觉的交互推迟到定位之后，定位本身无法利用文字条件。Grounding DINO的想法是让文字从一开始就参与检测：文字与图像在特征阶段相互融合，用与文字最相关的图像位置初始化查询，解码时查询同时读取图像和文字，最后把类别分类换成查询与文字词元的对齐分数\scite{vision-2303-05499}。它的检测部分沿用DINO（一种DETR改进型检测器），因此集合匹配与无NMS输出都继承自DETR。

输入是一幅图像和一段文字。文字可以是用句点分隔的类别名列表，例如“轿车. 货车. 公交车.”，也可以是一个指代表达。文字经BERT-base编码为至多256个词元特征$\bm{T}\in\mathbb{R}^{n_t\times d}$；图像经Swin Transformer主干得到多尺度特征，展平为$\bm{V}\in\mathbb{R}^{n_v\times d}$。为了防止并列的类别名相互干扰，原文采用子句级文字表示：注意力掩码阻止不相关的类别名彼此注意，每个类别仍保留各自的逐词特征。

模型由三段计算组成。第一段是特征增强器，共6层，每层对图像特征做可变形自注意力、对文字特征做自注意力，再做图像到文字和文字到图像两个方向的交叉注意力，使车辆附近的图像特征带上“货车”的语义，文字特征也带上图中实际出现的内容。第二段是语言引导的查询选择：对每个图像位置$i$计算它与所有词元的最大相似度$\max_r\bm{v}_i^\top\bm{t}_r$，取得分最高的900个位置，用它们的位置初始化查询的参考框，查询的内容部分仍是可学习向量。此时选中的只是候选位置，不能当成已经确认的货车。第三段是6层跨模态解码器，每层依次做查询自注意力、对图像特征的可变形交叉注意力、对文字特征的交叉注意力和前馈网络，逐层修正框。输出是900个查询各自的归一化框，以及一个$900\times n_t$的对齐logit矩阵，其中第$(i,r)$项是查询$i$与词元$r$的点积。表\ref{tab:vision-recognition-gdino-structure}按分段格式列出这些组件。

\begin{table*}[tbp]
 \centering
 \small
 \begin{tabularx}{\linewidth}{@{}llXl@{}}
 \toprule
 部分 & 输入形状 & 结构与层数 & 输出形状\\
 \midrule
 图像主干 & $H\times W\times3$ & Swin-T（T版）或Swin-L（L版）；T版取步幅8至32的三个尺度，再把步幅32的特征下采样得到步幅64，共4个尺度；L版直接取步幅4至32的4个尺度 & 4个尺度，展平为$n_v\times256$\\
 文字主干 & 至多256个词元 & BERT-base，子句级注意力掩码 & $n_t\times256$\\
 特征增强器 & 图像与文字特征 & 6层；每层图像可变形自注意力、文字自注意力、图像到文字与文字到图像交叉注意力、前馈网络 & 形状不变\\
 查询选择 & 增强后的特征 & 按与词元的最大相似度取前900个图像位置初始化参考框 & 900个查询\\
 跨模态解码器 & 900个查询 & 6层；每层查询自注意力、图像可变形交叉注意力、文字交叉注意力、前馈网络 & 900个框与$900\times n_t$对齐logit\\
 \bottomrule
 \end{tabularx}
 \caption{Grounding DINO的分段结构，隐藏维数$d=256$。原文给出的总参数量为T版1.72亿、L版3.41亿。}
 \label{tab:vision-recognition-gdino-structure}
\end{table*}

训练数据混合了三类：带类别框的检测数据（类别名拼成文字输入）、带短语与框对应的定位数据，以及由已有模型为图文对产生伪框的图像描述数据。每条训练记录包含图像、文字、目标框以及框与文字中词元的对应，一个对象可以对应多个词元。训练沿用DETR式二分匹配，匹配代价中分类、$L_1$与广义IoU三项的权重为2.0、5.0和2.0。匹配后对每个查询—词元logit施加焦点损失，使匹配查询与其对应词元的分数高、与其他词元的分数低；框用$L_1$和广义IoU约束，三项损失权重依次为1.0、5.0和2.0，每个解码层和编码器输出都计算辅助损失。T版的预训练数据为Objects365，或再加上GoldG（Flickr30k Entities与Visual Genome的短语定位数据）和约四百万对伪标注图文数据；L版还加入了OpenImages、COCO和指代表达数据。图像和文字主干从各自的预训练权重初始化，融合与解码部分在上述数据上联合训练。使用时不再为每次查询更新权重。

推理时，对每个查询取它在各词元上的最大对齐分数作为框分数，保留超过框阈值的查询；再在该查询的词元分数中找出超过文字阈值的词元，拼回短语作为标签；最后把归一化中心宽高乘以原图尺寸并转换为四边。框阈值、文字阈值和短语恢复规则都会改变最终输出，评价时须固定。若模型只在“货车”词元上高分而“红色”得分弱，最终记录不能悄悄把整个属性短语写成已核实的事实。

能接收“红色货车”“公交车旁的货车”这类字符串，只说明接口容纳了条件，不能证明属性和关系已被正确解析。可以分别遮去参照物、更换颜色词或保留多个同类目标，检查预测是否按所需证据改变。完整的表达定位协议在\ref{chap:vision-multimodal-understanding}章讨论。

\subsubsection{OWOBJ}
\label{sec:vision-recognition-unknown-discovery}
\label{sec:vision-recognition-owobj}

仓库画面中出现一个尚未加入类别表的搬运装置。固定类别检测器可能把它认作货车，也可能当成背景；Grounding DINO允许用户写出装置名称，但用户此时未必知道应当查询什么。\emph{未知对象发现}（unknown object discovery）要求在没有名字时先提出“这里可能有一个对象”，保留区域供后续核验。它与开放词汇定位的差别在于名称是否已作为输入；它与拒识也不同，拒识只回答已有裁剪是否属于已知类，发现还要确定未知对象在哪里。

问题的根源在训练分配。DETR式训练中，没有匹配到已知标注的查询一律学习“无对象”，而这些查询中既有真实背景，也有恰好未被标注的对象；把它们全部用零监督，检测器就学会把一切未标注的东西判为背景。OWOBJ的想法是给检测器加一个\emph{对象性}（objectness）信号，即一个候选是否像一个独立对象，而不是它属于哪个已知类；未匹配查询不再一律接受零标签，而是按其特征与已知对象特征分布的接近程度获得软标签\scite{vision-owobj-2025}。

OWOBJ不另造检测器，而是作为插件接入已有的Faster R-CNN式或DETR式检测器，推理流程不变。在开放世界检测实验中，它接在两个以可变形DETR为基础的方法之上，主干为ResNet-101；少样本检测实验接在Faster R-CNN式的DeFRCN上，开放词汇检测实验接在以CLIP为主干的CORA上。以可变形DETR式的实现为例，主干给出多尺度特征，经可变形编码器和解码器，得到每个查询的嵌入$\bm{q}_i\in\mathbb{R}^d$。查询之后接三个分支：框分支是3层前馈网络；分类分支输出$Z+1$维，其中$Z$维对应已知类，多出的一维表示未知或对象性；对象性模块本身不含可学习的层，只维护查询嵌入的一个高斯分布，用来给多出的那一维生成软标签。查询数$N$和嵌入维数$d$沿用所接的基线检测器，论文正文未给出具体数值。用集合$C_{\mathrm{k}}$表示训练已标注类别，输出可以是已知记录$(\bm{b},c,s)$，也可以是未命名记录$(\bm{b},\mathrm{unknown},s_{\mathrm{o}})$。原方法由匹配查询特征的移动平均估计均值$\bm{\mu}$与正定对角协方差$\bm{\Sigma}$，用
\begin{equation}
 \begin{aligned}
 d_{\mathrm{M}}(\bm{q})=
 (\bm{q}-\bm{\mu})^\top\bm{\Sigma}^{-1}(\bm{q}-\bm{\mu}),\qquad
 s_{\mathrm{o}}(\bm{q})=\exp[-d_{\mathrm{M}}(\bm{q})]
 \end{aligned}
 \label{eq:vision-owobj-objectness}
\end{equation}
生成对象性软标签。\term{马氏距离}（Mahalanobis distance）按各特征方向的正常变动尺度比较偏离，避免数值尺度大的通道独占距离。$s_{\mathrm{o}}$在距离为零时为1，距离增大时下降；它是训练设计中的对象性信号，不是已经校准的真实对象概率。

训练在原检测器的类别与框损失之外加入三类约束。匹配查询继续接受已知类别和框监督，并通过减小$d_{\mathrm M}$使对象特征在分布中更集中；未匹配查询用$s_{\mathrm{o}}$监督对象性输出。原工作还用变分近似的先验匹配项约束分布，高斯先验的方差随估计后验调整并加入方差偏置$\beta^2$（原文消融中$\beta=0.5$最好），以处理少标注条件下的分布收缩；已知与未匹配查询之间的能量间隔项$\max(0,E_{\mathrm u}-E_{\mathrm k}+\delta)$（$\delta=0.2$）则改变分类得分的相对约束，这里的能量来自分类logit的对数指数和。移动统计按$\bm\mu^{(t+1)}=\rho\bm\mu^{(t)}+(1-\rho)\bar{\bm q}$更新，$0\leq\rho<1$控制旧批次的保留程度，对角方差同样更新并保持正数下界。训练从所接检测器的预训练权重出发；使用时沿用该检测器的框和对象性输出筛选未知候选，式\eqref{eq:vision-owobj-objectness}只在训练中提供软标签，不能拿新图的真实标注重新估计对象分布。

作为一维教学例子，若$\mu=0$、$\Sigma=1$，两个未匹配查询为0.2和2，则软对象性分别为$e^{-0.04}\approx0.961$和$e^{-4}\approx0.018$。给前者较强的对象监督，可以避免直接把未知候选抹掉；但若背景特征也靠近这个分布，它同样会被提出。因此对象性需要同时检查未知召回和背景误检，不能由“保留更多区域”推断发现已经正确。

发现之后可以暂存原图编号、框、得分和当前类别表，再由人工或后续证据命名新类，形成新增训练记录。把新类加入检测器属于\term{类别增量学习}（class-incremental learning），后续学习既要扩展输出集合，又要保留旧类能力，可参考\BookChapterRef{chap:ket-update-maintenance}{卷八“知识的更新和维护”章}；它不是这一次未命名检测已经完成的步骤。仓库装置被误报为已知货车、墙面被误报为未知对象、加入装置类后货车退化，是三种不同的错误。

六个模型的对比见表\ref{tab:vision-recognition-detection-compare}，比较固定在同一街景检测任务和同一输出形式$(\bm{b},c,s)$上。从YOLO到FCOS，改变的是候选的密度和参考方式，去重始终依靠NMS；DETR把去重移进训练，换来更简洁的推理，付出的是收敛速度和小目标；Grounding DINO与OWOBJ都保留DETR式集合输出，分别改变了类别从哪里来和未标注对象怎样监督。

\begin{table*}[tbp]
 \centering
 \small
 \begin{tabularx}{\linewidth}{@{}lXXXX@{}}
 \toprule
 模型 & 负责目标的输出 & 类别来源 & 去重方式 & 主要代价\\
 \midrule
 YOLO & 中心所在格子的框槽 & 固定类别表 & NMS & 粗网格，同格多目标漏检\\
 RetinaNet & 与目标IoU足够大的多尺度锚框 & 固定类别表 & NMS & 锚框设定，背景失衡需焦点损失\\
 FCOS & 落在框内且尺度相符的位置 & 固定类别表 & 中心度加NMS & 仍依赖分配规则与NMS\\
 DETR & 二分匹配选中的查询 & 固定类别表加无对象 & 训练中的唯一匹配 & 收敛慢，小目标弱，查询数上限\\
 Grounding DINO & 匹配查询，与词元对齐 & 使用时给出的文字 & 唯一匹配 & 阈值与短语恢复影响输出\\
 OWOBJ & 匹配查询加对象性分支 & 已知类加未命名对象 & 沿用所接检测器 & 未知召回与背景误检需权衡\\
 \bottomrule
 \end{tabularx}
 \caption{同一检测任务下六个模型在关键设计上的差别。每一行相对上一行只改动一到两项设计；表中不给出精度数值，因为原论文的数据集与设置各不相同。}
 \label{tab:vision-recognition-detection-compare}
\end{table*}

\subsection{像素类别的预测}
\label{sec:vision-recognition-segmentation}
\label{sec:vision-recognition-dense-pixels}

框内常含背景，也可能含另一辆车；需要准确范围时，就要给每个像素一个判断。\term{语义分割}（semantic segmentation）给每个像素分配类别，同类车辆可以连成一片；\term{实例分割}（instance segmentation）给不同车辆各自的掩码；\term{全景分割}（panoptic segmentation）同时覆盖道路、天空等不可数区域，使每个像素获得类别，可数对象的像素再获得实例编号。两辆车的并集可以是正确的车辆语义区域，却是错误的单个实例，因此三种标签不能未经说明相互转换。图\ref{fig:vision-segmentation-outputs}沿用同一个离散街景，显示增加实例编号与背景类别各带来什么信息。本小节讨论只输出类别、不区分实例的像素预测；区分实例的掩码留给下一小节。

\begin{figure}[htbp]
 \centering
 % 同一离散场景的三种分割输出；字符在灰度下保留类别和身份区别。
\begin{tikzpicture}[font=\figurefont\small,x=4.5mm,y=-4.5mm]
 \foreach \offset/\title in {0/语义分割,8/实例分割,16/全景分割}{
  \begin{scope}[xshift=\offset*4.5mm]
   \node[anchor=south] at (3,0) {\title};
   \fill[BookPaper] (0,0.4) rectangle (6,4.4);
   \draw[BookMuted,line width=0.7pt] (0,0.4) rectangle (6,4.4);
   \draw[BookRule,line width=0.5pt] (0,1.4)--(6,1.4);
   \ifnum\offset=8
    \node[text=BookMuted] at (3,3.6) {不标背景};
   \else
    \node[text=BookMuted] at (3,3.6) {道路};
    \node[text=BookMuted] at (3,0.9) {天空};
   \fi
   \filldraw[fill=BookTeal!18,draw=BookTeal,line width=0.8pt]
      (0.5,1.8) rectangle (2.8,2.8);
   \ifnum\offset=0
    \filldraw[fill=BookTeal!18,draw=BookTeal,line width=0.8pt]
      (3.2,1.8) rectangle (5.5,2.8);
    \node at (1.65,2.3) {车};
    \node at (4.35,2.3) {车};
   \else
    \filldraw[fill=BookAmber!16,draw=BookAmber,dashed,line width=0.8pt]
      (3.2,1.8) rectangle (5.5,2.8);
    \node at (1.65,2.3) {车1};
    \node at (4.35,2.3) {车2};
   \fi
  \end{scope}
 }
\end{tikzpicture}

 \caption{三种分割结果的区别。语义分割保留类别，实例分割分开可数对象，全景分割共同保留背景类别与对象编号；文字编号和边界线型使身份区别在灰度打印中仍可辨认。图中区域为离散教学构造，框状图元表示教学掩码，不是检测框。}
 \label{fig:vision-segmentation-outputs}
\end{figure}

像素类别预测的输入是$H\times W\times3$的图像，输出是$H\times W\times K$的类别得分图，逐像素取最大即得类别图。难点在于分辨率与上下文之间的矛盾：判断道路中央的一块灰色像素是路面还是车顶，需要很大范围的上下文；画出路杆和车窗的细边界，又需要接近原图的分辨率。卷积网络靠逐层下采样扩大感受范围，代价正是丢失位置细节，直接把最后的低分辨率特征放大只会增加采样位置，补不回细边界。

经典分割方法走的是另一条路：它们在原分辨率上工作，边界由像素之间的低层差异决定，但不判断类别。\ref{sec:vision-features-otsu}节的Otsu阈值法在灰度直方图上选一个使类间方差最大的门限，把像素分成两类，只在前景和背景灰度分布分得开时有效。\ref{sec:vision-features-watershed}节的分水岭把梯度幅值看成地形，从各个极小值“注水”，水域相遇处就是边界；它得到闭合轮廓，却对噪声敏感，容易过分割，通常要先用标记控制注水的起点。\ref{sec:vision-features-graph-segmentation}节的图割把每个像素看成图的顶点，一元项衡量像素颜色与前景、背景颜色模型的匹配程度，成对项惩罚相似的相邻像素被分开；对二值标记，这一能量的全局最小值可由最小割求得，GrabCut再用用户给的框迭代更新颜色模型。\ref{sec:vision-features-superpixel}节的SLIC超像素在颜色和位置的联合空间里做局部$k$均值，把图像切成大小相近、贴合边界的小块。这些方法共同的长处是边界精细、不需要训练；共同的短处是不知道“路面”和“车顶”是两类：颜色相近的不同类别会被合并，纹理复杂的同一物体会被切碎。学习方法补上的正是语义，而经典方法的边界线索仍以别的形式保留下来：DeepLab系列的早期版本曾在网络输出之后接一个带成对项的条件随机场来贴合边界，交互分割和抠图也继续使用图割与颜色传播的思想。

DeepLabv3+对这一矛盾的回答是在编码端用空洞卷积扩大上下文而不再继续下采样，在解码端引入浅层特征恢复边界。BIT处理的是同一矛盾的另一种形态：输入变成同一地点两个时期的影像，类别变成“变化／不变”，上下文不仅来自同一幅图的周围，还来自另一时期的对应区域；它用少量语义词元在两期之间交换上下文。两者的输出形式相同，都是逐像素类别图，差别在于类别的含义和上下文的来源。

\subsubsection{DeepLabv3+}
\label{sec:vision-recognition-deeplab}

在DeepLabv3+之前，DeepLabv3已经用空洞卷积和多尺度池化取得了较强的上下文，但它直接把输出步幅为16的得分图双线性放大16倍，物体边界因而粗糙；另一类编码器—解码器网络逐级恢复分辨率，边界较好，上下文却受编码器下采样限制。DeepLabv3+把两者结合：以DeepLabv3作为编码器提供上下文，再加一个简单的解码器，用浅层特征修正边界\scite{vision-1802-02611}。

\term{空洞卷积}（atrous convolution）在卷积核的相邻采样点之间留出间隔，空洞率$r$表示间隔为$r$。例如一维三点核读取$x_{i-r},x_i,x_{i+r}$，参数个数仍为三，覆盖范围却从3扩大到$2r+1$。它让网络在不降低特征分辨率、不增加参数的前提下读到更大范围。

网络按编码器与解码器分段，各段形状见表\ref{tab:vision-recognition-deeplab-structure}。主干（ResNet-101或改造的Xception）把最后几个阶段的下采样换成空洞卷积，输出步幅为16，即$H/16\times W/16\times2048$的特征。\term{空洞空间金字塔池化}（atrous spatial pyramid pooling, ASPP）模块在这一特征上并行使用五个分支：一个$1\times1$卷积、三个空洞率分别为6、12、18的$3\times3$卷积（原文图1），以及一个整图平均池化分支；各分支输出拼接后经$1\times1$卷积得到256通道的编码器输出。不同空洞率对应不同大小的邻域，整图池化提供全局场景信息。以输出步幅16、空洞率18为例，一个$3\times3$核的两个最远采样点相隔$2\times18=36$个特征位置，约合输入上的576像素。解码器先把编码器输出双线性放大4倍到$H/4\times W/4$；从主干步幅为4的浅层取出特征（ResNet-101取第2阶段残差块在步长2之前的最后一层），用$1\times1$卷积压到48通道，以免浅层通道数过多压过语义特征——原文在8至64通道之间比较，48通道最好；两者拼接后经两个256通道的$3\times3$卷积细化（一个或三个都稍差），再用$1\times1$卷积输出$K$个类别得分，最后双线性放大4倍回到输入尺寸。

\begin{table*}[tbp]
 \centering
 \small
 \begin{tabularx}{\linewidth}{@{}llXl@{}}
 \toprule
 部分 & 输入形状 & 结构与层数 & 输出形状\\
 \midrule
 主干 & $513\times513\times3$ & ResNet-101，最后一个阶段改用空洞卷积，输出步幅16；或改造的对齐Xception & 浅层$129\times129\times256$；深层$33\times33\times2048$\\
 ASPP & $33\times33\times2048$ & $1\times1$卷积；空洞率6、12、18的$3\times3$卷积；整图池化；拼接后$1\times1$卷积 & $33\times33\times256$\\
 解码器：上采样 & $33\times33\times256$ & 双线性放大4倍 & $129\times129\times256$\\
 解码器：浅层分支 & $129\times129\times256$ & $1\times1$卷积 & $129\times129\times48$\\
 解码器：融合 & $129\times129\times304$ & 拼接后2个$3\times3$卷积（256通道），$1\times1$卷积输出类别 & $129\times129\times K$\\
 输出 & $129\times129\times K$ & 双线性放大4倍 & $513\times513\times K$\\
 \bottomrule
 \end{tabularx}
 \caption{DeepLabv3+的分段结构。输出步幅、通道数和解码器层数取自原文；输出尺寸按原文训练裁剪$513\times513$推算，ResNet-101各阶段通道数沿用\ref{sec:vision-ivm-resnet}节。原文以ResNet-101为主干时，不带解码器的计算量约为810亿次乘加，加上解码器后约为1013亿次；参数量论文未给出。}
 \label{tab:vision-recognition-deeplab-structure}
\end{table*}

改造的Xception有三处变化：中间流的模块重复16次（原文图4），使网络更深；所有最大池化换成带步长的深度可分离卷积，使任意位置都可以改用空洞卷积来保持分辨率；每个$3\times3$深度卷积之后都加批归一化和ReLU。\term{深度可分离卷积}（depthwise separable convolution）把一个标准卷积拆成两步：先对每个通道单独做$3\times3$卷积，再用$1\times1$卷积混合通道，参数和计算量因此大幅下降。原文在ASPP和解码器中也改用深度可分离卷积，乘加次数下降约三至四成，精度基本不变。

训练使用逐像素交叉熵。对训练图$i$的有效位置集$\Omega_i$，类别标注为$y_{i,p}$，按类别softmax得到$p_{i,p,r}$，则
\begin{equation}
 L_{\mathrm{pixel}}=
 -\frac{\sum_i\sum_{p\in\Omega_i}\log p_{i,p,y_{i,p}}}
        {\sum_i|\Omega_i|}\text{。}
 \label{eq:vision-dense-pixel-training}
\end{equation}
无效标签的像素不进入$\Omega_i$。每个位置的得分梯度是预测类别分布减真实类别指示，再经上采样、解码器和ASPP传回主干。主干从分类权重初始化，ASPP和解码器由像素监督学习。缩放、裁剪图像时，标签须用最近邻等保持类别的采样，不能用插值制造不存在的类别编号。推理时先把类别得分图插值回原图，再逐像素取最大；直接插值整数类别编号没有概率含义。训练与测试可以使用不同的输出步幅（例如测试用8以换取更细的结果），这改变了计算量和精度，应共同记录。

道路内部的类别主要依赖上下文，细路杆依赖分辨率；若路杆在输入缩小时已经消失，解码器的浅层特征也无法补回。两车相邻时，即使边缘清楚，纯类别图也不分配实例身份。

\subsubsection{BIT}
\label{sec:vision-recognition-bit}

遥感应用中的目标常常是找出新增的建筑，而不是给单幅影像中的每栋楼分类。\term{变化检测}（change detection）的输入是同一地区两个时期、已经配准的影像，输出是逐像素的“变化／不变”图。常见的孪生卷积网络分别提取两期特征再比较，但卷积的感受范围有限，难以利用远处的上下文判断一块区域的变化是否有意义；对两期全部像素做自注意力，计算又随像素数平方增长。BIT（bitemporal image transformer）的想法是先把每期特征压缩成少量语义词元，在这个小空间里用Transformer建模两期之间的上下文，再把结果投回像素\scite{vision-2103-00208}。

两期影像$\bm{X}_1,\bm{X}_2$经共享权重的ResNet-18主干得到特征图，展平为$\bm{F}_1,\bm{F}_2\in\mathbb{R}^{n\times d}$，$n$是空间位置数。语义分词器为每期计算空间注意力$\bm{A}\in\mathbb{R}^{n\times k}$，每列在位置上做softmax归一化，得到$k$个加权汇总的词元$\bm{T}=\bm{A}^{\top}\bm{F}\in\mathbb{R}^{k\times d}$；原文取$k=4$。每个词元可以理解为“建筑区域”“道路区域”这类高层概念在该期影像中的汇总。两期词元拼成$2k$个，加位置编码后经Transformer编码器相互注意，使一期的概念能参考另一期的对应内容。编码后的词元再拆回两期，各自经Transformer解码器投回像素：每个像素特征作为查询，读取本期的$k$个词元，得到上下文增强的特征$\bm{F}_1',\bm{F}_2'$。最后把两期特征上采样到原图尺寸，计算$|\bm{F}_1'-\bm{F}_2'|$，由浅层卷积预测头输出两类得分。

表\ref{tab:vision-recognition-bit-structure}给出原文主模型的分段形状。ResNet-18原有5个阶段、各下采样2倍；BIT把后面阶段的步长改为1，主模型只保留前4个阶段，接一个输出32通道的逐点卷积压缩特征，再双线性放大，使特征的下采样倍数为4。Transformer很小：词元数$L=4$，编码器1层，解码器8层，每层8个头、每头8维，前馈网络为两层、中间宽度64并用GELU，各子层前置层归一化；只给词元加可学习位置编码，解码器只有交叉注意力和前馈网络，不在像素之间做自注意力。预测头是两个带批归一化的$3\times3$卷积，输出通道依次为32和2。整个模型只有约355万参数，$256\times256$输入下约43.5亿次浮点运算；同样主干、不加Transformer的基线为338万参数，可见语义词元带来的额外开销很小。

\begin{table*}[tbp]
 \centering
 \small
 \begin{tabularx}{\linewidth}{@{}llXl@{}}
 \toprule
 部分 & 输入形状（每期） & 结构与层数 & 输出形状（每期）\\
 \midrule
 主干 & $256\times256\times3$ & ResNet-18前4个阶段，第4阶段步长改为1；两期共享权重 & $32\times32\times256$\\
 压缩与上采样 & $32\times32\times256$ & 逐点卷积到32通道，双线性放大2倍 & $64\times64\times32$\\
 语义分词器 & $4096\times32$ & 逐点卷积得到4张空间注意力图，各自softmax后加权求和 & $4\times32$\\
 Transformer编码器 & 两期拼接$8\times32$ & 1层，8头 & $8\times32$，再拆回两期\\
 Transformer解码器 & 像素$4096\times32$与本期4个词元 & 8层，只有交叉注意力与前馈网络 & $64\times64\times32$\\
 预测头 & 两期特征放大到$256\times256\times32$后取绝对差 & 2个$3\times3$卷积（32、2通道），逐像素softmax & $256\times256\times2$\\
 \bottomrule
 \end{tabularx}
 \caption{BIT主模型（原文记为BIT\_S4）的分段结构。层数、宽度与参数量取自原文；张量形状按$256\times256$输入推算，ResNet-18第4阶段的256通道沿用其标准配置。}
 \label{tab:vision-recognition-bit-structure}
\end{table*}

训练使用配准影像与像素变化标签，在有效像素上以式\eqref{eq:vision-dense-pixel-training}的交叉熵学习“变化／不变”。每次训练对两期图像和变化标签施加同样的裁剪与翻转，保证同一位置仍指同一地表。主干从ImageNet权重初始化，分词器、Transformer和预测头由变化数据学习。推理时保持两期的顺序与配准，把分数恢复到原网格，再用阈值选出变化区域。

新建筑应计为变化，树木由绿变黄是否计为变化则取决于标注定义。若两期影像整体错开几个像素，差异分支也会在边缘产生响应，因此配准是任务前提。词元交换能够利用上下文关联，却不能把传感器、季节或光照差异自动解释成真实的建设活动。

两个模型的对比如表\ref{tab:vision-recognition-pixel-compare}。BIT相对DeepLabv3+改变了输入（一幅变两幅）和上下文的组织方式（空洞卷积的固定邻域变为少量全局词元），换来了两期之间的全局比较和较低的注意力成本；代价是只有$k$个词元时细小变化可能被汇总掉，且整个方法依赖配准。

\begin{table}[htbp]
 \centering
 \small
 \begin{tabularx}{\linewidth}{@{}lXX@{}}
 \toprule
 比较项 & DeepLabv3+ & BIT\\
 \midrule
 输入 & 单幅图像 & 同一地区两期配准影像\\
 类别含义 & 语义类别 & 变化／不变\\
 上下文来源 & 多空洞率卷积与整图池化 & 两期共享的少量语义词元\\
 边界恢复 & 解码器融合步幅4的浅层特征 & 词元投回像素后上采样\\
 主要前提与风险 & 细结构须在输入中保留 & 配准误差与非目标变化\\
 \bottomrule
 \end{tabularx}
 \caption{两种像素类别预测模型的比较。两者输出都是逐像素类别图，差别在类别含义与上下文的组织。}
 \label{tab:vision-recognition-pixel-compare}
\end{table}

\subsection{对象掩码的预测}
\label{sec:vision-recognition-mask-sets}

对象掩码的输出是一组二值图，每张图对应一个实例，并附类别或得分。它比像素类别图多出的信息是身份：两辆相邻的车必须分在两张掩码里。难点因此有两层：一是确定有哪些对象，即检测问题中“哪个输出负责哪个目标”在像素层面的重现；二是确定要哪些对象，即目标由类别表、用户点击还是概念文字指定。

四个模型对这两层问题给出了不同的组合。Mask R-CNN先检测框，再在每个框内预测掩码，把实例分开的任务交给检测器，掩码分支只需在框内区分前景和背景。Mask2Former取消了框，沿用DETR的集合匹配思路，让每个查询直接预测一张整图掩码，同一结构可以输出语义、实例和全景结果。SAM改变了目标指定方式：不再依赖类别表，而是由用户的点或框指定一个对象，输出不带类别名的掩码。SAM 3再把指定方式从一个局部点击扩展为一个概念短语，要求找出图中所有符合概念的实例。

\subsubsection{Mask R-CNN}
\label{sec:vision-recognition-mask-rcnn}

两阶段检测器Faster R-CNN已经能为每个对象给出框和类别，自然的扩展是在框内再预测一张掩码。但它从特征图上截取区域特征时，会把框的连续坐标两次取整到特征格点；这对框分类影响不大，对像素级掩码却会造成可见的错位。Mask R-CNN的想法是在检测分支旁并行增加一个小型掩码分支，并用区域对齐代替取整，使框内位置与特征准确对应\scite{vision-1703-06870}。

网络的第一阶段是\term{区域建议网络}（region proposal network, RPN）：在FPN各层的每个位置放置锚框，预测对象性得分和框修正，解码并去重后得到约一千个候选框。第二阶段对每个候选用\term{区域对齐}（RoIAlign）取特征：把候选框等分为固定数目的小格，在每格中按连续坐标取若干采样点，用双线性插值读取特征值，再平均或取最大，得到固定尺寸的区域特征。设四个相邻格点的特征值为$f_{0,0},f_{0,1},f_{1,0},f_{1,1}$，采样点在水平方向偏移$a$、竖直方向偏移$b$，二者均在$[0,1]$内，则
\begin{equation}
 f(a,b)=(1-a)(1-b)f_{0,0}
 +a(1-b)f_{0,1}+(1-a)bf_{1,0}+abf_{1,1}\text{。}
 \label{eq:vision-roialign}
\end{equation}
用数值$(0,4,2,6)$与$(a,b)=(0.25,0.5)$可得$f=2$，既保留横向的四分之一偏移，也保留纵向中点；若直接取最近格点，这些偏移信息就会消失。原文在每个小格中取4个规则分布的采样点，再取平均，并指出结果对采样点的具体位置和个数不敏感。以ResNet-FPN为主干时，两个头的结构见表\ref{tab:vision-recognition-maskrcnn-heads}：框分支读取$7\times7\times256$的区域特征，经两个1024维全连接层输出类别和框修正；掩码分支读取$14\times14\times256$的区域特征，经4个$3\times3$卷积和一个$2\times2$、步长2的反卷积，输出$28\times28\times K$，每个类别一个二值掩码通道。原文另有以ResNet第4阶段特征为主干的C4版本：区域特征为$7\times7\times1024$，先经第5阶段残差块得到$7\times7\times2048$，框分支在其平均池化后输出，掩码分支经一次反卷积得到$14\times14\times K$，分辨率比FPN版低一半。

\begin{table*}[tbp]
 \centering
 \small
 \begin{tabular}{@{}lllrrl@{}}
 \toprule
 阶段或层组 & 层类型 & 卷积核与步长 & 输出通道 & 重复次数 & 输出尺寸\\
 \midrule
 框分支 & 区域对齐 & — & 256 & 1 & $7\times7$\\
        & 全连接 & — & 1024 & 2 & $1\times1$\\
        & 全连接 & — & $K$；$4K$ & 1 & $1\times1$\\
 掩码分支 & 区域对齐 & — & 256 & 1 & $14\times14$\\
          & 卷积 & $3\times3$ & 256 & 4 & $14\times14$\\
          & 反卷积 & $2\times2$，步长2 & 256 & 1 & $28\times28$\\
          & 卷积 & $1\times1$ & $K$ & 1 & $28\times28$\\
 \bottomrule
 \end{tabular}
 \caption{Mask R-CNN（ResNet-FPN）的两个区域头，取自原文图4。隐藏层均接ReLU；COCO上$K=80$。框分支最后一层并列输出类别得分与每类的框修正。}
 \label{tab:vision-recognition-maskrcnn-heads}
\end{table*}

训练时，与某个真实框IoU不低于0.5的候选为正例，并记住匹配的实例编号。总目标为$L=L_{\mathrm{cls}}+L_{\mathrm{box}}+L_{\mathrm{mask}}$，再加上RPN的对象性与框损失。掩码损失只在正例上、只对其匹配真实类别的那个通道计算逐像素sigmoid二元交叉熵；目标是匹配实例的掩码裁剪到候选框并缩放为$28\times28$。这样各类别的掩码通道不在同一像素上竞争，类别判断交给框分支。一个正候选若匹配错了实例，即使“车辆”类别正确，也会学到另一辆车的掩码。主干从分类权重初始化，各部分联合训练。原文输入短边为800像素，FPN版每幅图采样512个区域，正负比例1比3；8块GPU、每块2幅图，共16万次迭代。

推理时FPN版保留1000个候选，经框分支筛选、去重后，只对得分最高的100个框运行掩码分支，取预测类别对应的掩码通道，以0.5为阈值二值化，并缩放放回原图。若一个$28\times28$掩码对应原框$(20,10,80,50)$，归一化区域位置$(\xi,\eta)$回到$(20+60\xi,10+40\eta)$。恢复后的掩码、类别和框共享实例编号。区域对齐解决采样量化误差，不能补回候选阶段漏掉的车辆；$28\times28$的掩码分辨率对大对象的细边界也有限。

\subsubsection{Mask2Former}
\label{sec:vision-recognition-mask2former}

Mask R-CNN依赖框：框漏了，掩码就没有；道路、天空这类不可数区域又不适合用框表示，语义分割和全景分割只好另用一套网络。Mask2Former的想法是沿用DETR的集合预测，但每个查询输出的不是框而是一张整图掩码加一个类别，这样语义、实例和全景三种任务都可以写成“预测一组带类别的掩码”\scite{vision-2112-01527}。它相对早期掩码分类方法的关键改动是掩码注意力：每个查询只读取自己当前预测的前景区域。

网络由三部分组成。主干产生多尺度特征；像素解码器（默认为6层多尺度可变形注意力）在$1/32$、$1/16$、$1/8$三个尺度上交换信息，再经一次上采样和横向连接得到$1/4$分辨率的逐像素嵌入$\bm{f}_p\in\mathbb{R}^{d}$。Transformer解码器维护$N$个查询（默认100），共9层，按$1/32$、$1/16$、$1/8$的顺序循环读取三个尺度的特征。每层先做掩码交叉注意力，再做查询间自注意力和前馈网络。设第$i$个查询解码成掩码向量$\bm{e}_i\in\mathbb{R}^d$，它在位置$p$的掩码概率为$\sigma(\bm{e}_i^\top\bm{f}_p)$。\term{掩码注意力}（masked attention）把上一层预测的掩码缩放到当前尺度并以0.5二值化，在注意力得分中加入屏蔽项$M_{i,j}$：前景位置为0，其余为$-\infty$，再作softmax。这样每个查询只聚合自己可能对应区域内的特征，而不是被整图背景稀释。若某个查询的掩码为空，所有位置都被屏蔽，分母失效，实现需恢复可读位置。与早期掩码分类方法相比，Mask2Former还把查询特征改为可学习的初始值，并在每层给特征加正弦位置编码和可学习的尺度编码，去掉了解码器中的随机失活。

\begin{table}[htbp]
 \centering
 \small
 \begin{tabular}{@{}lcccc@{}}
 \toprule
 主干 & 查询数 & 解码层数 & 参数量 & 计算量\\
 \midrule
 ResNet-50 & 100 & 9 & 44M & 226G\\
 ResNet-101 & 100 & 9 & 63M & 293G\\
 Swin-L & 200 & 9 & 216M & 868G\\
 \bottomrule
 \end{tabular}
 \caption{Mask2Former在COCO全景分割上的配置，取自原文。像素解码器均为6层多尺度可变形注意力；计算量为每幅图的浮点运算次数。Transformer的宽度、头数与前馈网络宽度在正文中只以符号出现，论文未给出具体数值。}
 \label{tab:vision-recognition-mask2former-config}
\end{table}

表\ref{tab:vision-recognition-mask2former-config}列出原文报告的三个规模。参数增长主要来自主干；解码部分在三种配置中结构相同，只有Swin-L的全景与实例模型把查询数增加到200。

训练沿用DETR式二分匹配，匹配代价和损失都由类别项和掩码项组成。Dice项度量软预测与真实前景的重叠，对于预测$p_j$和真实标签$m_j$，一种带平滑量$\epsilon>0$的形式为
\begin{equation}
 L_{\mathrm{Dice}}=
 1-\frac{2\sum_j p_jm_j+\epsilon}
 {\sum_j p_j+\sum_j m_j+\epsilon}\text{。}
 \label{eq:vision-dice}
\end{equation}
二元交叉熵逐像素约束预测，Dice从区域整体约束重叠，二者对小对象和大对象的敏感程度不同。原文的类别、二元交叉熵和Dice权重分别为2、5、5，“无对象”类权重为0.1。掩码损失不在全图计算，而是在$K=12544$个（即$112\times112$个）采样点上计算，匹配时均匀采样，最终损失时按重要性采样偏向预测不确定的点，以降低高分辨率掩码的显存开销；原文报告这使每幅图的训练显存从18GB降到6GB。每层解码输出都计算辅助损失。COCO上的训练使用AdamW，学习率$10^{-4}$（主干乘0.1），批大小16，训练50轮，输入经大尺度抖动后裁剪为$1024\times1024$。

推理时按任务组织输出。语义输出将各查询的类别概率与掩码概率相乘，对查询求和，再逐像素选类；实例输出保留查询编号，按类别得分与掩码质量筛选；全景输出还要处理区域重叠、无对象和不可数类别。两个车辆查询在语义图中可以合并，在实例图中不能合并。原工作按任务与数据集分别训练，统一架构不等于一份参数对所有任务通用；早期掩码错误也可能限制后续读取，小对象和被遮挡对象需要单独检查。

\subsubsection{SAM}
\label{sec:vision-recognition-sam}

前两个模型都只能分割训练类别表中的对象。用户有时只想选中图上的一个区域，甚至说不出它的类别名，例如一块污渍或车门上的贴纸。SAM（segment anything model）定义了\term{可提示分割}（promptable segmentation）任务：给定图像和一个提示（点、框或粗略掩码），输出提示所指对象的有效掩码；提示有歧义时，至少给出一个合理的结果\scite{vision-2304-02643}。它的想法是把昂贵的图像编码与廉价的提示解码分开，同一幅图编码一次后，用户每改一次提示只重跑轻量解码器，从而支持实时交互。

三个组件的配置见表\ref{tab:vision-recognition-sam-config}，其中ViT-H的层数与宽度沿用ViT原文的Huge配置\scite{dosovitskiy2020vit}。图像编码器是以掩码自编码预训练的ViT-H，图像块边长16，输入缩放并填充到$1024\times1024$，因此得到$64\times64=4096$个词元。对4096个词元做全局注意力代价很高，SAM的编码器大多数块只在$14\times14$的窗口内做注意力，另有4个均匀分布的块做全局注意力，让远处信息仍能交换。编码器之后接一个颈部：$1\times1$卷积把通道降到256，再接一个256通道的$3\times3$卷积，两者之后各有层归一化，输出$64\times64\times256$的图像嵌入。

提示编码器把点的坐标经位置编码映射为256维向量，并加上“正点”“负点”或“框左上角”“框右下角”类型的学习嵌入；粗略掩码以图像嵌入4倍的分辨率输入，经两个$2\times2$、步长2的卷积（输出4和16通道）和一个$1\times1$卷积变成256通道，与图像嵌入逐位置相加，没有掩码提示时加一个学习的“无掩码”嵌入。掩码解码器是两层改造的Transformer，宽度256、8个头、前馈网络中间宽度2048，交叉注意力中查询、键、值降到128维以节省计算。输出词元（3个掩码词元加1个IoU词元）与提示词元拼在一起，每层依次做：词元自注意力、词元到图像的交叉注意力、逐词元前馈网络、图像到词元的交叉注意力。最后一次的“图像读词元”让图像嵌入也吸收了提示信息。解码后，图像嵌入经两个$2\times2$、步长2的反卷积（输出64和32通道）上采样4倍到$256\times256$，每个掩码词元经3层MLP得到32维的动态权重，与上采样特征逐位置做内积，得到一张$256\times256$的掩码logit。

\begin{table*}[tbp]
 \centering
 \small
 \begin{tabular}{@{}lccccccl@{}}
 \toprule
 组件 & 图像块大小 & 词元数 & 层数 & 宽度 & 头数 & MLP倍率 & 参数量\\
 \midrule
 图像编码器（ViT-H） & 16 & 4096 & 32 & 1280 & 16 & 4 & 约6.3亿\\
 提示编码器 & — & 每点或角点1个 & — & 256 & — & — & 论文未单列\\
 掩码解码器 & — & 输出与提示词元 & 2 & 256 & 8 & 8 & 论文未单列\\
 \bottomrule
 \end{tabular}
 \caption{SAM的配置。ViT-H的层数、宽度、头数与MLP倍率沿用ViT原文的Huge配置，SAM原文图13标注其编码器约636M参数；SAM编码器中4个块为全局注意力、其余为$14\times14$窗口注意力。掩码解码器的MLP倍率按前馈宽度2048与宽度256折算。}
 \label{tab:vision-recognition-sam-config}
\end{table*}

一个提示可能有多种合理解释：落在车门上的点可以指车门，也可以指整辆车。SAM为单个提示同时输出三张掩码（大致对应局部、部件、整体），另有一个头为每张掩码预测它与目标的IoU，用于排序。训练时只对三张中与真实掩码损失最小的那一张反传掩码损失，这让不同输出分工到不同粒度，而不是被平均成模糊的中间结果；掩码损失是焦点损失与Dice损失按20比1线性组合，IoU头用均方误差监督。训练模拟交互：先从真实掩码中采一个点或框，之后在当前预测的错误区域采新点，并把上一轮的掩码logit作为提示输入，每个掩码共模拟11轮。训练数据来自作者的数据引擎，模型辅助人工标注，再逐步转为全自动标注，最终形成约1100万张图像、11亿个掩码的SA-1B数据集。最终模型用AdamW训练9万次迭代，批大小256，在256块A100上约需68小时。

使用时先缓存图像嵌入，点或框改变只重跑提示编码和解码；掩码logit上采样到输入尺寸，去掉填充并恢复原图尺寸，再以0为阈值二值化。用户在漏掉的车尾补正点、在误选的路面补负点，可以逐步消除歧义；IoU预测可以推荐候选，但选择车门还是整车仍由用户意图决定。报告结果时应说明提示来源和点击次数：用真实标注框作提示与由检测器自动给框，具有不同的信息和错误条件。SAM的掩码不附带类别名称。

\subsubsection{SAM 3}
\label{sec:vision-recognition-sam3}

SAM的一个点只选择一个对象。“找出所有校车”则需要先判断这个概念是否出现，再枚举所有符合概念的实例；用SAM完成这件事，必须由用户逐个点击，或者先用开放词汇检测器给框。SAM 3把指定方式从局部提示扩展为概念：输入一个简短名词短语和若干带正负标签的图像示例框，输出所有匹配实例的框和掩码。原报告称这一任务为可提示概念分割（promptable concept segmentation, PCS）\scite{vision-2511-16719}。

模型的核心想法是把“概念是否出现”与“哪个实例匹配”分开。普通DETR式检测器让每个查询同时回答这两个问题：一辆普通客车的查询既要判断自己是车，又要判断全图是否存在校车，局部信息不足以回答后者。SAM 3增加一个共享的存在词元，与对象查询一同进入解码器，由存在预测头得到$p_{\mathrm{pres}}$；查询$i$的局部分数$s_{i\mid\mathrm{pres}}$只表达“若概念存在，我是否为其一例”，最终分数为$s_i=p_{\mathrm{pres}}s_{i\mid\mathrm{pres}}$。

检测器的结构沿用DETR式查询解码。文字短语经文字编码器得到词元特征，示例框经编码后同时携带位置、框内视觉内容和正负标签，二者共同组成概念条件。图像特征来自与文字编码器联合预训练的视觉编码器，经融合编码器读取概念条件；对象查询在解码器中细化框，掩码头再为每个查询产生实例掩码。负示例告诉模型哪些相似对象不属于当前概念，而不是把该区域的像素删除。

原报告附录给出了各组件的规模，见表\ref{tab:vision-recognition-sam3-structure}。视觉编码器是在54亿图文对上对比训练的感知编码器（Perception Encoder），共32层，其中4层做全局注意力，其余层在窗口内做注意力：$1008\times1008$的输入被分成$3\times3$个336像素见方的窗口，每个窗口$24\times24$个词元。融合编码器的每个块依次做图像自注意力、对提示词元的交叉注意力和前馈网络；检测解码器沿用DETR式结构，并加入以框为参照的位置偏置、逐层修正框等改进。存在词元与对象查询一同进入解码器，输出经MLP头给出所有查询共享的$p_{\mathrm{pres}}$。

\begin{table*}[tbp]
 \centering
 \small
 \begin{tabularx}{\linewidth}{@{}llXl@{}}
 \toprule
 部分 & 输入 & 结构与层数 & 规模或输出\\
 \midrule
 视觉编码器 & $1008\times1008$图像 & 感知编码器，32层（4层全局注意力，其余窗口注意力） & 约4.5亿参数\\
 文字编码器 & 名词短语 & 因果Transformer，上下文至多32个词元 & 约3亿参数\\
 融合编码器 & 图像特征与提示词元 & 6个块：自注意力、对提示的交叉注意力、前馈网络 & 条件化图像特征\\
 检测解码器 & 200个对象查询加存在词元 & 6个块，DETR式 & 每查询的框、局部分数；共享的存在分数\\
 分割头 & 解码后的查询与多尺度特征 & 改编自MaskFormer，语义与实例分割共用 & 每查询一张实例掩码\\
 跟踪器 & 视频帧 & 沿用SAM~2的记忆注意力与记忆库，与检测器共用视觉编码器 & 跨帧掩码\\
 \bottomrule
 \end{tabularx}
 \caption{SAM~3的组件，取自原报告附录。检测器与跟踪器合计约1亿参数，全模型约8.5亿参数；图像块边长原文未直接写出。}
 \label{tab:vision-recognition-sam3-structure}
\end{table*}

训练分阶段进行。先用大规模图文对预训练视觉和文字编码器；再以图像、短语和掩码组成的数据训练检测器及两侧编码器，数据既有伪标注也有人工核验的标注；高质量阶段加入存在头，并根据漏检采正示例、根据高分误检采负示例，累加示例以模拟交互纠错。存在头用二元交叉熵覆盖含概念的正图和难负概念；概念存在时，实例按类别和范围代价做二分匹配，再以局部得分、框的$L_1$与广义IoU、掩码的焦点与Dice损失训练查询。概念缺席的图像不向局部分类分支施加逐查询的负梯度，这与普通DETR的无对象训练不同。视频跟踪阶段冻结主干训练跟踪器，不应倒推为前面的图像训练也一直冻结。使用时更换“校车”文字只改变条件，不重新学习权重。

若存在分数为0.9，两个实例的局部分数为0.8和0.6，最终分数为0.72和0.54；若存在分数为0.05，同样的局部分数只剩0.04和0.03。保留超过阈值的查询，将各掩码放回原图并保持独立实例编号。共享存在分支的一个错误会同时改变所有查询：把“校车”误判为不存在，多个真实实例会一起被压低，形成相关的漏检。评价因此分别检查概念有无、每实例召回和掩码边界。简短概念和示例框扩大了目标的指定方式，但不足以推断任意关系嵌套或长指令都已解析；视频中的身份传播在\ref{chap:vision-video-recognition}章继续讨论。

四个模型的对比见表\ref{tab:vision-recognition-mask-compare}。从Mask R-CNN到Mask2Former，改变的是实例的来源（先有框，到查询直接给掩码），换来了三类分割的统一与不依赖框，代价是训练和显存更重；从Mask2Former到SAM，改变的是目标的指定方式（类别表到用户提示），换来了类别无关的交互分割，代价是失去类别名；SAM 3再把提示从一个点扩展到一个概念，恢复了语义，代价是存在判断错误会牵连所有实例。

\begin{table*}[tbp]
 \centering
 \small
 \begin{tabularx}{\linewidth}{@{}lXXXX@{}}
 \toprule
 模型 & 目标的指定 & 实例的来源 & 输出是否带类别 & 主要代价或风险\\
 \midrule
 Mask R-CNN & 固定类别表 & 先检测框，框内预测掩码 & 是 & 框漏检即掩码缺失，掩码分辨率有限\\
 Mask2Former & 固定类别表 & 查询经集合匹配直接给整图掩码 & 是 & 训练成本高，早期掩码错误限制读取\\
 SAM & 点、框或粗掩码 & 每个提示一组歧义候选 & 否 & 依赖提示质量，需交互消歧\\
 SAM 3 & 概念短语与正负示例 & 查询枚举全部匹配实例 & 是（概念） & 存在判断错误牵连全部实例\\
 \bottomrule
 \end{tabularx}
 \caption{同一街景掩码任务下四个模型的比较。表中只比较接口与机制，不给出不同论文设置下的精度。}
 \label{tab:vision-recognition-mask-compare}
\end{table*}

\subsection{前景透明度的预测}
\label{sec:vision-recognition-matting}

发丝和半透明边缘的一个像素可能同时收到前景与背景的光，二值掩码只能选择保留或删除，无法表达这种混合。\term{图像抠图}（image matting）估计每个像素的前景贡献比例$\alpha\in[0,1]$。在线性颜色空间的简化合成模型中，每个像素满足
\begin{equation}
 \bm{x}=\alpha\bm{f}+(1-\alpha)\bm{b}\text{，}
 \label{eq:vision-matting}
\end{equation}
其中$\bm{x},\bm{f},\bm{b}\in\mathbb{R}^3$是观测、前景与背景颜色。每个像素只有三个颜色观测，却有前景、背景和$\alpha$共七个未知数，问题本身欠定，需要额外条件或先验。路线的差别就在于额外条件从哪里来：传统方法假设颜色在局部平滑变化，由用户画出的\term{三分图}（trimap）给出确定前景、确定背景和未知区域三类标记，再在未知区域内传播颜色；学习方法仍可使用三分图作为条件，但把“怎样从图像和条件推出$\alpha$”交给网络，从大量合成样本中学习。不使用三分图的方法则让网络自行判断前景，代价是对目标指定的歧义更敏感。下面以Deep Image Matting说明有三分图条件的学习路线。

此前的抠图方法主要依赖局部颜色统计：当前景和背景颜色相近，或背景纹理复杂时，颜色本身不足以区分，结果在发丝处出现明显误差；另一个限制是缺少足够大的真实$\alpha$数据集来训练学习方法。Deep Image Matting从两方面入手：构造大规模合成训练集，并用深度网络从图像的结构和语义信息中推断$\alpha$，而不只依赖颜色\scite{vision-1703-03872}。

网络分两个阶段。第一阶段是以VGG-16为基础的编码器—解码器，输入是RGB图像与三分图拼接成的$H\times W\times4$张量，编码器逐级下采样提取特征，解码器逐级上采样输出$H\times W\times1$的初步$\hat\alpha$。第二阶段是一个小型全卷积细化网络，输入是原图与初步$\hat\alpha$，预测残差以锐化边缘细节，输出最终$\alpha$。确定前景和确定背景区域直接取1和0，学习的重点是未知带中的混合比例。

逐层结构见表\ref{tab:vision-recognition-matting-layers}。编码器有14个卷积层和5个最大池化层，取VGG-16的前14个卷积层（第14层由全连接层fc6改写为卷积）初始化，第四个输入通道的权重初始化为零。解码器比编码器小得多：6个卷积层和5个反池化层，最后接一个输出单通道的预测层；每个反池化层把特征放大回上一级的分辨率，与编码器的5次池化对应。细化网络是4个$3\times3$卷积，前3层64通道并接ReLU，最后一层输出1通道；它不做下采样，输入的第四通道（初步$\hat\alpha$，缩放到$[0,1]$）直接加到输出上，网络只学习残差。

\begin{table*}[tbp]
 \centering
 \small
 \begin{tabular}{@{}lllrrl@{}}
 \toprule
 阶段或层组 & 层类型 & 卷积核与步长 & 输出通道 & 重复次数 & 输出尺寸\\
 \midrule
 编码器 & 卷积与池化 & VGG-16前14层 & 至512 & 14＋5 & $H/32\times W/32$\\
 解码器 & 卷积 & $1\times1$ & 512 & 1 & $H/32\times W/32$\\
        & 反池化与卷积 & $5\times5$ & 512，256，128，64，64 & 5 & 逐级至$H\times W$\\
        & 预测层 & $5\times5$ & 1 & 1 & $H\times W$\\
 细化网络 & 卷积 & $3\times3$ & 64 & 3 & $H\times W$\\
          & 卷积 & $3\times3$ & 1 & 1 & $H\times W$\\
 \bottomrule
 \end{tabular}
 \caption{Deep Image Matting的逐层结构。编码器与解码器的层数取自原文正文，解码器各层的核大小与通道数取自原文图3；输入为$H\times W\times4$（RGB加三分图），训练时$H=W=320$。细化网络的输入为原图与初步$\hat\alpha$，共4通道。}
 \label{tab:vision-recognition-matting-layers}
\end{table*}

训练数据由合成得到：把已有精确$\alpha$和前景颜色的对象，按式\eqref{eq:vision-matting}合成到大量不同的背景图上；原文用493个前景对象各配100张来自COCO和VOC的背景，得到49300张训练图，测试集则由50个对象各配20张背景组成；训练时以未知像素为中心随机裁剪$320\times320$的块；三分图由真实$\alpha$的未知区域随机膨胀生成，使网络适应不同宽度的未知带。第一阶段的目标在未知区$U$上组合两项：
\[
 \begin{aligned}
 L={}&w_\ell\frac1{|U|}\sum_{p\in U}\sqrt{(\hat\alpha_p-\alpha_p)^2+\epsilon^2}\\
 &+(1-w_\ell)\frac1{|U|}\sum_{p\in U}
 \sqrt{\lVert\hat\alpha_p\bm f_p+(1-\hat\alpha_p)\bm b_p-\bm x_p\rVert_2^2+\epsilon^2}\text{，}
 \end{aligned}
\]
其中$\epsilon$是避免零点不可导的小常数，原文取$w_\ell=0.5$。第一项直接约束贡献比例，第二项是合成损失，用已知的前景和背景重新合成，检查该比例的颜色后果；两项通过同一混合公式连接。编码器从分类网络权重初始化，先训练第一阶段，再固定它训练细化网络，最后联合微调。三分图可以由目标$\alpha$构造以供训练，但使用时必须实际取得提示，不能把真实$\alpha$隐藏在输入中。

使用时输出$\alpha$；要把对象放到新背景$\bm{b}'$上，还需估计或取得前景颜色$\bm{f}$，再合成$\alpha\bm{f}+(1-\alpha)\bm{b}'$。例如灰度前景为0.8、原背景为0.2、$\alpha=0.25$，观测为0.35。换成0.6的背景，正确合成为$0.25\times0.8+0.75\times0.6=0.65$；若二值化为背景，结果为0.6，发丝的贡献消失；若直接把观测0.35当成前景颜色，又会混入旧背景。三分图若误把发丝标成确定背景，网络会被错误条件限制。合成训练数据与真实照片在光照和噪声上存在差距，边界、贡献比例和最终合成质量应分别核验。

\section{数量、姿态与空间测量}
\label{sec:vision-recognition-measurement}

“有几辆车”“人的手腕在哪里”“路牌离相机多远”“这几张照片是从哪里拍的”分别要求数量、带部位语义的位置、深度和相机位姿。它们都可以借用上一节的对象范围，但输出已经是测量量：有单位、有坐标系，误差也要按测量量计算，不能由检测或分割指标代替。本节按测量量划分为四个小节。数量只需要规定个体和计入规则；姿态要给每个关节一个位置，并说明它属于哪个人；深度要给每个像素一个距离；相机与场景几何则要同时确定拍摄位置和三维结构。从数量到几何，所需的几何约定越来越多，下面先回顾后三个小节共用的相机模型。

要把手腕或路牌放到三维空间中，需要\ref{sec:vision-formation-camera}节建立的相机模型。这里只回顾后文要用的结论和记号，推导与标定方法见\ref{sec:vision-formation-camera-intrinsics}节和\ref{sec:vision-formation-calibration}节。本章忽略镜头畸变（实际使用前先用标定参数校正），像素坐标向右、向下增加，相机光轴向前。

世界点$\bm{P}_{\mathrm w}\in\mathbb{R}^3$经外参$\bm{R},\bm{t}$变到相机坐标$\bm{P}_{\mathrm c}=(x_{\mathrm c},y_{\mathrm c},z)^\top=\bm{R}\bm{P}_{\mathrm w}+\bm{t}$，其中$\bm{R}$是旋转矩阵；$\bm{t}$是世界到相机变换的平移，相机中心的世界坐标则是$-\bm{R}^{\top}\bm{t}$。内参矩阵$\bm{K}$含像素焦距$f_x,f_y$与主点$(c_x,c_y)$，于是像素$(u,v)$满足
\begin{equation}
 z\begin{pmatrix}u\\v\\1\end{pmatrix}
 =\bm{K}(\bm{R}\bm{P}_{\mathrm w}+\bm{t}),\qquad
 \bm{K}=\begin{pmatrix}f_x&0&c_x\\0&f_y&c_y\\0&0&1\end{pmatrix}\text{。}
 \label{eq:vision-camera-projection}
\end{equation}
$z>0$表示点在相机前方。本章反复用到这个式子的逆向读法：令$\tilde{\bm{p}}=(u,v,1)^\top$，沿这条像素射线的相机点都是$\bm{P}_{\mathrm c}(z)=z\bm{K}^{-1}\tilde{\bm{p}}$，$z>0$待定，所以像素不能单独确定深度。光轴深度$z$与到相机中心的直线距离$\rho$也不是同一个量：
\begin{equation}
 \rho=\lVert\bm{P}_{\mathrm c}\rVert_2
 =z\sqrt{1+\left(\frac{u-c_x}{f_x}\right)^2
            +\left(\frac{v-c_y}{f_y}\right)^2}\text{。}
 \label{eq:vision-depth-range-measurement}
\end{equation}
教学中取$f_x=f_y=800$像素、主点$(320,240)$和像素$(520,240)$。若$z=4$米，反投影得到$(1,0,4)$米，直线距离为$\sqrt{17}\approx4.123$米；若$z=8$米，点$(2,0,8)$仍投到同一像素。两次投影相同，测量量却不同。

多视图增加了射线约束，但还需区分\emph{相对尺度}（relative scale）与\emph{度量尺度}（metric scale）：前者只确定长度比例，后者还确定一单位对应多少米。把全部世界点和相机平移同时乘$a>0$，相机点与$z$也乘$a$，式\eqref{eq:vision-camera-projection}的像素不变。已知相机基线、已知物体长度或带单位的深度观测可以确定这个比例；只把首相机设为原点，固定的是参照位置和朝向，没有测出尺度。网络从训练分布中学到的尺度是附加先验，须在当前设备与场景中验证。

同一个世界点进入第二幅图时，消去世界点，可得\ref{sec:vision-formation-two-view}节的两相机相对位姿：
\begin{equation}
 \begin{aligned}
 \bm{P}_{\mathrm c,2}
 &=\bm{R}_{2,1}\bm{P}_{\mathrm c,1}+\bm{t}_{2,1},\\
 \bm{R}_{2,1}&=\bm{R}_2\bm{R}_1^\top,\qquad
 \bm{t}_{2,1}=\bm{t}_2-\bm{R}_2\bm{R}_1^\top\bm{t}_1\text{。}
 \end{aligned}
 \label{eq:vision-relative-camera-measurement}
\end{equation}
将第一幅图的反投影点代入，再用第二幅图的$\bm{K}_2$投影，就得到一个可检查的跨图位置。裁剪和缩放也必须进入内参：若横纵缩放为$a_x,a_y$，裁剪原点为$(u_0,v_0)$，则新焦距为$(a_xf_x,a_yf_y)$，新主点为$(a_x(c_x-u_0),a_y(c_y-v_0))$。把新像素配旧内参会直接改变射线方向。落在图像内且位于相机前方只是可见的必要条件；\emph{遮挡}（occlusion）指同一射线上更近的表面挡住更远的表面，跨图比较一个投影点时还要检查深度是否相容。以下姿态、深度和几何恢复都沿用这些约定。

\subsection{对象数量的估计}
\label{sec:vision-recognition-counting}

计数的输入是图像和一个计数规则（数什么类别、在哪个区域、被遮挡一半的算不算），输出是一个数，通常还希望附带每个对象的位置。难点随密度变化：路口两三辆车时，逐个找出来再数最可靠；广场上几百个人头相互遮挡时，单个人的范围难以标注，检测器也难以分开。

两条路线的差别在于是否先形成个体。基于检测实例汇总的计数先用上一节的检测器得到对象清单，再按规则数清单中的项；它保留每个对象的位置和身份，可以逐个核验，但在密集遮挡中会漏检和重复。基于密度图积分的计数不形成个体，而是预测每个位置分摊多少计数，再对整幅图求和；它适合密集人群，代价是失去个体清单。

\subsubsection{基于检测实例汇总的计数}
\label{sec:vision-recognition-instance-counting}

这条路线不需要新模型：检测器可以是\ref{sec:vision-recognition-detection}节的任何一个，计数完全由汇总规则决定。给定去重后的检测集合，按类别、得分和区域筛选，再把满足条件的指示量相加。若区域为$R$，以框中心$\bm{o}_i$是否在$R$内为计入规则，则
\begin{equation}
 \hat n_R=\sum_i
 \mathbf{1}\{c_i\in C_{\mathrm{target}}\}
 \mathbf{1}\{s_i\geq\eta\}
 \mathbf{1}\{\bm{o}_i\in R\}\text{，}
 \label{eq:vision-detection-count}
\end{equation}
其中$c_i$是类别，$s_i$是得分，$\eta$是得分阈值。“中心落入”与“任意可见部分落入”是两种不同规则，边界上的对象会得到不同计数，规则须事先规定。

阈值$\eta$直接决定计数的偏差：阈值高则漏数，阈值低则多数，它应在验证集上按计数误差而不是检测精度选择。拥挤场景中NMS可能删掉相邻的真实对象，DETR式检测器则受查询数上限约束。例如真实有5辆车，漏掉1辆，另1辆被重复预测，总数仍为5；只看计数误差会掩盖两种错误，因此还应检查对象匹配。跨帧统计时，同一辆车会在多帧中出现，需由\ref{chap:vision-video-recognition}章的跟踪维持身份。

\subsubsection{基于密度图积分的计数}
\label{sec:vision-recognition-mcnn}

人群过密时，标注每个人的框既昂贵又不准确，检测器在严重遮挡下也会大量漏检；同一画面中，近处的人头可能有几十像素，远处的只有几个像素。MCNN用头部点标注代替框，学习一张\term{密度图}（density map），即每个位置分摊多少计数贡献；对密度图求和就是人数\scite{vision-mcnn2016}。它针对尺度变化的想法是用几个感受范围不同的卷积分支并行处理同一幅图，各自适应不同大小的人头。

首先需要把点标注变成密度监督。若标注点为$\bm{p}_i$，用归一化核构造
\begin{equation}
 D(\bm{p})=\sum_{i=1}^{n}G_{\sigma_i}(\bm{p}-\bm{p}_i),
 \qquad
 \sum_{\bm{p}}G_{\sigma_i}(\bm{p}-\bm{p}_i)=1\text{。}
 \label{eq:vision-density-target}
\end{equation}
每个核的和为1，所以整幅密度图之和等于标注人数。原工作采用几何自适应核：若第$i$点最近$k$个邻点的平均距离为$\bar d_i$，取$\sigma_i=\beta\bar d_i$，原文取$\beta=0.3$。拥挤处邻点近，核就窄；稀疏处核较宽，以此用点间距近似透视造成的人头大小变化。孤立点或点数不足时需回退为固定核宽。

网络有三列结构相同的卷积网络，每列都是“卷积—池化—卷积—池化—卷积—卷积”共4个卷积层和2个$2\times2$最大池化，激活函数为ReLU；三列只在卷积核大小和滤波器个数上不同，逐层配置见表\ref{tab:vision-recognition-mcnn-layers}。大核列从$9\times9$起始，适合近处的大人头；小核列从$5\times5$起始，适合远处的小人头。原文有意让大核的列使用较少的滤波器，以控制参数量。两次池化使三列的输出都是输入的$1/4$分辨率，于是可以在通道维拼接成$8+10+12=30$个通道，再经一个$1\times1$卷积映射为单通道密度图$\hat D$。这个$1\times1$卷积相当于为三列学习融合权重，它取代了全连接层，因此网络可以接收任意尺寸的输入，原文也不把图像缩放到统一尺寸，以免缩放扭曲密度。训练目标是与目标密度的平方误差$L=(2n)^{-1}\sum_{i=1}^n\sum_p(\hat D_i(p)-D_i(p))^2$，标签密度也须下采样到$1/4$分辨率，并保持总和不变。原文先分别预训练每一列（直接把该列第4个卷积层的输出映射为密度图），再用这些权重初始化三列并联合微调；数据增强是从每幅图裁出9个、各为原图四分之一大小的块。推理时对$\hat D$求和得到人数，结果可以是非整数。

\begin{table*}[tbp]
 \centering
 \small
 \begin{tabular}{@{}lllrrl@{}}
 \toprule
 阶段或层组 & 层类型 & 卷积核与步长 & 输出通道 & 重复次数 & 输出尺寸\\
 \midrule
 大核列 & 卷积，池化 & $9\times9$；$2\times2$，步长2 & 16 & 1 & $H/2\times W/2$\\
        & 卷积，池化 & $7\times7$；$2\times2$，步长2 & 32 & 1 & $H/4\times W/4$\\
        & 卷积 & $7\times7$ & 16，8 & 2 & $H/4\times W/4$\\
 中核列 & 卷积，池化 & $7\times7$；$2\times2$，步长2 & 20 & 1 & $H/2\times W/2$\\
        & 卷积，池化 & $5\times5$；$2\times2$，步长2 & 40 & 1 & $H/4\times W/4$\\
        & 卷积 & $5\times5$ & 20，10 & 2 & $H/4\times W/4$\\
 小核列 & 卷积，池化 & $5\times5$；$2\times2$，步长2 & 24 & 1 & $H/2\times W/2$\\
        & 卷积，池化 & $3\times3$；$2\times2$，步长2 & 48 & 1 & $H/4\times W/4$\\
        & 卷积 & $3\times3$ & 24，12 & 2 & $H/4\times W/4$\\
 融合 & 拼接，卷积 & $1\times1$ & 1 & 1 & $H/4\times W/4$\\
 \bottomrule
 \end{tabular}
 \caption{MCNN的逐层结构，取自原文图3。输入为$H\times W$的人群图像；“16，8”表示两个连续卷积层的输出通道依次为16和8。卷积保持空间尺寸，下采样只来自两次最大池化。原文未给出参数总量。}
 \label{tab:vision-recognition-mcnn-layers}
\end{table*}

保持总和是密度标签的关键约束。将每个密度单元分成四个，四个新单元应分享原来的质量，不能每个都复制原值；连续情形下，缩放$(x',y')=a(x,y)$须取$D'(x',y')=D(x'/a,y'/a)/a^2$，变量代换才给出$\int D'=\int D$。裁剪训练时先按中心规则确定哪些点计入，再为这些点在裁剪域构造归一化核；直接把原密度图裁一块，会把边外对象的核尾也计入。三个点各贡献1，整图密度之和为3；若一个边缘点的核有一半落到图外，直接截断后该点只贡献约0.5。

总和正确仍可能把本属左半边的人数移到右半边，应同时检查局部计数和空间误差。新的拍摄高度和密度范围可能超出训练条件；密度图也没有恢复每个人的身份或完整位置。两条路线在同一路口与广场图像上的取舍可以概括为：检测汇总在稀疏场景中给出可核验的清单，在密集场景中系统性漏数；密度积分在密集场景中更稳定，却只能核验区域总量，不能列出个体。

\subsection{人体姿态的估计}
\label{sec:vision-recognition-pose}

\term{姿态估计}（pose estimation）的输入是含人的图像，输出是每个人的关节位置。人体关节不同于图像中任意可匹配的角点：左腕、右肘具有固定的部位含义，多人场景还必须知道这些部位属于谁。输出可以是图像平面上的二维关节、相机坐标中的三维关节，也可以是完整的人体表面网格。难点随输出维数增加：二维要解决多人归属，三维要面对节首的深度歧义，网格还要确定关节之间的身体形状。

三个模型对应三级输出。OpenPose在整图上预测二维关节，并用部位亲和场把关节组装成不同的人；Martinez基线把二维关节作为输入，用一个简单网络回归三维关节，将二维定位误差与三维提升误差分开；SAM 3D Body则从图像直接恢复参数化人体网格，并接受关节点或掩码提示来纠正结果。

\subsubsection{OpenPose}
\label{sec:vision-recognition-openpose}

多人姿态的常见做法是自顶向下：先检测每个人，再在每个人的框内定位关节。它的计算量随人数线性增长，人与人重叠时检测框和裁剪质量又会直接影响关节结果。OpenPose采用自下而上路线：在整图上一次性找出所有关节候选，再判断哪些候选属于同一个人\scite{vision-1611-08050}。组装所依据的线索是\term{部位亲和场}（part affinity field, PAF）：对每一种肢段（例如前臂），在真实肢段经过的带状区域内写入从一端关节指向另一端关节的单位向量，其他位置为零。两个候选关节之间若真是同一人的前臂，连线经过的像素上的PAF方向就应与连线方向一致。

网络先用VGG-19的前10层把图像变成特征$\bm{F}$（空间步幅为8）。之后是多阶段的双分支卷积网络：第一阶段从$\bm{F}$同时预测$J$张关节置信图$\bm{S}^{(1)}$和$C$个肢段的PAF $\bm{L}^{(1)}$（每个肢段两通道，共$2C$通道）；第$t$阶段把$\bm{F}$、$\bm{S}^{(t-1)}$、$\bm{L}^{(t-1)}$拼接作为输入，重新预测两组输出。后续阶段因此能利用前一阶段的关节和肢段结果修正彼此，例如由清晰的肘部和前臂方向推断模糊的手腕。
各阶段的层配置见表\ref{tab:vision-recognition-openpose-stages}。第一阶段每个分支是3个$3\times3$卷积加2个$1\times1$卷积；从第二阶段起，每个分支换成5个$7\times7$卷积加2个$1\times1$卷积，更大的卷积核扩大了每层的感受范围，使后续阶段能在步幅8的特征上看到整条肢段两端的证据。原文共用6个阶段，每个阶段的两组输出都接受监督。测试时把$1080\times1920$的视频帧缩放到$368\times654$，在GTX 1080上九个人的画面中卷积网络约需99.6毫秒，后续的组装只需约0.58毫秒。原文的期刊扩展版把每个$7\times7$卷积换成3个连续的$3\times3$卷积并拼接其输出，先跑若干只预测PAF的阶段，再跑预测置信图的阶段；这里按会议版讲解。

\begin{table*}[tbp]
 \centering
 \small
 \begin{tabular}{@{}lllrrl@{}}
 \toprule
 阶段或层组 & 层类型 & 卷积核与步长 & 输出通道 & 重复次数 & 输出尺寸\\
 \midrule
 特征提取 & VGG-19前10层 & $3\times3$，含池化 & 论文未给出 & 1 & $H/8\times W/8$\\
 阶段1（每分支） & 卷积 & $3\times3$ & 论文未给出 & 3 & $H/8\times W/8$\\
                 & 卷积 & $1\times1$ & 末层为$J$或$2C$ & 2 & $H/8\times W/8$\\
 阶段$t\geq2$（每分支） & 卷积 & $7\times7$ & 论文未给出 & 5 & $H/8\times W/8$\\
                 & 卷积 & $1\times1$ & 末层为$J$或$2C$ & 2 & $H/8\times W/8$\\
 \bottomrule
 \end{tabular}
 \caption{OpenPose（会议版）的结构，取自原文图3与正文。阶段$t\geq2$的输入是$\bm F$、$\bm S^{(t-1)}$与$\bm L^{(t-1)}$的拼接，共6个阶段。两个分支结构相同、参数独立，分别输出$J$张置信图和$C$个肢段的PAF；中间层通道数和$J,C$的具体取值在原文正文中未给出。}
 \label{tab:vision-recognition-openpose-stages}
\end{table*}

训练从关键点标注构造两类目标。关节置信图在每个标注位置放一个高斯峰，同类多人的峰取逐点最大值；PAF在肢段带内写入单位方向，多人肢段重叠处取平均。每个阶段都对两组输出计算加权平方误差，权重$W(\bm{p})$在未标注的人身上为零，以免把漏标的人当作负例；每个阶段都接受监督，缓解深层网络的梯度消失。

推理时先在置信图上做非极大值抑制得到关节候选。对候选肘部$\bm{a}$和腕部$\bm{b}$，令方向$\bm{u}=(\bm{b}-\bm{a})/\lVert\bm{b}-\bm{a}\rVert_2$，沿连线对PAF积分：
\begin{equation}
 s(\bm{a},\bm{b})=
 \int_0^1
 \bm{L}((1-\lambda)\bm{a}+\lambda\bm{b})^\top
 \bm{u}\,\mathrm{d}\lambda\text{。}
 \label{eq:vision-paf-score}
\end{equation}
实现中在有限个点上采样求平均，$\bm{a}=\bm{b}$的退化候选应排除。对每一种肢段，所有肘部候选与腕部候选构成一个二分图，以$s$为边权求最大权匹配，保证一个关节在该肢段中只连一次；再按人体骨架树把共享关节的肢段拼成完整的人。假设五个采样点的方向内积为$(1,1,0.8,1,0.7)$，平均为0.9；连到另一人的手腕时变为$(0.1,0,-0.2,0,0.1)$，平均为0。两个手腕都定位正确，也仍需这一关联步骤才能把手臂接对。投影重合、严重遮挡和交叉肢段会使证据混杂，因此关节定位与人体组装应分开评价。

\subsubsection{Martinez基线}
\label{sec:vision-recognition-lifting}

从单幅图像估计三维姿态的端到端网络，误差可能来自两处：二维上没看准关节，或者看准了却无法推出深度。端到端训练把这两种误差混在一起，难以判断瓶颈。Martinez等人的基线把问题拆开：假定二维关节已经给出，只用一个简单的全连接网络把二维坐标“提升”为三维坐标\scite{vision-1705-03098}。实验结果表明这种简单提升已能取得有竞争力的误差，说明相当一部分困难在于二维检测，而不在三维推断本身。

输入是按固定部位顺序排列的$j$个二维关节，展平为$2j$维向量$\bm{x}_i$，并用训练集统计量做零均值单位方差规范化；输出是相机坐标中相对根关节（骨盆）的三维关节，展平为$3j$维向量$\bm{y}_i$。网络先用线性层把输入映射到1024维，接两个残差块，每块包含两组“线性层、批归一化、ReLU、随机失活”，残差连接为$\bm{h}'=\bm{h}+F_{\bm{\theta}}(\bm{h})$，最后一层线性层输出$3j$维。全网共6个线性层，约400万至500万个可训练参数；每层权重的范数被约束不超过1，随机失活比例为0.5，原文报告残差连接使误差下降约10\%。逐层结构见表\ref{tab:vision-recognition-martinez-layers}。单个残差块内的计算顺序是“线性层$\to$批归一化$\to$ReLU$\to$随机失活$\to$线性层$\to$批归一化$\to$ReLU$\to$随机失活$\to$与块输入相加”。网络很小，原文在Titan Xp上一次前向约2毫秒。

\begin{table*}[tbp]
 \centering
 \small
 \begin{tabular}{@{}lllrrl@{}}
 \toprule
 阶段或层组 & 层类型 & 卷积核与步长 & 输出通道 & 重复次数 & 输出尺寸\\
 \midrule
 输入 & 二维关节 & — & $2j$ & — & 向量\\
 输入层 & 线性 & — & 1024 & 1 & 向量\\
 残差块 & 线性（见正文） & — & 1024 & 2块\ensuremath{\times}2组 & 向量\\
 输出层 & 线性 & — & $3j$ & 1 & 向量\\
 \bottomrule
 \end{tabular}
 \caption{Martinez基线的结构，取自原文图1与正文。网络只有全连接层，“卷积核与步长”一列不适用；每个残差块含两组“线性、批归一化、ReLU、随机失活”。关节数$j$在原文中以符号出现。}
 \label{tab:vision-recognition-martinez-layers}
\end{table*}

训练以均方误差$(1/n)\sum_i\lVert g_{\bm{\theta}}(\bm{x}_i)-\bm{y}_i\rVert_2^2$学习（Adam，初始学习率0.001并指数衰减，批大小64，200轮），数据是带动作捕捉三维真值的室内数据集，二维输入可以是真值投影，也可以是二维检测器的输出。使用时必须复现训练时的规范化，再把输出逆规范化；关节顺序错位会把一个部位的输入当成另一个部位，维度却仍然匹配，因而不会报错。输出是相对根关节的坐标，不含人到相机的距离，要得到后者还需另估根平移。

节首的像素$(520,240)$可以对应$(1,0,4)$或$(2,0,8)$。网络用训练数据中的人体比例和姿态分布在这些解释中选一种，向前或向后弯曲的肢段也可能有相近的二维投影，选择同样来自先验。评价要说明是否允许旋转、平移和尺度对齐：相似变换对齐后的误差较小，不意味着原始输出已有正确朝向或距离。已知相机的多视图关节可以通过射线求交估计三维点，增加的是观测，而单目提升网络增加的只是先验。

\subsubsection{SAM 3D Body}
\label{sec:vision-recognition-sam3d-body}

两个三维骨架可以有相同的关节位置，却对应不同的肩宽、躯干厚度和手指形状；若要把人体轮廓投回图像比对，或驱动一个虚拟人，关节点还不够。\emph{人体网格恢复}（human mesh recovery, HMR）估计由顶点和面片组成的人体表面，同时解释姿态、体形与相机。已有方法在遮挡、罕见姿态和手部细节上容易失败，而且结果一旦出错，用户无法指出错误。SAM 3D Body的想法是把SAM式的提示接口引入人体网格恢复：模型可以直接预测，也可以接受二维关键点或掩码提示，按提示修正结果\scite{vision-sam3d-body-2026}。

输出采用参数化人体表示Momentum Human Rig（MHR），它将骨架结构、关节姿态和表面形状分开控制。记这些参数为$\bm{\beta}_{\mathrm{sk}},\bm{\theta}_{\mathrm{pose}},\bm{\beta}_{\mathrm{surf}}$，人体模型资产分别产生关节和顶点：
\[
 \bm{J}=\mathcal{J}(\bm{\beta}_{\mathrm{sk}},
                   \bm{\theta}_{\mathrm{pose}}),\qquad
 \bm{V}=\mathcal{V}(\bm{\beta}_{\mathrm{sk}},
                   \bm{\theta}_{\mathrm{pose}},
                   \bm{\beta}_{\mathrm{surf}})\text{，}
\]
其中$\bm{J}\in\mathbb{R}^{j\times3}$、$\bm{V}\in\mathbb{R}^{v\times3}$，$j,v$由模型资产规定。这两个式子只说明输出分工，不是把MHR的全部变形计算写成一个线性式。网络只需预测少量参数，输出就自动遵守人体连接结构；把关节或顶点经相机投影，才得到可与原图比对的位置和轮廓。

模型先取得人体裁剪及裁剪到原图的变换，把裁剪缩放到$512\times512$，由共享图像编码器得到空间特征；相机内参由现成的视场角估计器提供。原文给出两个规模：以ViT-H为编码器的版本约6.32亿参数，以DINOv3为编码器的版本约8.4亿参数。身体解码器的查询包含初始MHR与相机估计、辅助二维和三维关节查询，以及可选的提示。二维关键点提示写成$(x,y,\ell)$，$\ell$指出关节身份，位置编码与学习的关节嵌入使“左腕在这里”不同于一个无部位含义的点；掩码提示经卷积编码后加到图像特征上。具体地，解码器的输入词元依次是1个MHR与相机词元、$N$个关键点提示词元、$J_{2\mathrm D}$个二维和$J_{3\mathrm D}$个三维辅助关节词元，以及2个手部位置词元，宽度均为$D$；它们经交叉注意力读取图像特征后，第一个词元的输出经MLP回归姿态、形状、骨架和相机参数，再经MHR资产恢复网格。解码器的层数、宽度$D$以及MHR各参数向量的维数，论文均未给出。手部在整人裁剪中占的像素很少，模型另设可选的手部裁剪和独立手部解码器（共用编码器权重）；为避免直接替换手部参数连带改变肘部，原方法用手部解码得到的腕位置和身体解码得到的肘位置再次提示身体分支，细化后按人体连接树合并。

训练数据约700万幅图像，来自SA-1B、合成数据和多个手部与多视图数据集；三维依据来自多阶段标注：人工关键点、可微拟合、多视图几何和密集关键点检测以不同组合产生监督，另有合成和多视图数据。二维关节在裁剪坐标中监督；身体三维关节减去骨盆、手关节减去手腕后比较，只有已标注的手部参与对应目标。关节使用带学习不确定性权重的$L_1$项，姿态与形状参数使用$L_2$项，另加关节限位和手框误差。训练中随机进行多轮提示以模拟纠错，并提高用户给定关节的监督权重。

使用时人体框可由外部检测器给出，框误差属于完整流程的一部分。用户点先变到裁剪坐标，预测的关节和网格投影后再用裁剪逆变换回原图。教学中设左腕被遮住，无提示预测投回原图$(120,150)$，用户根据额外依据给点$(128,146)$，像素残差为$\sqrt{8^2+4^2}\approx8.94$。加入提示后应重跑解码、恢复网格并重新投影，检查这个残差以及肘腕连接、手形和轮廓是否共同改变；这里只定义需要检查的量，不假定提示后的误差必然下降。轮廓相符的网格仍可能把手腕放在不同深度，未知真实身高时也不能把模板的身体尺寸当成该人的测量结果。连续姿态与身份关联由\ref{chap:vision-video-recognition}章承接。

三个模型的对比见表\ref{tab:vision-recognition-pose-compare}。Martinez基线相对OpenPose改变的是输出维数，它把OpenPose这类二维结果当作输入，换来三维关节，付出的是对人体先验的依赖和单人输入的限制；SAM 3D Body再把输出从关节换成网格，并引入提示纠错，换来表面与手部，付出的是对人体模型资产、多阶段标注和裁剪流程的依赖。

\begin{table}[htbp]
 \centering
 \small
 \begin{tabularx}{\linewidth}{@{}lXXX@{}}
 \toprule
 比较项 & OpenPose & Martinez基线 & SAM 3D Body\\
 \midrule
 输入 & 整幅多人图像 & 单人二维关节 & 人体裁剪与可选提示\\
 输出 & 二维关节与人体归属 & 相对根关节的三维关节 & 参数化人体网格与相机\\
 关键机制 & 置信图加PAF组装 & 全连接残差回归 & 提示查询解码与MHR\\
 深度来源 & 不估计 & 人体姿态先验 & 图像与人体先验\\
 主要局限 & 交叉肢段与遮挡 & 依赖二维输入与对齐方式 & 尺度歧义、依赖模型资产\\
 \bottomrule
 \end{tabularx}
 \caption{三种姿态输出的比较。三者回答的不是同一个问题，比较时应固定关节集合、坐标与对齐方式。}
 \label{tab:vision-recognition-pose-compare}
\end{table}

\subsection{场景深度的估计}
\label{sec:vision-recognition-depth}

\term{深度估计}（depth estimation）的输入是一幅或多幅图像，输出是参考图每个像素的光轴深度$z$。节首的式\eqref{eq:vision-camera-projection}已经说明，单个像素只确定一条射线，射线上的$z$需要额外信息。信息来源决定了路线：单目路线只有一幅图，靠物体形状、透视、遮挡关系和训练中学到的先验推断深度，通常只能可靠地给出相对远近；多视图路线有多幅已知相机的图像，同一表面在不同图中的位置提供跨射线约束，在相机基线已知时可以得到度量深度。Depth Anything V2和MVSNet分别代表这两条路线；两者的输出都要说明是深度、逆深度还是相对量，以及是否已经确定度量尺度。

多视图路线有一段很长的经典前史，MVSNet正是在它上面做的改动，这里先回顾要点，完整推导见\ref{sec:vision-formation-multi-view}节。最简单的情形是平行双目：两台内参相同的相机沿水平方向相隔基线$b$，同一点在左右图中只在横坐标上相差视差$\Delta u$，由相似三角形得到$z=fb/\Delta u$，其中$f$是像素焦距。取$f=800$像素、$b=0.2$米、$\Delta u=40$像素，得$z=4$米；视差误差1像素约引起0.1米的深度误差，而深度越远视差越小，同样的视差误差带来的深度误差越大。经典立体匹配因此把问题化为沿极线搜索对应：对左图每个像素，在右图同一行的若干候选视差上计算匹配代价（例如窗口内灰度差的平方和或归一化互相关），得到一个“像素$\times$视差”的代价体，再加上相邻像素视差应当平滑的约束，选出代价最小的视差。多于两台相机、而且不平行时，平面扫描法把候选视差换成参考相机中的一组深度平面，每个深度平面通过单应把其他图像变换到参考视角，再比较外观是否一致。这条路线在纹理丰富处精确，但手工的匹配代价在弱纹理、反光和遮挡处不可靠，平滑约束又会抹掉细结构。MVSNet保留了平面扫描与代价体这两个经典结构，把手工代价换成学习的特征和三维卷积。

\subsubsection{Depth Anything V2}
\label{sec:vision-recognition-depth-anything}

单目深度模型需要大量带深度标签的训练图，真实数据的深度却来自深度传感器或多视图重建，往往有空洞、噪声，物体边缘尤其粗糙；用这些标签训练，模型会学到模糊的边界。合成图像有精确深度，但外观与真实照片存在差距。Depth Anything V2的想法是把两种数据的长处分开使用：只用合成图像训练一个大教师，得到精细的几何；再让教师给大量无标注真实图像产生伪标签，学生模型从这些伪标签学习真实外观\scite{vision-2406-09414}。伪标签学习的一般依据可参考\BookChapterRef{chap:data-derived-knowledge-acquisition}{卷八“基于数据本身的知识获取”章}。

网络由DINOv2编码器和密集预测变换器（dense prediction transformer, DPT）解码器组成。DINOv2的图像块边长为14\scite{vision-2304-07193}，以输入$518\times518$为例，编码器输出$37\times37$个词元；DPT从编码器的四个不同层取出词元，各自重组为二维特征图并缩放到不同分辨率，再自深而浅逐级融合、上采样，最后输出$H\times W\times1$的逐像素相对逆深度。DPT把多层空间特征融合成逐位置数值，而不是把全图汇总向量当作深度图。原文按DINOv2的四个规模训练了四个学生模型，配置见表\ref{tab:vision-recognition-dav2-config}。

\begin{table*}[tbp]
 \centering
 \small
 \begin{tabular}{@{}lccccccl@{}}
 \toprule
 模型 & 图像块大小 & 词元数 & 层数 & 宽度 & 头数 & MLP倍率 & 参数量\\
 \midrule
 Depth Anything V2-S（ViT-S） & 14 & 1369 & 12 & 384 & 6 & 未列出 & 约25M\\
 Depth Anything V2-B（ViT-B） & 14 & 1369 & 见注 & 768 & 12 & 未列出 & 论文未给出\\
 Depth Anything V2-L（ViT-L） & 14 & 1369 & 24 & 1024 & 16 & 未列出 & 约335M\\
 Depth Anything V2-G（ViT-g） & 14 & 1369 & 40 & 1536 & 24 & 未列出 & 约1.3B\\
 \bottomrule
 \end{tabular}
 \caption{Depth Anything V2的配置。参数量取自Depth Anything V2原文（含DPT解码器）；编码器的块数、宽度与头数取自DINOv2原文表17，其中ViT-B的块数在该表中印为18，与ViT-B通常的12层不同，此处不作推断。DINOv2原文只说明ViT-S/B/L用普通MLP、ViT-g用SwiGLU作前馈网络，未列出前馈宽度。词元数按$518\times518$输入、块边长14计为$37\times37$。}
 \label{tab:vision-recognition-dav2-config}
\end{table*}

训练分三步。教师（最大规模的编码器）只用约59.5万张合成RGB和精确深度训练，批大小64、16万次迭代；教师为大量真实无标注图像预测伪深度；不同规模的学生只用这些约6200万张真实图与伪标签对训练，批大小192、48万次迭代，编码器与解码器的学习率分别为$5\times10^{-6}$和$5\times10^{-5}$。基础检查点预测尺度、平移不变的相对逆深度：比较预测与标签之前，先消去两者之间的整体尺度和偏移。设预测为$\hat q_p$、目标为$q_p$，一个教学用的平方对齐目标为
\begin{equation}
 \min_{a,b}\frac1{|\Omega|}
       \sum_{p\in\Omega}(a\hat q_p+b-q_p)^2\text{，}
 \label{eq:vision-relative-depth-alignment}
\end{equation}
其中$\Omega$是有效监督位置，$a,b$是比较时的对齐量。令导数为零，若预测有非零方差，得$a=\operatorname{Cov}(\hat q,q)/\operatorname{Var}(\hat q)$、$b=\bar q-a\bar{\hat q}$；全图预测为常数时尺度无从估计。这个平方例子解释对齐消去了什么信息，原论文的尺度平移不变损失在具体形式上与此不同。训练还加入梯度匹配项，比较相邻位置残差的差，防止整体对齐合适而边界过于平滑。尺度平移不变项与梯度匹配项的权重比为1比2。对伪标签样本，原工作忽略损失最大的10\%区域以降低噪声影响。例如预测$(1,2,3)$与目标$(3,5,7)$在$a=2,b=1$下完全重合；对齐误差为零并不说明原预测的数值已经是逆米单位。

使用相对检查点时，按模型尺寸和归一化处理图像，解码出$\hat q$，插值回原图有效区域，并保留其相对逆深度身份。若三个位置的相对逆深度为$(0.8,0.4,0.2)$，可以判断第一个最近，却不能未经校准说深度为$(1.25,2.5,5)$米；颜色图只是显示映射，也不能从色值读米。度量适配用带单位的深度标签微调已有网络，使输出在该数据、最大深度和相机条件下对应实际距离；即便如此，换相机或场景仍可能带来尺度误差。

\subsubsection{MVSNet}
\label{sec:vision-recognition-mvsnet}

单目深度依赖先验，遇到训练中少见的场景结构就没有依据。如果有多张重叠照片且相机已知，同一表面在不同图中的位置可以直接约束深度。传统多视图立体方法用手工设计的相似度比较图像块，在弱纹理和反光区域表现差，也难以融合多张图的证据。MVSNet的想法是在学习的特征上做这件事：对参考图的每个像素枚举一组候选深度，把各源图的特征按候选深度变换到参考视角，用方差度量多视图是否一致，形成代价体，再用三维卷积从中推断深度\scite{vision-1804-02505}。

每个候选深度$d_k$对应参考相机中$z=d_k$的一个平面。参考像素$\tilde{\bm p}_r$在该深度下为$d_k\bm K_r^{-1}\tilde{\bm p}_r$，进入源相机得到$\bm R_{s,r}d_k\bm K_r^{-1}\tilde{\bm p}_r+\bm t_{s,r}$。由于$\bm e_3^\top\bm K_r^{-1}\tilde{\bm p}_r=1$，其中$\bm e_3=(0,0,1)^\top$，除去齐次比例$d_k$得到
\begin{equation}
 \begin{aligned}
 \tilde{\bm p}_s&\sim\bm H_{s,r}(d_k)\tilde{\bm p}_r,\\
 \bm H_{s,r}(d_k)=
 \bm K_s\left(\bm R_{s,r}
 +\frac{\bm t_{s,r}\bm e_3^\top}{d_k}\right)\bm K_r^{-1}\text{。}
 \end{aligned}
 \label{eq:vision-mvsnet-homography-process}
\end{equation}
这个映射对整张平面成立，因而可以用一个$3\times3$矩阵变换整幅源特征图，并对$d_k$可微。特征图尺度不同于图像时，先按步幅调整内参，再用双线性插值读取源特征；越界或落到相机后方的位置不得当成有效观测。

网络结构如下。$n_{\mathrm v}$幅图（1幅参考、其余为源图）经共享的8层二维卷积网络，各自得到$H/4\times W/4\times32$的特征图。对$D$个候选深度，把每幅源特征图按式\eqref{eq:vision-mvsnet-homography-process}变换到参考视角，得到$n_{\mathrm v}$个$D\times H/4\times W/4\times32$的特征体。将参考特征和变换后的各源特征记为$\bm{f}_r$，逐元素计算方差：
\begin{equation}
 \begin{aligned}
 \bar{\bm{f}}&=\frac1{n_{\mathrm{v}}}\sum_{r=1}^{n_{\mathrm{v}}}\bm{f}_r,\\
 \bm{c}(d_k,u,v)&=\frac1{n_{\mathrm{v}}}\sum_{r=1}^{n_{\mathrm{v}}}
   (\bm{f}_r-\bar{\bm{f}})\odot(\bm{f}_r-\bar{\bm{f}})\text{，}
 \end{aligned}
 \label{eq:vision-mvs-cost}
\end{equation}
其中$\odot$表示逐元素相乘。方差对视图数没有限制，任意多幅源图都可以合成一个\term{代价体}（cost volume），即形状为$D\times H/4\times W/4\times32$、记录每个位置在每个候选深度下多视图有多不一致的张量。三维U形卷积网络在空间和深度三个方向上整理代价体，输出单通道的$D\times H/4\times W/4$得分，沿深度轴做softmax得到概率$\pi_k$，以期望$\hat d=\sum_k\pi_kd_k$回归深度。这种期望回归对概率可微，结果也不必等于某个离散候选。最后一个浅层细化网络读取参考图像和初步深度，预测残差以修正边界。各段的结构与形状见表\ref{tab:vision-recognition-mvsnet-structure}。

\begin{table*}[tbp]
 \centering
 \small
 \begin{tabularx}{\linewidth}{@{}llXl@{}}
 \toprule
 部分 & 输入形状 & 结构与层数 & 输出形状\\
 \midrule
 特征提取（每图） & $512\times640\times3$ & 8层二维卷积，第3、6层步长为2，每个尺度2层；除末层外均接批归一化与ReLU；各视图共享权重 & $128\times160\times32$\\
 单应变换与代价 & $n_{\mathrm v}$张特征图 & 对$D$个深度平面按式\eqref{eq:vision-mvsnet-homography-process}变换，按式\eqref{eq:vision-mvs-cost}求方差 & $D\times128\times160\times32$\\
 代价整理 & 代价体 & 四个尺度的三维U形网络，第一层三维卷积把通道从32降到8，每个尺度2个卷积层，末层输出1通道 & $D\times128\times160$\\
 深度回归 & 得分体 & 沿深度轴softmax，求期望深度 & $128\times160$\\
 细化 & 初步深度与缩放的参考图（4通道） & 3个32通道二维卷积加1个单通道卷积，输出残差 & $128\times160$\\
 \bottomrule
 \end{tabularx}
 \caption{MVSNet的分段结构。层数与通道取自原文正文；形状按原文DTU训练设置推算：输入裁剪为$640\times512$，$n_{\mathrm v}=3$，深度假设从425毫米到935毫米、间隔2毫米，共$D=256$个。特征网络各中间层的通道数在正文中未给出。}
 \label{tab:vision-recognition-mvsnet-structure}
\end{table*}

同一像素的教学标量特征在正确假设下为$(1,1.1,0.9)$，方差约0.0067；在错误假设下读到$(1,2,3)$，方差为$2/3$。若候选深度为2、4、6米，概率为0.1、0.7、0.2，回归得到4.2米。双峰概率会把结果放在两个表面之间，因而概率的集中程度也应作为置信度检查。

训练数据包含参考图、若干重叠源图、已知相机和有效参考深度（原文使用DTU室内物体数据集，约2.7万个训练样本，测试时用5幅图、分辨率$1600\times1184$）。在有效像素上以$L_1$损失比较初步深度和细化深度与真值，两项相加，更新二维特征网络、三维代价整理网络和细化网络。使用时相机由当前图集提供；得到深度后按概率和多视图几何一致性过滤，再反投影合并点云。代价体的存储为$O(Dh'w'd)$，随候选数和分辨率增长，这是高分辨率重建的主要瓶颈。与本小节开头的平行双目一样，已知的相机基线正是MVSNet输出度量尺度的来源。

\begin{figure}[htbp]
 \centering
 % 针孔投影的几何关系；示意深度候选，不按物理长度绘制。
\begin{tikzpicture}[font=\figurefont\small,x=1cm,y=1cm,
 ray/.style={draw=BookBlue,line width=0.9pt},
 alt/.style={draw=BookAmber,dashed,line width=0.9pt}]
 \coordinate (a) at (0,0);
 \coordinate (b) at (7,0);
 \coordinate (p) at (2,3);
 \coordinate (q) at (3,4.5);
 \draw[ray] (a)--(q);
 \draw[ray] (b)--(p);
 \draw[alt] (b)--(q);
 \draw[draw=BookInk,line width=1.2pt] (-0.3,1)--(2.2,1);
 \draw[draw=BookInk,line width=1.2pt] (4.5,1)--(7.5,1);
 \fill[BookInk] (a) circle(1.1mm);
 \fill[BookInk] (b) circle(1.1mm);
 \fill[BookTeal] (p) circle(1.1mm);
 \draw[BookAmber,fill=white,line width=0.9pt] (q) circle(1.1mm);
 \fill[BookBlue] (0.667,1) circle(1mm);
 \fill[BookBlue] (5.333,1) circle(1mm);
 \draw[BookAmber,fill=white,line width=0.9pt] (6.111,1) circle(1mm);
 \node[below] at (a) {参考相机};
 \node[below] at (b) {源相机};
 \node[above right] at (0.667,1) {同一参考像素};
 \node[below] at (5.333,0.9) {$u'_1$};
 \node[below] at (6.111,0.9) {$u'_2$};
 \node[left] at (p) {深度$d_1$};
 \node[right] at (q) {深度$d_2$};
 \node[right] at (7.5,1) {源图平面};
 \draw[draw=BookMuted,<->,line width=0.7pt] (0,-0.8)--(7,-0.8)
   node[midway,below] {已知相机相对位置};
\end{tikzpicture}

 \caption{已知相机下，参考像素沿射线形成不同深度假设，分别投到源图的不同位置。图为针孔几何教学示意，只表示对应关系，不按物理长度绘制。跨图比较可排除不一致的假设，但遮挡或重复纹理会使判断不唯一。}
 \label{fig:vision-depth-hypotheses}
\end{figure}

图\ref{fig:vision-depth-hypotheses}展示两个深度假设如何查询不同的源图位置。弱纹理使多个候选代价相近，遮挡使源图看到另一表面，反光破坏外观一致性，相机误差使正确深度也无法对齐。两条路线在同一街景上的比较可以概括为：Depth Anything V2只要一幅图，覆盖所有像素，边界细致，但默认输出相对量，尺度来自先验或额外适配；MVSNet需要多幅图和已知相机，在相机基线已知时给出度量深度，但在弱纹理、遮挡区域可能无有效结果，代价体的存储也限制分辨率。

\subsection{相机与场景几何的估计}
\label{sec:vision-recognition-geometry}

MVSNet把相机当作已知条件。若手头只有一组普通照片，相机未知，就要同时估计每张图的内外参和场景的三维结构。输入是一组RGB图像，输出是每图的相机、深度或点云，以及它们所在的公共坐标系。难点在于相机和结构相互依赖：知道相机才能三角化点，知道点才能求相机。经典的回答是\ref{sec:vision-formation-multi-view}节的运动恢复结构，这里回顾它的步骤和失效条件。第一步在每幅图上检测SIFT等局部特征，在图对之间匹配，并用RANSAC估计基础矩阵或本质矩阵来剔除错误匹配；第二步从一对重叠充分的图开始，分解本质矩阵得到相对位姿，三角化出初始三维点；第三步逐幅加入新图，用已有三维点与新图特征的对应求新相机的位姿，再三角化新点。每加入若干幅图，就做一次光束法平差：对相机参数$\{\bm K_i,\bm R_i,\bm t_i\}$和三维点$\{\bm X_j\}$联合最小化重投影误差$\sum_{(i,j)}\lVert\bm p_{i,j}-\pi(\bm K_i,\bm R_i,\bm t_i,\bm X_j)\rVert_2^2$，其中$\bm p_{i,j}$是第$j$个点在第$i$幅图中的观测，$\pi$按式\eqref{eq:vision-camera-projection}投影。这个流程的每一步都有明确的几何意义，在纹理丰富、重叠充分时可以达到很高的精度；但它是一条串联的流水线，任何一步失败都会向后传递。弱纹理使特征点太少，重复图案使匹配出错，视图太少或基线太短使三角化病态，动态物体又违反静态场景的假设。

本小节的五个模型都用前馈网络直接预测几何，差别在于一次处理多少视图、信息怎样到达。DUSt3R一次只看两张图，直接输出两张图在同一坐标系中的点图，多张图则靠逐对预测加全局优化拼起来。VGGT让任意多张图在一次前向中相互注意，直接输出所有相机和深度，省去了逐对再优化，但注意力成本随视图数平方增长。VGG-T$^3$针对这一成本，把全局注意力中随视图增长的键值压缩到一个固定大小的网络里。InfiniteVGGT面对另一种条件：视图逐帧到达，不能等全集，它为流式处理维护一个有限的缓存。FF3R则在几何之外同时恢复可查询的语义，使三维结构能回答“椅子在哪里”。

\subsubsection{DUSt3R}
\label{sec:vision-recognition-dust3r}

上面的经典流程中，特征点匹配不足，相机就估不出来，后续一切都无从谈起。DUSt3R的想法是跳过显式的相机估计，让网络对一对图像直接回归\term{点图}（point map）：每个有效像素对应一个三维点，两张点图都表达在第一幅图的相机坐标中\scite{vision-2312-14132}。点图同时隐含了深度、两图的相对位姿和像素对应，相机可以事后从点图中恢复。

两幅图各自切块后经共享的ViT编码器，得到两组词元。两个Transformer解码分支每层先做自注意力，再对另一分支的词元做交叉注意力，从而反复交换跨图信息。每个分支最后由DPT式头输出$H\times W\times3$的点图和$H\times W\times1$的置信度；第二张图的点图同样表达在第一相机坐标中，这一约定迫使网络在内部推断两图的相对位姿。编码器和解码器从CroCo跨视图补全预训练初始化。原文的编码器为ViT-Large规模，解码器为ViT-Base规模，头部最终采用DPT；块边长、解码块数和参数总量论文未给出。训练分三段：先在$224\times224$上用线性头训练，再改用最长边512像素、长宽比随批次随机变化的图像，最后换成DPT头。训练数据约850万对图像，来自室内、室外、物体和驾驶等八个数据集。一对图像的前向在H100上约需40毫秒。

训练记录含同一场景的图像对，以及由深度、内参和外参得到的真实点图和有效位置；使用时不输入相机，并不意味着训练不曾使用几何标签。原工作分别用有效点到原点的平均距离规范化预测和目标，得到尺度无关的点误差$e_i$。置信加权目标为$\sum_i(c_i e_i-\alpha\log c_i)$，$\alpha>0$，置信度以$c_i=1+\exp h_i>1$参数化。对固定$e_i>0$，该项对$c_i$的导数为$e_i-\alpha/c_i$，误差大的点倾向于较小的置信，而对数项使网络不能把所有置信都压到最小。以$\alpha=1$为例，$e=0.1$时驻点为$c=10$；$e=2$时驻点0.5不在允许范围，便取下边界。这是可靠性权重，不是校准过的正确率。天空等无效区域按标签排除。

多于两张图时，每一对都以自己的首图为参考，不能直接拼接。DUSt3R把有足够重叠的图对连成图，边$e$保存两张点图$\bm X^{v,e}$及置信；对每幅视图求共同点图$\bm\chi^v$，对每条边求尺度$s_e>0$与刚体变换$(\bm R_e,\bm t_e)$，使$\sum_{e,v,p}c_p^{v,e}\lVert\bm\chi_p^v-s_e(\bm R_e\bm X_p^{v,e}+\bm t_e)\rVert_2$最小，并固定尺度规范以防所有点缩到零。这是对场景变量的迭代优化，网络权重不变。相机内参可以从点图拟合：例如假设主点在中心，由$(u-c_x,v-c_y)\approx f(x/z,y/z)$对高置信点拟合$f$。若甲乙图预测建筑宽度为2个单位，乙丙图为5个单位，对齐要先估计约2.5的相对尺度；没有额外度量依据时，共同场景只有约定尺度。逐对预测、全局对齐两个阶段使处理大量图像时计算重复，弱重叠和动态物体也会干扰对齐。

\subsubsection{VGGT}
\label{sec:vision-recognition-vggt}

DUSt3R一次只看两张图，$n_{\mathrm v}$张图需要处理许多对，再做一次全局优化；每一对看不到其他视图的证据，三张图共同才能消除的歧义在逐对阶段无法利用。VGGT的想法是让所有视图在一次前向中共同参与计算，直接输出每张图的相机、深度、点图和跟踪特征，不再依赖后续优化\scite{vision-2503-11651}。

每张图经DINOv2初始化的编码器切成带位置的图像块词元，并附一个相机词元和若干辅助寄存词元；首图使用与其他图不同的学习词元，标记它为坐标参照。主干由交替注意力组成：帧内层只在同一张图的词元之间注意，保留各图的结构；全局层在所有视图的全部词元之间注意，交换跨视角证据。两种层交替堆叠，相机词元因此始终与所属图像关联。最后，相机头读取相机词元，输出九维编码（四维旋转四元数、三维平移和横纵视场角）；DPT式稠密头把多层图像特征恢复为空间网格，输出深度、点图、置信和跟踪特征。视场角$\varphi_x,\varphi_y$在中心主点条件下给出$f_x=w/[2\tan(\varphi_x/2)]$、$f_y=h/[2\tan(\varphi_y/2)]$，因此得到的是当前网络输入网格上的内参。深度$D_i$是第$i$相机自身的光轴深度，点图和外参则用首相机作为公共参照，二者可以互相换算：
\begin{equation}
 \bm{P}_{1,i}(u,v)=
 \bm R_i^\top\left[
 D_i(u,v)\bm K_i^{-1}(u,v,1)^\top-\bm t_i\right]\text{。}
 \label{eq:vision-vggt-unprojection-process}
\end{equation}
这里$\bm R_i,\bm t_i$沿用世界到第$i$相机的约定，公共世界就是首相机。

表\ref{tab:vision-recognition-ff-geometry-config}给出VGGT的配置。主干共24个注意力块，每块含一个帧内自注意力层和一个全局自注意力层，即帧内与全局各24层；宽度1024、16个头，沿用DINOv2中ViT-L的设置，并使用QK归一化与LayerScale，全网没有交叉注意力。单个块的计算可以按ViT的块写法表示为两次Transformer更新（原文未说明层归一化的具体位置，这里按ViT的前置写法）：先把$n_{\mathrm v}$幅图的词元按图分组，对每组做$\bm z'=\bm z+\operatorname{MSA}(\operatorname{LN}(\bm z))$、$\bm z''=\bm z'+\operatorname{MLP}(\operatorname{LN}(\bm z'))$；再把全部图的词元拼成一个序列，做同样的两步。帧内层与全局层的差别只在注意力的范围。相机头在相机词元上再加4个自注意力层和一个线性层；DPT头取第4、11、17、23块的词元。每帧除图像块词元外还有1个相机词元和4个寄存词元，寄存词元的输出被丢弃。

\begin{table*}[tbp]
 \centering
 \small
 \begin{tabular}{@{}lccccccl@{}}
 \toprule
 模型 & 图像块大小 & 每图词元数 & 层数 & 宽度 & 头数 & MLP倍率 & 参数量\\
 \midrule
 DUSt3R编码器 & 论文未给出 & — & ViT-Large & — & — & — & 论文未给出\\
 DUSt3R解码器 & — & — & ViT-Base（两分支） & — & — & — & 论文未给出\\
 VGGT主干 & 14 & $37^2+5$ & 帧内24＋全局24 & 1024 & 16 & 未列出 & 约12亿（全模型）\\
 VGG-T$^3$ & 14 & 同VGGT & 同VGGT & 1024 & — & 快权重4 & 论文未给出\\
 \bottomrule
 \end{tabular}
 \caption{前馈几何模型的配置。DUSt3R原文只给出编码器与解码器的规模名称；VGGT的“每图词元数”按最长边518、正方形输入计为$37\times37$个图像块词元加1个相机词元与4个寄存词元。VGG-T$^3$沿用VGGT的层数，只把每个全局注意力层换成测试时训练模块，故不再有注意力头；其快权重是输入输出1024维、中间层4倍宽的SwiGLU网络。}
 \label{tab:vision-recognition-ff-geometry-config}
\end{table*}

训练记录是同一场景的多张RGB、每图相机、有效深度和公共坐标点图，跟踪监督还包含查询点的跨图位置及可见性；这些几何来自传感器、合成引擎或运动恢复结构。训练先把全部量变到首相机坐标，再用有效点到原点的平均距离除相机平移、点坐标和深度。与DUSt3R在比较时才归一化预测不同，VGGT要求网络直接输出这套规范化的结果。相机编码用Huber误差监督（小残差为平方、大残差为线性），深度和点图用置信加权的位置误差与空间梯度误差监督，跟踪用位置误差和二元可见性误差监督。主干与任务头联合训练，每个场景随机采样2至24幅图，跟踪损失的权重为0.05；优化器为AdamW，共16万次迭代，在64块A100上训练约9天。

使用时，一次前向得到全部相机和稠密量。教学中首相机为单位旋转、零平移，主点$(320,240)$、焦距800、像素$(520,240)$与深度4个规范单位恢复点$(1,0,4)$，仍不能将4个单位写成4米。可用式\eqref{eq:vision-vggt-unprojection-process}从深度构造公共点，也可直接读取点图；两种输出不应未经核验混拼。跟踪接口把原图查询点变到处理网格，在查询图特征中插值取其描述，再与各图跟踪特征相关，估计跨图位置和可见性。全局注意力比较$(n_{\mathrm v}p)^2$个词元对（$p$为每图词元数），图像数增多时计算和显存都迅速增长，这是下一个模型要解决的问题。

\subsubsection{\texorpdfstring{VGG-T$^3$}{VGG-T3}}
\label{sec:vision-recognition-vgg-t3}

建筑的照片从三张增到数百张时，VGGT的全局注意力成为瓶颈：$n_{\mathrm v}$幅图、每图$p$个词元，比较次数为$(n_{\mathrm v}p)^2$，键和值本身也占$O(n_{\mathrm v}pd)$的空间。高效注意力实现可以不保存完整得分矩阵，却不能减少比较次数。VGG-T$^3$的想法是在使用时把每个全局层中随视图增长的键值集合压缩到一个固定大小的多层感知机（multi-layer perceptron, MLP）中：这个MLP把键映射为对应的值，其参数就保存了当前场景；之后查询只需送入MLP读取，计算量随词元数线性增长\scite{vision-vgg-t3-2026}。

\term{测试时训练}（test-time training, TTT）指在处理当前输入时临时更新一小组参数。对一个全局层，设当前图集产生$N=n_{\mathrm v}p$个键$\bm{k}_i$、查询$\bm{q}_i$和目标值$\bm{v}'_i$，维数均为$d$；场景权重为$\bm{\phi}$，MLP记为$T_{\bm{\phi}}$。内层目标和读取为
\begin{equation}
 \begin{aligned}
 \bm{\phi}_*&\approx\argmin_{\bm{\phi}}
   \sum_{i=1}^{N}
   \ell\bigl(T_{\bm{\phi}}(\bm{k}_i),\bm{v}'_i\bigr),\\
 \bm{o}_i&=T_{\bm{\phi}_*}(\bm{q}_i)\text{。}
 \end{aligned}
 \label{eq:vision-vgg-t3-fast-weights}
\end{equation}
其中$\ell$是逐词元的拟合损失，衡量MLP从键读出的结果与目标值是否一致；近似号表示实际只做规定次数的更新。原文实现沿用一种已有的大块测试时训练设计，把$\ell$取为$T_{\bm\phi}(\bm k_i)$与$\bm v'_i$的内积形式；下面的教学例为了便于手算改用平方误差，两者都是对各词元求和的损失，这一点是后文分批计算的依据。目标值$\bm{v}'_i$中图像块部分先在各图网格上经过短二维卷积混合邻域，相机和寄存词元不经过这项卷积；邻域混合使拟合目标带有局部上下文，避免只是复现同一词元的两个线性投影。查询和键还做$L_2$归一化。

模型以VGGT权重初始化，保留图像编码、帧内层和预测头，把全局注意力层换成上述测试时训练模块，并适配其投影和新增参数。每层先从训练得到的初始化取得场景权重，按键值残差计算梯度，更新后把查询送入MLP，所得词元进入下一层；全部层完成后由相机、深度和点图头解码。快权重MLP是输入输出均为1024维、中间层宽度为4倍的SwiGLU网络；图像块的值在各图网格上经一个$3\times3$的二维卷积混合邻域。原文内层采用Muon优化器（5次Newton–Schulz迭代、内层学习率0.1），训练时内层更新一步，测试默认更新两步，更多步数没有进一步改善。外层只训练全局层的查询、键、值和输出投影以及新增参数，VGGT原有的其他参数全部冻结；在8块A100上训练10万步，原文估计约为VGGT训练成本的12\%。外层训练沿用VGGT的几何监督，使初始化和投影学会“适合快速拟合”的表示。

求和形式的内层目标允许把全集分批拟合。固定一次更新开始时的$\bm{\phi}$，把词元划成批次$B_s$，则$\nabla_{\bm{\phi}}L=\sum_s\sum_{i\in B_s}\nabla_{\bm{\phi}}\ell(T_{\bm{\phi}}(\bm{k}_i),\bm{v}'_i)$。这一等式来自有限和的微分线性性，要求各批在同一组权重上计算梯度并累加后再更新；逐批改权重就不是同一次全集梯度。未处理的批可以放在主机内存或分到多个设备，降低单卡显存。用标量线性模型$T_\phi(k)=\phi k$和平均平方误差作教学简化，两个键值样本为$(1,2)$、$(2,4)$；在$\phi=0$处梯度为$-10$，步长0.1更新到$\phi=1$，再更新一次到1.5，查询1.5读出2.25，尚未达到精确映射给出的3。有限次拟合与精确保存全部键值之间的差距，就是固定容量带来的近似误差。原文在其基准中指出相机预测较不精确，并采用直接点图头评估点图；原图、编码特征和输出点图的存储仍随视图数增长。在A100上，原文摘要报告重建1000幅图约需54秒，比使用全局注意力的基线快11.6倍。

\subsubsection{InfiniteVGGT}
\label{sec:vision-recognition-infinitevggt}

前两个模型都假定完整图集已经到齐。沿走廊行走时，第$t$幅图已经到达，未来的画面尚未发生，当前输出不能借未来视图纠正自身。\emph{因果几何预测}（causal geometry prediction）限制当前帧只读取自身与已到达的历史。直接保留每层所有历史帧的键值可以复用旧特征，但缓存随$t$无限增长；只保留最近几帧，又会删去初始参照与回访线索。InfiniteVGGT的想法是维护一个固定预算的滚动缓存，按键的多样性决定保留哪些历史，并始终保留首帧作为锚点\scite{vision-infinitevggt-2026}。

InfiniteVGGT建立在已有的因果几何模型StreamVGGT上，不重新训练权重，因此结构与VGGT式模型相同（表\ref{tab:vision-recognition-ff-geometry-config}），只是全局注意力改成只读过去的因果注意力：图像仍经帧内和因果跨帧注意力，输出相机、深度、点图和跟踪特征；变化只在于哪些历史键值可读。各层各头把第一帧的键值固定为锚点，因为几何输出以首相机为参照；其余历史与新到词元构成可裁减的候选集。相邻画面含大量重复内容，原方法用归一化键的方向差异选择保留内容，从而不必计算完整注意力矩阵。对一个层和头，令候选键的单位向量为$\hat{\bm{k}}_i$，平均向量为$\bm{\mu}=m^{-1}\sum_i\hat{\bm{k}}_i$，多样性评分为
\begin{equation}
 s_i=-\frac{\hat{\bm{k}}_i^\top\bm{\mu}}
                 {\lVert\bm{\mu}\rVert_2}\text{。}
 \label{eq:vision-infinitevggt-diversity}
\end{equation}
与平均方向差别大的键得分高，与其值一起保留。这个评分近似识别冗余，不直接度量三维信息量；平均向量为零时需要数值退化处理。不同层的多样性不同，原文把层平均评分$\bar s_\ell$经softmax变成预算比例$b_\ell=\exp(\bar s_\ell/\tau)/\sum_j\exp(\bar s_j/\tau)$，按总预算$B$分配各层保留数，再按分数选取键值，最后加回首帧锚点。每到一帧，执行一次追加、评分、裁减、保留；整个过程不更新网络权重。每帧约新增1000个词元，不加裁减时100帧内就累积约$10^5$个。原文消融中每层每头的初始预算取1万、2.5万和5万个词元，在2.5万附近效果趋于饱和；温度$\tau$的取值论文未给出。在单块A100上，用余弦相似度评分时每帧约0.17秒、峰值显存约14.5GB，改用注意力权重评分则每帧多约120毫秒。

教学中令候选键方向为$(1,0)$、$(1,0)$、$(0,1)$，平均为$(2/3,1/3)$；前两者与平均方向的余弦为$2/\sqrt5\approx0.894$，第三者为$1/\sqrt5\approx0.447$，第三者评分更高。预算只容两项时，它和一个重复方向留下。这说明裁减能保留不同方向，却没有证明被删的方向在之后回访时不重要。固定缓存控制的是工作状态，每帧输出的点图若全部写盘，总存储仍随时间增长。走廊返回起点时，应比较短程相机误差、累计漂移和回访位置重合；保留首帧有助于维持参照，但不是显式的回环检测与全图校正。原文为此构建的Long3D基准有5个序列，长度从约2100帧到约9500帧；这些有限序列可以提供长程测试，不能证明无限时长下的稳定。

\subsubsection{FF3R}
\label{sec:vision-recognition-ff3r}

房间两幅照片都标出“椅子”，还不能知道它们是否指向同一张椅子的同一表面。前四个模型只恢复几何，要回答“椅子在哪里”，常见做法是先重建几何，再把二维语义分割结果投到三维；但几何错位会把椅面语义放到墙上，逐场景优化又很慢。FF3R的想法是在一次前馈中同时预测几何、外观和开放词汇语义，并用跨视角的语义一致性反过来约束几何\scite{vision-ff3r-2026}。

输入是至少两幅多视图图像。几何分支沿用VGGT式的交替注意力词元；语义分支用冻结的CLIP-LSeg编码器取得逐位置、与文字对齐的语义特征。只在几何主干的最后一层，几何词元作为查询、语义词元作为键和值做一次交叉注意力，使几何解码获得语义上下文。解码器输出深度和相机，并为每个像素预测一个三维高斯：中心由深度反投影确定，协方差描述空间扩展，颜色参数描述外观，不透明度描述遮挡贡献，另附语义向量$\bm{f}_i$。高斯是可投影、可聚合的场景表示，不应把其中心直接当作精确的物体表面。

为用图像监督这些三维属性，要把它们投回观察相机。每个高斯的中心用节首的投影得到二维位置，协方差经投影的局部线性近似成为像平面上的椭圆；按深度先后计算对像素的贡献。若第$i$项在该像素的合成系数为$a_i\in[0,1]$，前方透过比例为$T_i=\prod_{j<i}(1-a_j)$，则颜色和特征分别按$\sum_iT_i a_i\bm c_i$与$\sum_iT_i a_i\bm f_i$汇总。两者共用贡献权重，因此语义特征的监督也会反向影响所预测的几何。完整的新视角渲染在\ref{sec:vision-generation-novel-view}节展开。

若两视图看到同一表面，它们的语义特征应在几何对应下相容。使用式\eqref{eq:vision-relative-camera-measurement}，将目标图像像素$\tilde{\bm{p}}_t$在预测深度$D_t$下反投影，再投到上下文图$c$：
\begin{equation}
 \bm{p}_c=
 \Pi\!\left[
 \bm{K}_c\left(
 \bm{R}_{c,t}D_t(\bm{p}_t)\bm{K}_t^{-1}\tilde{\bm{p}}_t
 +\bm{t}_{c,t}\right)\right]\text{，}
 \label{eq:vision-ff3r-warp}
\end{equation}
其中$\Pi$用第三坐标除前两坐标。按$\bm{p}_c$插值采样源特征，计算$1-\cos(\bm{f}_t,\bm{f}_{c\to t})$，只在投影落入图像且深度相容的位置求平均，再交换两图求另一方向；边界和遮挡筛选避免把不曾看见同一表面的位置强行拉到一起。

训练从同一场景采样上下文视图，几何主干和深度头由VGGT权重初始化，新增的融合与输出部分随机初始化，CLIP-LSeg保持冻结。完整目标包含渲染RGB的像素与感知误差、渲染特征与CLIP-LSeg特征的余弦误差、上述双向跨视角特征误差，以及以VGGT输出为伪标签的相机和深度蒸馏项（深度只在高置信位置比较）。原文的总目标为$L_{\mathrm{rgb}}+0.1L_{\mathrm{feat}}+0.1L_{\mathrm{warp}}+L_{\mathrm d}+10L_{\mathrm p}$，其中$L_{\mathrm{rgb}}$是$L_1$误差加0.05倍的感知误差，$L_{\mathrm p}$是相机位姿的Huber蒸馏误差。训练用DL3DV-10K中的约6500个场景、$448\times448$的图像，每次至多16个视图，在8块A100上约两天；高斯头的输出通道、语义特征维数和参数总量论文未给出。因此“无需人工相机、深度和语义标注”不等于只有RGB残差，也不等于没有预训练监督。视图密集时，可按三维位置把高斯分到体素，用几何置信$C_i$和与体素平均语义的余弦距离$d_i^{\mathrm{sem}}$共同加权：
\begin{equation}
 \begin{aligned}
 w_{i\to s}=
 \frac{\exp(C_i-\lambda d_i^{\mathrm{sem}})}
 {\sum_{j:\,v(j)=s}\exp(C_j-\lambda d_j^{\mathrm{sem}})},\qquad
 \bar{\bm{a}}_s=\sum_{i:\,v(i)=s}w_{i\to s}\bm{a}_i\text{，}
 \end{aligned}
 \label{eq:vision-ff3r-voxel}
\end{equation}
其中$v(i)$是体素编号，$\bm{a}_i$可以是位置、颜色或语义。若两项置信均为1、语义距离为0与1、$\lambda=1$，权重为$e/(e+1)\approx0.731$与0.269，偏离项贡献下降但未消失。使用时，把“椅子”编码到同一语义空间，与三维或投影后的特征比较，再按阈值恢复区域。仅颜色渲染自然，不能证明椅子的语义或尺度正确；未见的背面仍只有先验预测。作者项目页标注其为CVPR 2026 Findings论文。

五个模型的对比见表\ref{tab:vision-recognition-geometry-compare}，比较固定在同一组建筑或房间照片上。VGGT相对DUSt3R改变了一次处理的视图数，换来单次前向和多视图共同推断，代价是注意力随视图数平方增长；VGG-T$^3$把全局注意力换成测试时拟合的固定容量映射，换来线性增长，代价是近似误差和测试时的额外优化步；InfiniteVGGT改变的是信息到达方式，换来流式处理和固定工作内存，代价是被裁减的历史无法再读取；FF3R在几何上增加语义，换来可查询的三维表示，代价是更复杂的训练目标并依赖VGGT伪标签。

\begin{table*}[tbp]
 \centering
 \small
 \begin{tabularx}{\linewidth}{@{}lXXXX@{}}
 \toprule
 模型 & 视图的处理方式 & 主要输出 & 跨视图信息的保存 & 主要代价或局限\\
 \midrule
 DUSt3R & 逐对前馈，再全局对齐 & 成对点图与置信 & 每对的交叉注意力 & 对数随视图数增长，需额外优化\\
 VGGT & 全部视图一次前馈 & 相机、深度、点图、跟踪 & 全局注意力的全部键值 & 计算随视图数平方增长\\
 VGG-T$^3$ & 全部视图，测试时拟合 & 同VGGT，点图为主 & 每层一个固定容量MLP & 拟合近似误差，相机较弱\\
 InfiniteVGGT & 逐帧到达，因果处理 & 每帧相机与点图 & 首帧锚点加多样性缓存 & 删去的历史不可恢复，长程漂移\\
 FF3R & 多视图一次前馈 & 几何、三维高斯与语义 & 几何读取语义词元 & 训练目标复杂，依赖伪标签\\
 \bottomrule
 \end{tabularx}
 \caption{五种前馈几何模型的比较。尺度均为规范化的约定尺度，若无额外度量依据，都不直接给出米制结果。}
 \label{tab:vision-recognition-geometry-compare}
\end{table*}

\section{图像匹配与检索}
\label{sec:vision-recognition-retrieval}

前几节把一幅图像变成类别、范围和测量。检索的输入是一幅查询图像（或其中一块区域）和一个大图库，输出是按相关程度排序的图库图像列表。它面临两个相互牵制的难点。第一，什么算“相关”要事先规定：找相似的街景和找同一辆车的其他照片，正例完全不同。第二，图库可能有百万张图，逐张做精细的局部匹配与几何验证不可承受；只用一个紧凑向量比较又很快，却容易被外观相似的其他实例混淆。

常用的解决方式是按步骤分工。先规定检索相关性，确定正例、忽略项和评价协议；再用代价低的表示从全库中召回少量候选，宁可多召回，不可漏掉；最后只对这些候选做代价高的局部匹配和几何验证，重新排序。三个步骤的顺序反映了信息由粗到细、成本由低到高的安排，后一步只能在前一步留下的候选中纠错，不能找回前一步漏掉的图。

\subsection{检索相关性的规定}
\label{sec:vision-recognition-relevance}

“找类似的红色汽车”和“找同一辆车的其他照片”使用不同的正例。前者容许不同实例，只要外观或语义相近；后者需要车身划痕、车牌等区分个体的证据，同时允许颜色因光照而变化。\term{实例检索}（instance retrieval）指后一种任务：正例是同一个物体或地点的照片，而不是同类物体。若没有预先规定相关性，高相似度排序可能朝错误的目标优化：一个按类别语义训练的表示会把所有红色货车排在前面，这对实例检索恰恰是干扰。

查询的范围也要规定。整幅照片包含建筑、树木和天空，局部查询只指定一块招牌，两者容许的背景、遮挡和视角变化不同。图库中的近重复图会使任务过于容易，干扰图的规模则决定排序有多难。Revisited Oxford和Paris数据集修订了地标实例检索的标注，按难度把正例分为容易和困难两档，并把难以判断的图像标为忽略项，还提供了约一百万张干扰图\scite{vision-1803-11285}。评价前应固定查询范围、正例定义、忽略项、图库版本和干扰图规模；文字查询图像的检索由\ref{chap:vision-multimodal-understanding}章继续讨论。

\subsection{候选图像的召回}
\label{sec:vision-recognition-candidates}

召回的输入是查询图像和预先处理好的图库，输出是前$k$个候选及其图号。难点在于既要快，又不能漏掉视角差异大的正例。图库图像的表示在建库时离线计算一次，查询时只计算查询图的表示，然后在库中搜索。

表示有两种形式。全局表示把整幅图汇总为一个固定维向量，比较只需一次内积，配合近似最近邻索引可以在百万级图库中毫秒级返回候选；GeM是一种为实例检索学习全局向量的方法。局部表示为每幅图保留几百到几千个关键点及其描述向量，信息更细，可以描述招牌一角这样的局部，但直接逐对比较代价高；它在建库时一并提取存储，既可经视觉词汇聚合后用于投票式召回，也是下一步重排的输入。SuperPoint是一种学习得到的局部点检测与描述方法。近似索引带来的召回损失与表示本身的区分能力应分开检查。

这两种表示都有经典前身，而且学习方法正是沿着经典流程的结构逐处替换的。\ref{sec:vision-features-bovw}节的视觉词袋先在每幅图上检测SIFT关键点（\ref{sec:vision-features-sift}节），得到几百到几千个128维描述子；再在大量描述子上用$k$均值聚出一个视觉词汇表，把每个描述子量化到最近的“视觉词”。一幅图于是变成一个视觉词频率向量，按TF-IDF加权：在少数图中出现的词更有区分力，权重更大。检索时用倒排索引只访问与查询共享视觉词的图库图，因而可以扩展到百万级图库。VLAD进一步保留了量化丢掉的信息：它不只数每个视觉词出现几次，而是把落到第$c$个词的描述子与该词中心的残差$\bm x-\bm\mu_c$累加，拼成一个$K\times128$维向量，再归一化，得到更紧凑也更有区分力的全局描述。召回之后，经典流程用\ref{sec:vision-recognition-ransac}节的RANSAC做几何验证：只有能被同一单应或基础矩阵解释的匹配才算数。这条流水线的弱点在前两步：SIFT在大光照变化和弱纹理处检测不稳定，量化把相近但不同的描述子合并或把同一结构拆开，词频统计又看不到空间排列。GeM把“局部特征加聚合”换成卷积特征加可学习的池化，并用三维重建自动产生训练对；SuperPoint把手工的关键点检测和描述换成网络；LightGlue把最近邻匹配换成带上下文的注意力匹配。几何验证保留了下来，因为它不依赖任何外观假设。

\subsubsection{GeM}
\label{sec:vision-recognition-gem}

卷积网络的最后一层特征是一个空间网格，需要汇总成一个向量。平均池化平等计入每个位置，背景中的天空和路面会稀释目标；最大池化只保留每个通道最强的响应，又对单个噪声响应过于敏感。另一个问题是监督：分类网络学到的是类别语义，而实例检索需要区分同类中的不同个体，人工标注“哪些照片是同一建筑”成本很高。GeM针对这两点，提出可学习指数的广义均值池化，并利用三维重建自动挖掘训练图对\scite{vision-1711-02512}。

\term{广义均值池化}（generalized-mean pooling, GeM）对第$k$通道的非负特征$x_{p,k}$计算
\begin{equation}
 g_k=\left(\frac1n\sum_{p=1}^{n}x_{p,k}^{\,q_k}\right)^{1/q_k},
 \qquad q_k>0\text{。}
 \label{eq:vision-gem}
\end{equation}
$q_k=1$为平均，$q_k$趋于无穷时趋近最大值，介于其间时强响应权重更大，但其他位置仍有贡献。对响应$(1,1,1,4)$，平均为1.75，$q=3$时为$(67/4)^{1/3}\approx2.559$，最大值为4。实现时在幂运算前给特征加下界，避免零点的数值问题；指数可以固定，也可以与编码器一同学习。

完整的流程是：全卷积主干去掉最后的池化和全连接层，把图像变成$h\times w\times C$的非负特征（ReLU之后），GeM得到$C$维向量，可选地接一个白化变换，再做$L_2$归一化得到单位向量$\bar{\bm g}$。原文比较了三种主干：AlexNet给出256维、VGG16给出512维、ResNet-101给出2048维描述。以ResNet-101和$1024\times1024$的测试图为例，末层特征约为$32\times32\times2048$，GeM把这1024个位置汇成一个2048维向量。指数$p$可以每通道一个，也可以全部通道共享；原文最终采用共享且可学习的单个$p$，初始化为3，VGG16微调后收敛到约2.92。白化是在训练集描述上学到的线性判别投影，作为后处理接在GeM之后，默认不降维。单位向量满足$\lVert\bm{u}-\bm{v}\rVert_2^2=2-2\bm{u}^{\top}\bm{v}$，按欧氏距离或内积排序结果相同。

训练图对来自运动恢复结构：对大量无标注的地标照片做三维重建，共享足够多三维点的图对作为正对，不同重建模型中的图像作为负对；负对中优先选取当前描述最相近的困难负例，并随训练用当前参数重新挖掘。原文从约740万张网络照片重建出713个三维模型、16.3万余张图，其中551个模型（约13.3万张图）用于训练、162个用于验证，并去掉与Oxford、Paris测试地标重叠的模型；每个训练元组含1个查询、1个正例和5个困难负例，训练图最长边不超过362像素。因此“不需要人工标注”并不等于没有监督依据。一条训练图对$(\bm X,\bm X')$的同实例标签为$y\in\{0,1\}$，经共享编码、GeM和归一化后距离为$d=\lVert\bar{\bm g}-\bar{\bm g}'\rVert_2$，对比目标为
\begin{equation}
 \ell=\tfrac12 y d^2+
 \tfrac12(1-y)[\max(0,m-d)]^2\text{。}
 \label{eq:vision-gem-contrastive-process}
\end{equation}
正对的梯度推动距离缩短；负对只在$d<m$时受到推开的梯度，超过间隔$m$就不再计入。原文的间隔$m$对AlexNet、VGG16和ResNet-101分别取0.7、0.75和0.85。教学中$m=1$，正对距离0.4的目标为0.08，负对距离0.4的目标为0.18，负对距离1.2为零。编码器从分类权重初始化，训练共同更新编码器与池化指数。

使用时为每幅图库图像保存归一化描述与图号，对查询做完全相同的处理，按内积取前$k$个候选。指数、白化和归一化须一致，只对图库做白化而忘了查询，会改变距离定义。同一建筑的两张照片可能因视角差异而距离较远，另一栋同样玻璃幕墙的建筑却可能很近；全局描述只提供候选，不能证明局部细节一致。已有的DINOv3向量也可以作为召回基线，比较时应固定输入、归一化与图库。

\subsubsection{SuperPoint}
\label{sec:vision-recognition-superpoint}

局部特征需要两样东西：在不同照片中能重复检测到的位置，以及能把对应位置配起来的描述。SIFT等手工方法用设计好的规则完成这两件事\scite{vision-lowe2004}，在模糊、光照变化和弱纹理处常常失效；用学习方法改进时，困难在于“兴趣点”没有自然的人工标注，人很难一致地标出所有可重复的角点。SuperPoint的想法是用自监督产生兴趣点标签：先在具有已知角点的合成几何图形上训练检测器，再对真实图像施加多个已知的单应变换，把各次检测结果变换回原图并汇总，作为更稳定的伪标签；检测与描述共用一个编码器联合训练\scite{vision-1712-07629}。这里的兴趣点没有“左腕”那样的对象部位语义，只要求适合跨图重复辨认。

输入是$H\times W$的灰度图，VGG式编码器经三次下采样得到$H/8\times W/8\times128$的特征。编码器有8个$3\times3$卷积层，通道依次为64、64、64、64、128、128、128、128，每两层之后接一个$2\times2$最大池化（共三个）。两个头各是一个256通道的$3\times3$卷积加一个$1\times1$卷积，检测头的$1\times1$卷积输出65通道，描述头输出256通道；全网各层都接ReLU和批归一化（输出层除外）。逐层结构见表\ref{tab:vision-recognition-superpoint-layers}。

\begin{table*}[tbp]
 \centering
 \small
 \begin{tabular}{@{}lllrrl@{}}
 \toprule
 阶段或层组 & 层类型 & 卷积核与步长 & 输出通道 & 重复次数 & 输出尺寸\\
 \midrule
 编码器1 & 卷积；池化 & $3\times3$；$2\times2$，步长2 & 64 & 2 & $H/2\times W/2$\\
 编码器2 & 卷积；池化 & $3\times3$；$2\times2$，步长2 & 64 & 2 & $H/4\times W/4$\\
 编码器3 & 卷积；池化 & $3\times3$；$2\times2$，步长2 & 128 & 2 & $H/8\times W/8$\\
 编码器4 & 卷积 & $3\times3$ & 128 & 2 & $H/8\times W/8$\\
 检测头 & 卷积 & $3\times3$；$1\times1$ & 256；65 & 1 & $H/8\times W/8$\\
 描述头 & 卷积 & $3\times3$；$1\times1$ & 256；256 & 1 & $H/8\times W/8$\\
 \bottomrule
 \end{tabular}
 \caption{SuperPoint的逐层结构，取自原文第6节。重复次数指每组中$3\times3$卷积的个数；两个头共享编码器。检测头输出经softmax、去掉“无点”通道并重排后得到$H\times W$的点分数图，描述头输出经双三次插值和$L_2$归一化得到每点的256维描述，这两步都没有可学习参数。}
 \label{tab:vision-recognition-superpoint-layers}
\end{table*}

检测头输出$H/8\times W/8\times65$：每个粗格对应原图一个$8\times8$区域，64个通道表示点落在其中哪个像素，第65个通道表示“该区域没有点”。对65个通道做softmax，去掉无点通道，再把64个通道重排回$8\times8$的空间位置，就得到原图分辨率的点分数图，不需要反卷积。描述头输出$H/8\times W/8\times256$的粗描述，在关键点位置做双三次插值采样，并$L_2$归一化。最终保存每个点的$(x,y,\text{得分},\bm{d})$。

训练分三步：在合成形状上训练检测头（这一阶段的模型称为MagicPoint）；用单应适应为真实图像生成伪标签；用伪标签和已知单应同时训练检测和描述。检测损失是每个粗格上65类的交叉熵。描述损失利用已知单应：一幅图与其变换图之间，对应的粗格描述内积应大于正间隔$m_p=1$，不对应的粗格应小于负间隔$m_n=0.2$，形成铰链损失；由于不对应的粗格对远多于对应的，正对项乘以$\lambda_d=250$来平衡，描述损失整体再乘0.0001后与检测损失相加。单应适应对每幅图采样$N_h=100$个随机单应，原文做了两轮。真实图像训练使用约8万张COCO图像，缩放为$240\times320$的灰度图；在Titan X上处理$480\times640$的图像约需13毫秒。单应适应提供的只是平面变换下的监督，不能覆盖任意三维视差。

使用时按阈值和局部峰值抑制保留点，并可按分数限制点数。若输入缩小为原图的一半，网络点$(100,60)$还原为原图$(200,120)$，描述向量本身不变。保存点坐标所在网格和图号，后续匹配器才能在同一坐标系中做几何验证；没有点时应返回空集合，不虚造对应。窗格重复时，检测器可以在许多窗角上稳定响应，但描述可能也很相似；“点能重复检测”与“点能匹配到正确窗格”是两个不同要求，后者留给重排步骤。

两种召回表示的对比见表\ref{tab:vision-recognition-recall-compare}。它们不是互相替代的方案，而是在召回步骤中为后续计算准备的两份不同信息。

\begin{table}[htbp]
 \centering
 \small
 \begin{tabularx}{\linewidth}{@{}lXX@{}}
 \toprule
 比较项 & GeM & SuperPoint\\
 \midrule
 每图存储 & 一个$C$维单位向量 & 数百至数千个点及256维描述\\
 比较成本 & 一次内积，可用近似索引 & 点集之间的匹配，代价高\\
 监督来源 & 三维重建挖掘的正负图对 & 合成形状与单应适应的伪标签\\
 擅长 & 大图库中快速召回 & 局部对应，供几何验证\\
 易错情形 & 外观相似的不同实例 & 重复纹理、大视差\\
 \bottomrule
 \end{tabularx}
 \caption{召回步骤中两种表示的比较。前者决定哪些图像进入候选，后者为候选的重排准备局部证据。}
 \label{tab:vision-recognition-recall-compare}
\end{table}

\subsection{候选图像的重排}
\label{sec:vision-recognition-verification}

重排的输入是查询图和召回的前$k$个候选，每张图都附有局部点与描述；输出是重新排序的列表，并附每张候选与查询之间的对应点和几何变换。难点在于外观相似的其他实例：同样的玻璃幕墙、同样的窗格会产生许多描述相近的点对。重排分两步：先在两图的点之间建立对应，再检查这些对应能否被一个几何变换共同解释。LightGlue负责第一步，用注意力利用上下文给出点对；RANSAC负责第二步，在含大量错误的点对中寻找被多数支持的几何模型。几何通过之后，最终的相关性仍要回到查询定义判断。

\subsubsection{LightGlue}
\label{sec:vision-recognition-lightglue}

独立的最近邻匹配只比较两个描述，无法利用邻近点怎样共同移动的信息，重复窗格因此很容易配错。用注意力让两图的点相互交流可以解决这一问题，但此前的注意力匹配器对每一对图像都执行固定层数的计算，简单图对与困难图对花费相同。LightGlue的想法是在注意力匹配的基础上让计算量自适应：容易的图对提前结束，已经确定无法匹配的点提前丢弃\scite{vision-2306-13643}。它接收已有的关键点和描述（例如SuperPoint的输出），不负责从像素检测点。

两图分别有$n,m$个点，每个点的输入是归一化坐标和描述。网络有9个相同的层，宽度$d=256$，每个注意力单元4个头；每层先做图内自注意力，再做跨图交叉注意力。旋转位置编码只用在自注意力中：它按两点坐标之差旋转查询和键，使注意力得分只依赖相对位置，点因此能感知邻近点的空间排列；交叉注意力不加位置编码，因为两图之间没有共同的坐标系。配置见表\ref{tab:vision-recognition-lightglue-config}。

\begin{table}[htbp]
 \centering
 \small
 \begin{tabular}{@{}lcccccc@{}}
 \toprule
 模型 & 词元 & 层数 & 宽度 & 头数 & MLP倍率 & 参数量\\
 \midrule
 LightGlue & 每个关键点1个 & 9 & 256 & 4 & 论文未给出 & 论文未给出\\
 \bottomrule
 \end{tabular}
 \caption{LightGlue的配置。LightGlue不切图像块，词元是两图的关键点，数量随输入而变（原文微调时每图取2048个点）；每层含一次自注意力和一次交叉注意力。}
 \label{tab:vision-recognition-lightglue-config}
\end{table}每层之后都可以接匹配头：把更新后的描述投影为$d$维，以$1/\sqrt d$缩放两两内积得到$S_{i,j}$；另由线性层和sigmoid给出每个点的可匹配性$\mu_i,\nu_j\in[0,1]$。匹配概率为
\begin{equation}
 P_{i,j}=\mu_i\nu_j
 \frac{\exp S_{i,j}}{\sum_{j'=1}^{m}\exp S_{i,j'}}
 \frac{\exp S_{i,j}}{\sum_{i'=1}^{n}\exp S_{i',j}}\text{。}
 \label{eq:vision-lightglue-score}
\end{equation}
第一个分式要求$i$在对面所有候选中偏好$j$，第二个要求$j$在本侧所有候选中偏好$i$，可匹配性则允许被遮挡或越界的点放弃匹配。若$\mu_i=0.9,\nu_j=0.8$，两个方向的归一化分数为0.7和0.6，则$P_{i,j}=0.3024$。

训练数据是带真实点对的图像对，点对由深度和相机重投影得到，或由已知单应得到；遮挡和越界的点必须标为无对应。若真实匹配集合为$\mathcal M$，两侧无对应点集为$U,V$，目标为
\begin{equation}
 \begin{aligned}
 L_{\mathrm{match}}={}&-\frac1{|\mathcal M|}
 \sum_{(i,j)\in\mathcal M}\log P_{i,j}\\
 &-\frac1{|U|}\sum_{i\in U}\log(1-\mu_i)
 -\frac1{|V|}\sum_{j\in V}\log(1-\nu_j)\text{。}
 \end{aligned}
 \label{eq:vision-lightglue-training-process}
\end{equation}
空集合对应的项应跳过。这一目标在每一层都计算，使任何一层都能输出可用的匹配，这是提前退出的前提。之后另训练一个小的置信分类器，预测某个点在当前层的决定是否已与最终层一致，其梯度不回传到主干，避免为早退判断改动描述。匹配器的权重与所用的描述类型绑定，为一种特征训练的权重不能无条件用于另一种描述。

推理时每层结束后检查置信分类器。第$\ell$层的置信阈值为$\lambda_\ell=0.8+0.1e^{-4\ell/L}$，浅层要求更高的置信；当至少95\%的点超过阈值时提前退出。置信已超过阈值、可匹配性又低于0.01的点被判为不可匹配，从后续层中删去。原文报告容易的图对2至3层后就退出；训练先在约17万张图生成的合成单应图对上预训练，再在MegaDepth上微调。最终取$P$中行最大与列最大互相指向且超过阈值的点对，输出两幅原图上的坐标和匹配得分。任一图没有点时应返回空对应，不能为了让后续几何估计运行而强行填点。$P_{i,j}$是模型的配对评分，不是几何正确性的证明；比较速度时还需固定特征类型、图像尺寸、硬件和图对难度分布。

\subsubsection{RANSAC}
\label{sec:vision-recognition-ransac}

LightGlue的点对中仍会有错误，重复纹理可以产生看似可信但错误的对应。最小二乘法用全部点对拟合几何模型，少数严重错误就会把结果拉偏。\term{随机采样一致性}（random sample consensus, RANSAC）的想法是反过来：每次只随机抽取恰好能确定模型的最少点对，估计一个候选模型，再数有多少点对与它相符，保留获得最多支持的模型\scite{vision-ransac-1981}。只要某一次抽到的全是正确点对，就能得到好模型，而错误点对再多也只是不支持它。

几何模型取决于场景。平面上的点在两图间由\term{单应变换}（homography）联系，齐次坐标满足$\tilde{\bm{p}}'\sim\bm{H}\tilde{\bm{p}}$；一般三维场景的两视图则遵守\term{极线约束}（epipolar constraint）$\tilde{\bm{p}}'^\top\bm{F}\tilde{\bm{p}}=0$，$\bm{F}$为基础矩阵，它只约束对应点落在另一图的一条线上。以单应为例，消去齐次比例可写$\tilde{\bm p}'\times\bm H\tilde{\bm p}=\bm0$，每个点对提供两条独立线性约束，四对约束$\bm H$除尺度外的八个自由度；共线或重复的点会使估计退化。先把坐标平移缩放到合适的数值范围再求解，再变回原图。候选模型的对称转移误差为
\[
 e_i^2=\lVert\bm p_i'-\Pi(\bm H\tilde{\bm p}_i)\rVert_2^2
       +\lVert\bm p_i-\Pi(\bm H^{-1}\tilde{\bm p}_i')\rVert_2^2\text{，}
\]
要求$\bm H$可逆；误差低于像素阈值的点对为内点，其余为外点。

RANSAC不需要训练，它的“学习”是在当前点对上的搜索。算法从无模型开始，反复执行抽样、估计、计内点；若新模型的内点更多，或内点相同而残差更小，就替换当前最佳；最后用全部内点重新估计模型。抽样次数决定成功的可能：每次抽取$s$个点对，内点比例为$w$，全内点样本的概率约为$w^s$，$k$次都没抽到的概率为$(1-w^s)^k$。要求它不超过$1-p$，取对数得
\begin{equation}
 k\geq\frac{\log(1-p)}{\log(1-w^s)}\text{，}
 \label{eq:vision-ransac-iterations}
\end{equation}
分母为负，除法时不等号方向随之改变。取$w=0.5,s=4,p=0.99$，至少需要72次；内点率降低时次数迅速增加，$w=0$时不存在有限次数保证。实现常根据当前最佳模型的内点率动态更新$k$，这依赖估计可信，错误的高支持模型也可能使搜索过早结束。

原始RANSAC只用残差判断内点，忽略了点在空间上的相邻关系。Graph-Cut RANSAC在局部优化阶段为每个点对设二值标签$\ell_i$（内点或外点），把局部目标写成$\sum_iD_i(\ell_i;e_i)+\lambda\sum_{(i,j)\in E}w_{i,j}\mathbf1\{\ell_i\ne\ell_j\}$：一元项按残差惩罚不合适的标记，邻接项倾向于让空间上相邻的点对标签一致\scite{vision-1706-00984}。在满足二值可割条件时，图割求得这一次标记的最优解，再用选出的内点重估几何。它解决的是给定模型附近的局部标记问题，不是对所有模型和标签组合的全局最优；邻近的错误匹配也可能外观一致，仍须检查模型假设。

重排时以内点数量、内点覆盖查询范围的比例和残差为依据重新排序，排序规则应事先固定。建筑正面近似平面时，单应适用；有明显前后视差时，同一个单应无法对齐所有表面，应改用基础矩阵或分区建模。纯旋转、近共线点和重叠太少都会改变估计条件。一小块重复窗格获得的高内点数，也不能替代整幅查询的相关性。图\ref{fig:vision-retrieval-evidence}把全局召回、局部配对与几何验证分开列出，避免把一个阶段的高分误作整条证据链已经完成。

\begin{figure}[htbp]
 \centering
 % 流程节点间箭头均连接左右中心；下方为各阶段尚未排除的错误。
\begin{tikzpicture}[font=\figurefont\small,
 block/.style={draw=BookTeal,fill=BookTeal!8,minimum width=28mm,
   minimum height=12mm,align=center,inner sep=1.5mm},
 flow/.style={-{Stealth[length=2mm]},draw=BookInk,line width=0.9pt},
 note/.style={align=center,text width=28mm,font=\figurefont\footnotesize,
   text=BookMuted}]
 \node[block] (global) at (0,0) {全局描述\\筛选候选};
 \node[block] (local) at (3.7,0) {局部对应\\匹配细节};
 \node[block] (geo) at (7.4,0) {几何支持\\核验重排};
 \draw[flow] (global.east)--(local.west);
 \draw[flow] (local.east)--(geo.west);
 \node[note] at (0,-1.6) {可能只同类\\并非同一实例};
 \node[note] at (3.7,-1.6) {重复窗格\\可能相互错配};
 \node[note] at (7.4,-1.6) {几何匹配仍需\\满足查询定义};
 \foreach \x in {0,3.7,7.4}
   \draw[draw=BookRule,line width=0.7pt] (\x,-0.65)--(\x,-1);
\end{tikzpicture}

 \caption{图像检索中逐步保留更具体的证据。上方沿主流程减少候选，下方说明每步仍可能留下的错误。几何一致可以帮助重排，却仍要检查所得实例是否满足查询的相关性定义。}
 \label{fig:vision-retrieval-evidence}
\end{figure}

两种重排方法的对比见表\ref{tab:vision-recognition-rerank-compare}。LightGlue用学习得到的上下文消解描述的歧义，给出更多正确点对；RANSAC不学习任何参数，而是用几何模型筛掉剩余的错误。二者串联使用：没有好的点对，RANSAC需要的抽样次数会急剧增加；没有几何验证，LightGlue的高分点对也可能整体指向一个相似的错误实例。

\begin{table}[htbp]
 \centering
 \small
 \begin{tabularx}{\linewidth}{@{}lXX@{}}
 \toprule
 比较项 & LightGlue & RANSAC\\
 \midrule
 输入 & 两图的点与描述 & 两图的候选点对\\
 输出 & 点对及匹配概率 & 几何模型与内点集合\\
 依据 & 学习的图内与跨图上下文 & 单应或极线几何的一致性\\
 是否训练 & 需要带真实对应的图对 & 无需训练，按次数搜索\\
 失败情形 & 重复纹理、描述类型不符 & 内点率过低、退化配置、模型不符\\
 \bottomrule
 \end{tabularx}
 \caption{重排步骤中两种方法的分工。前者提高对应质量，后者检验对应能否被同一几何关系解释。}
 \label{tab:vision-recognition-rerank-compare}
\end{table}

\section{识别与检索结果的检验}
\label{sec:vision-recognition-evaluation}

本章的各种输出最终要组合使用，检验因此不能只看单项分数。组合系统需要让类别、框、掩码、属性和测量指向同一对象：若车辆检测正确，属性却来自旁车，逐模块的总体准确率仍可能很高。街景应逐实例保留标识并核验空间对应；工业零件应同时检查异常热图、整图得分与规定的缺陷，避免正常材质变化被业务阈值放大。

\ref{sec:vision-recognition-challenges}节已经给出各项指标的定义，这里关心的是报告这些指标时必须固定什么。检测的AP取决于插值方式、IoU阈值、尺寸分组和每图检测数量上限：PASCAL VOC式的单阈值AP与COCO式在10个阈值上平均的AP不是同一个量，COCO还规定每图最多计入100个得分最高的检测，并按面积小于$32^2$、介于$32^2$与$96^2$之间、大于$96^2$像素分出小、中、大三组。因此不能把任意一个“AP”视为同一指标，跨论文比较前要先核对协议。

区域任务有不同的分母。mIoU按类别平均，小类别与大类别同权，适合检查路杆这类细结构；按像素计的总体准确率则被道路和天空主导。实例分割的掩码AP先要匹配对象；式\eqref{eq:vision-pq}的PQ同时受范围与实例错误影响，前文的例子已说明只报告匹配区域的平均IoU会掩盖漏报。

检索的式\eqref{eq:vision-retrieval-ap}在相关项全部可见时才是完整的AP；只返回前若干项时，计算截断指标须说明分母和遗漏规则，无相关项的查询应按协议排除或单独评价。Revisited Oxford与Paris的“忽略项”在计算时既不算正例也不算负例，若误当负例，排在前面的难以判断的图就会无端压低AP。

开放条件需要单独的检验。未知发现在隐藏新类名称的条件下计算未知召回，同时记录未知被误报为已知的数量、背景假阳性，以及后续加类前后旧类的AP。教学中有5个未知物体，正确发现3个、误当已知1个、完全漏掉1个，未知召回为$3/5=0.6$；这不能由已知货车的AP代替。文字定位和概念分割要固定候选短语、正负示例及提示成本，使用真实框作提示取得的优势不能直接归于无提示系统。

测量量按单位和坐标检验。姿态比较固定关节集合、单位和坐标，以MPJPE计算$(1/j)\sum_k\lVert\hat{\bm J}_k-\bm J_k\rVert_2$。根对齐只去平移，相似变换对齐还去旋转和尺度，二者回答的问题不同。人体表面另比对应顶点误差、轮廓重投影和手足局部，少量关节误差不能证明全部网格准确。单目深度需区分相对排序、边界和米制误差，允许对齐的结果不能与未对齐的度量结果混为一表。离线场景比较固定视图集合、分辨率、首相机参照和拟合步数，分别报告相机、点图、重投影、完整度及总成本；流式场景只允许已到达的图，记录缓存预算、局部误差与长程漂移，一次全集对齐后的轨迹误差会隐藏部分累计误差。三维语义先在有效几何下找到同一表面，再比较各视角的查询归属；RGB渲染、语言相似度和几何一致各有误差来源，三者之一正确不能替其他两项通过。表\ref{tab:vision-recognition-protocols}汇总了各类输出需要固定的条件。

\begin{table}[htbp]
 \centering
 \small
 \begin{tabularx}{\linewidth}{@{}lXX@{}}
 \toprule
 输出 & 需要固定的条件 & 不能由它单独证明的内容\\
 \midrule
 类别、属性 & 类别表、未知项、对象范围 & 对象存在及属性归属都正确\\
 框、掩码 & 匹配规则、提示、可见范围 & 未观测部分及三维形状正确\\
 数量、姿态 & 计入规则、关节与对齐方式 & 实例清单、身份及原始朝向正确\\
 深度、点图 & 尺度、坐标、有效区域 & 距离带实际单位且适合测量\\
 检索排序 & 正例、图库、忽略项及难度 & 对象属性、文字和距离正确\\
 \bottomrule
 \end{tabularx}
 \caption{不同视觉成果所需的评价条件。右列列出仍需另行核验的推论，不表示这些任务必然无法提供相应能力。}
 \label{tab:vision-recognition-protocols}
\end{table}

\section{本章小结}
\label{sec:vision-recognition-summary}

回到开篇的路口照片。“这是街道”由场景分类给出，类别表固定时用冻结表示加分类头，类别在使用时给出时用图文相似度；“哪辆是红色”要先把属性限制在对象区域内；“车在哪里”按需要的精细程度给出框、像素类别、实例掩码或透明度；“有几辆、离多远”在定位之上引入计数规则、相机模型和尺度；“另一张照片里是不是同一辆车”则由全局召回、局部配对和几何验证逐步回答。每一步都把预测接回明确的对象、坐标和单位，这正是本章要建立的从像素到可核验结果的计算。

贯穿各节的是同一类设计选择：哪个输出负责哪个目标，目标由谁指定，结果在什么坐标和尺度下成立。检测从网格、锚框、逐位置距离走到集合匹配，再把类别表换成文字或对象性；掩码从框内预测走到查询掩码、点提示和概念提示；几何从逐对点图走到多视图联合、固定容量拟合和流式缓存。每一次改动都换来一种能力，也带来新的代价，表格中的比较说明了这些交换。

第一节提出的五类挑战，在本章得到的回答程度并不相同。外观变化带来的类别歧义，经典方法用梯度直方图和视觉词袋这类人工不变量缓和，学习方法用大规模预训练表示加分类头或文字类别处理，但背景捷径和细粒度部位仍需单独检查。尺度与遮挡带来的定位困难，从滑动窗口、图像金字塔和级联，演变为卷积的逐位置预测、特征金字塔和焦点损失，再到集合匹配；拥挤遮挡下的漏检并没有消失。稠密输出与分辨率的矛盾，由空洞卷积、编码器—解码器和查询式掩码缓解，而阈值、分水岭、图割和超像素给出的边界线索仍在交互分割和抠图中延续。单目观测带来的几何不确定无法由网络消除，只能由多视图观测或训练先验补足，本章在每一处都标明了尺度的来源；双目视差、平面扫描与光束法平差的结构，分别在MVSNet的代价体和前馈几何模型的全局对齐中延续下来。开放类别与长尾带来的泛化困难，靠文字类别、对象性、正常记忆和自动挖掘的训练对得到部分缓解，几何验证则是经典与学习两条路线共用的、与类别无关的核验手段。各模型的逐层结构表把这些设计落到了具体的层数、通道和张量形状上，读者可以据此复核每一步的输入输出。

仍未解决的部分同样清楚。候选高分不等于目标存在，区域正确不等于身份正确，相对深度不等于米制距离，几何支持也不等于任意业务相关性。开放概念增加了目标歧义，长尾类别需要更细的覆盖分析，遮挡和单目投影限制了可恢复的空间量；有限缓存不保证长期回访没有漂移，三维语义只有与位置和可见表面共同核验才可用。LVIS把大词表和少见实例纳入分割评价，MVTec AD把正常训练与测试缺陷定位分开，为这些问题提供了不同的观察条件\scite{vision-1908-03195,vision-mvtec2019}。连续帧中的身份、事件和检索留给\ref{chap:vision-video-recognition}章，文字与图像的联合定位和问答留给\ref{chap:vision-multimodal-understanding}章。

\paragraph{扩展阅读}
检测方面，Deformable DETR针对DETR收敛慢、小目标弱的问题，让每个查询只在参考点附近采样少量位置做注意力，从而能读取多尺度特征图；Grounding DINO中的可变形交叉注意力和OWOBJ实验所接的检测器都以它为基础\scite{vision-2010-04159}。DINO检测器在此之上加入含正负样本的去噪训练、混合查询选择和逐层修正参考框，Grounding DINO正是把这一检测器扩展到文字类别\scite{vision-2203-03605}。

分割、姿态与深度方面，PSPNet在主干末端对特征图做多种尺度的池化再拼接，用金字塔池化汇集全局上下文，可与DeepLabv3+以空洞卷积扩大感受野的做法对照\scite{vision-1612-01105}。HRNet不采用先降分辨率再恢复的结构，而是全程并行保持高分辨率分支，并在各分辨率之间反复交换信息，适合姿态和语义分割这类需要精确位置的输出\scite{vision-1908-07919}。ViTPose表明普通ViT主干加轻量解码器回归关节热图即可在COCO关键点上取得很强的结果；它采用先检测人、再逐人估计的自顶向下流程，与OpenPose自底向上组装肢体的方式不同\scite{vision-2204-12484}。MiDaS用对尺度和平移不变的损失混合多个标注口径不同的深度数据集训练，输出相对深度，与Depth Anything V2同样面向相对深度；ZoeDepth在相对深度预训练之上加入按场景类型选择的米制深度头，尝试把相对深度换算为以米为单位的距离\scite{vision-1907-01341,vision-2302-12288}。

匹配检索与基准方面，DELG在一个网络中同时输出GeM池化的全局描述子和经注意力筛选的局部描述子，把召回和重排所需的两类特征合到一次前向计算中\scite{vision-2001-05027}。SuperGlue是LightGlue的前身，它用图神经网络在两组关键点之间交替做自注意力和交叉注意力，再用Sinkhorn迭代求解带“未匹配”选项的最优传输\scite{vision-1911-11763}。LoFTR不检测关键点，先在粗网格上用Transformer建立稠密对应、再逐块细化，在弱纹理区域更有优势\scite{vision-2104-00680}。COCO的原始论文交代了本章多次使用的检测和实例分割标注怎样采集、类别怎样设计\scite{vision-1405-0312}；ADE20K提供150类的场景解析基准，常用于比较语义分割与全景分割模型\scite{vision-1608-05442}。
